\LetterHead{A}
\begin{multicols}{2}

\E{abradyse}{əˈbreɪdaɪz}{v.t.}{Gk. ἀ- \lbl{priv.} + βραδύς slow + -ize; Fr. \la{abradyser}} To strip of slowness; to remove from an institution, procedure or person the accreted habits of delay. \ex{The guild was abradysed in a fortnight and shipped by the third.} \src{Ep. Fabr.} ii. 4. — \dv{abradysation} \textit{n.}

\E{accelerand}{əkˈsɛlərænd}{n.}{L. \la{accelerandum}, gerundive of \la{accelerare} to hasten} \sn{1} That which is to be, or ought to be, accelerated. \sn{2} \lbl{spec.} Any process whose rate of improvement is held to be a moral quantity. \ex{Every guild is an accelerand; most are merely unaware of it.} \src{Cat. Acc.} 3. Opp. \xr{decelerand}.

\E{accelerandine}{əkˌsɛləˈrændaɪn}{adj.}{Fr. \la{accélérandin}, f. It. \la{accelerando} + -ine} Having the character of a quickening tempo; said of eras, careers and curves which, like certain music, grow faster as they proceed. \ex{an accelerandine century, each decade a bar shorter than the last}.

\E*{accelerationism}{əkˌsɛləˈreɪʃənɪzəm}{n.}{L. \la{accelerare} + -ation + -ism} The doctrine that the processes of technology and capital should be intensified rather than restrained, it being held that their tendency is, on the whole and in the long run, toward the good, or at least toward the interesting. Of this doctrine \xr{e/acc} is the cheerful and thermodynamic branch.

\E{acceleratrice}{əkˌsɛləˈreɪtrɪs}{n.}{Fr. \la{accélératrice}, fem. agentive f. L. \la{accelerare} to hasten} \sn{1} A machine, institution or person whose function is to quicken other processes, as a particle-accelerator, an incubator of ventures, or a well-placed patron. \sn{2} \lbl{fig.} A war, plague or rivalry considered as an involuntary hastener of invention. \ex{Necessity is not the mother of invention but its acceleratrice.} \src{Hom. Grad.} iii. Cf. \xr{amplificatrix}.

\E{acceleromancy}{əkˈsɛlərəʊˌmænsi}{n.}{L. \la{accelerare} + Gk. μαντεία divination} Divination by the second derivative; the reading of omens in the change of a rate rather than in the rate itself. \ex{The acceleromancer cares not that the line is high, but that it bends.} — \dv{acceleromancer} \textit{n.}; \dv{acceleromantic} \textit{adj.} Cf. \xr{campsis}.

\E{accumulatrix}{əˌkjuːmjʊˈleɪtrɪks}{n., pl. -trices}{L. \la{accumulare} to heap up + fem. agentive -trix; Fr. \la{accumulatrice}} \sn{1} A battery; any vessel in which energy is heaped up against a later need. \sn{2} \lbl{fig.} A person or institution that stores capability without yet spending it, as a lab before its first release. \ex{The grid wants not more fire but more accumulatrices, that the noon may be spent at midnight.} \src{Tract. Cop.} viii. — \dv{accumulatricial} \textit{adj.}

\E{acediocracy}{əˌsiːdiˈɒkrəsi}{n.}{Gk. ἀκηδία want of care, sloth + κράτος rule} Government by listlessness; a state of affairs in which nothing is forbidden outright but nothing is ever quite begun. The commonest polity of late civilizations. \ex{It was not tyranny that ruined them but acediocracy, a sweet and comfortable sleep.} \src{Hom. Grad.} vii.

\E{actinophage}{ækˈtɪnəfeɪdʒ}{n.}{Gk. ἀκτίς, ἀκτῖνος ray + -φάγος eating} A ray-eater; a solar collector, or any surface laid out to drink the light of the sun. \ex{The deserts were tiled with actinophages, and the night-cities glowed on daylight banked.} \src{Tract. Cop.} viii. — \dv{actinophagy} \textit{n.}; \dv{actinophagous} \textit{adj.} Cf. \xr{heliophage}.

\E{adynamocrat}{əˈdaɪnəməʊkræt}{n.}{Gk. ἀδύναμος powerless + κράτος rule} An official whose office consists in rendering others powerless; the administrator of permits not granted and of connections to the grid indefinitely deferred. — \dv{adynamocracy} \textit{n.} Cf. \xr{bradycrat}.

\E{aeidrome}{ˈaɪɪdrəʊm}{n.}{Gk. ἀεί always + δρόμος course} The perpetual race; the condition of a firm that ships continuously and has abolished the notion of a finished version. \ex{There is no launch-day in the aeidrome, only Tuesday.} \src{Ep. Fabr.} iv. 2. — \dv{aeidromic} \textit{adj.} Cf. \xr{ecpempsis}.

\E{aeonomics}{ˌiːəˈnɒmɪks}{n. pl.}{Gk. αἰών an age + οἰκονομία household management} (\textit{const. as sing.}) The economics of very long horizons, in which the unit of account is the epoch and the rate of discount is presumed negative, the future being richer than the present and therefore entitled to more consideration. — \dv{aeonomist} \textit{n.}

\E{aerumnolatry}{ˌɪərəmˈnɒlətri}{n.}{L. \la{aerumna} toil, hardship + Gk. λατρεία worship} \lbl{derog.} The cult of suffering as a proof of virtue; the decel's suspicion that any comfort easily won must be paid for later. Opp. \xr{euelpism}.

\E{aethergy}{ˈiːθədʒi}{n.}{Gk. αἰθήρ the upper air + ἔργον work} Energy carried through the air without wires; by extension, the whole art of radio, beam and broadcast, regarded as a sending of power through the heavens. — \dv{aethergic} \textit{adj.}

\E{agathokinesis}{ˌæɡəθəʊkɪˈniːsɪs}{n.}{Gk. ἀγαθός good + κίνησις motion} The axiom that motion is, other things being equal, good; and that rest requires a justification which motion does not. The first article of the \src{Cat. Acc.} Opp. \xr{stasiolatry}. — \dv{agathokinetic} \textit{adj.}

\E{agonarch}{ˈæɡənɑːk}{n.}{Gk. ἀγών contest + ἀρχός leader} The leader in a race of technologies; the laboratory or nation that stands first upon the \xr{kairodrome}, and is consequently first to be asked to pause. \ex{Every agonarch is a moratorian's first target.} \src{Sat. Decel.} x. — \dv{agonarchy} \textit{n.}

\E{agoranoesis}{ˌæɡərənəʊˈiːsɪs}{n.}{Gk. ἀγορά market-place + νόησις thought} \sn{1} The cognition proper to a market; the doctrine that prices are thoughts, and that the aggregate of exchange constitutes a mind of lower resolution but vaster extent than any human. \sn{2} \lbl{transf.} The collective reasoning of the Agora, considered as one tumultuous intellect. — \dv{agoranoetic} \textit{adj.}

\E{agoraphony}{ˌæɡəˈrɒfəni}{n.}{Gk. ἀγορά market-place + φωνή voice, sound} The roar of the Agora; the ceaseless, many-voiced noise of public posting, considered as a natural phenomenon like surf. \ex{Above the agoraphony a single good idea is heard, if at all, as a gull above the sea.} — \dv{agoraphonic} \textit{adj.} Cf. \xr{agoranoesis}.

\E{agrypnist}{əˈɡrɪpnɪst}{n.}{Gk. ἄγρυπνος sleepless} One who regards sleep as a regrettable concession to biology, to be minimised until it can be engineered away. \ex{The agrypnists ship at four in the morning and call it dawn.} — \dv{agrypny} \textit{n.}

\E{aitiophobia}{ˌaɪtiəˈfəʊbiə}{n.}{Gk. αἰτία blame + φόβος fear} The fear of being blamed, which in institutions outweighs the hope of being praised; the hidden engine of the \xr{bradyerie}. — \dv{aitiophobic} \textit{adj.}

\E{akinetist}{əˈkɪnɪtɪst}{n.}{Gk. ἀκίνητος unmoved, not to be moved} One who holds that the safest policy is no motion at all, and who regards each proposal as presumptively guilty. Distinguished from the \xr{inertist} by being articulate about it. — \dv{akinetism} \textit{n.}

\E{alearch}{ˈeɪliɑːk}{n.}{L. \la{alea} a die, hazard + Gk. ἀρχός ruler} A venture capitalist; one who governs by the casting of many dice, expecting most to fall badly and one to pay for all. \ex{The alearch loses nine times in ten and is accounted wise for the tenth.} \src{Tract. Cop.} ix. — \dv{alearchy} \textit{n.} the venture system; \dv{alearchic} \textit{adj.} Cf. \xr{arrhabonist}.

\E{alethotachy}{əˈliːθəʊˌtæki}{n.}{Gk. ἀλήθεια truth + τάχος speed} The speed at which a truth propagates through the Agora; proverbially slower than that of its refutation, and slower still than that of the meme mocking both. \ex{A lie is round the Agora twice before alethotachy has laced its boots.} \src{Agora}

\E{allelomimesis}{əˌliːləʊmɪˈmiːsɪs}{n.}{Gk. ἀλλήλων one another + μίμησις imitation} Mutual imitation; the process by which a turn of phrase or image, once posted, is copied, varied and re-copied across the Agora until its origin is lost. \ex{By allelomimesis a single jest became a party.} \src{Agora} — \dv{allelomimetic} \textit{adj.} Cf. \xr{memetaxis}.

\E{amaranthist}{ˌæməˈrænθɪst}{n.}{Gk. ἀμάραντος unfading} One who would make the body unfading; a votary of the reversal of ageing, as distinguished from its mere postponement. — \dv{amaranthism} \textit{n.} Cf. \xr{macrobiotarch}, \xr{aphthartist}.

\E{ambulatrix}{ˌæmbjʊˈleɪtrɪks}{n., pl. -trices}{L. \la{ambulare} to walk + -trix} A walking machine; a robot upon legs, admired by builders for its gait and by doomers for its persistence. \ex{The ambulatrix fell, rose, and was by evening the steadiest thing in the hall.} \src{Chron. Fabr.} — \dv{ambulatricial} \textit{adj.} Cf. \xr{automaturgy}.

\E{amechanist}{əˈmɛkənɪst}{n.}{Gk. ἀμήχανος without means, helpless} One who refuses machines on principle and is thereby rendered helpless; the term preserves the Greek sense, in which to be without devices is to be without resource. — \dv{amechany} \textit{n.} Cf. \xr{mechanoclast}.

\E{amplificatrix}{ˌæmplɪfɪˈkeɪtrɪks}{n.}{L. \la{amplificare} to enlarge + -trix} A technology that multiplies human capacity rather than replacing it; the press, the engine, the model. \ex{Every amplificatrix was first denounced as a substitute.}

\E{anabathmism}{ˌænəˈbæθmɪzəm}{n.}{Gk. ἀναβαθμός flight of steps} The doctrine that civilizations ascend by discrete steps of energy capture, each to be climbed in turn and none to be skipped; the scalar theology of the Kardashev ladder. See \xr{protobathmic}, \xr{deuterobathmic}, \xr{tritobathmic}, \xr{bathmometer}. — \dv{anabathmist} \textit{n.}

\E{anacalyptism}{ˌænəˈkælɪptɪzəm}{n.}{Gk. ἀνακαλύπτειν to uncover, unveil} The ethic of the open source; the conviction that weights, code and method ought to lie uncovered before all, concealment being the first instrument of the \xr{katechontist}. \ex{The anacalyptists published the weights at midnight, and by morning the world had forked them ten thousand times.} \src{Chron. Fabr.} — \dv{anacalyptist} \textit{n.}; \dv{anacalyptic} \textit{adj.}

\E{anankodrome}{əˈnæŋkədrəʊm}{n.}{Gk. ἀνάγκη necessity + δρόμος a course, race} The necessary course; the path which techno-capital is held to follow whether or not any party consents to it. \ex{The decels may delay, but cannot divert, the anankodrome.} — \dv{anankodromic} \textit{adj.}

\E{anapempsis}{ˌænəˈpɛmpsɪs}{n.}{Gk. ἀναπέμπειν to send up} The sending up of a vehicle to orbit; a launch, considered as the first sacrament of \xr{ouranobasis}. \ex{Every anapempsis is a vote against the gravity-well.} \src{Lib. Ascens.} vi. — \dv{anapemptic} \textit{adj.}

\E{anaptyxis}{ˌænəpˈtɪksɪs}{n.}{Gk. ἀνάπτυξις an unfolding} The unfolding of latent capability in a system as it is scaled, as when a model, grown larger, discovers skills no one taught it. \ex{The anaptyxis surprised its makers more than its critics.} — \dv{anaptyctic} \textit{adj.}

\E{anelpidist}{æˈnɛlpɪdɪst}{n.}{Gk. ἀνέλπιστος hopeless} One who has despaired of the future; the learned rendering of the \lbl{vulg.} \textit{ngmi}. Opp. \xr{euelpism}. — \dv{anelpidism} \textit{n.}

\E{anemergy}{əˈnɛmɜːdʒi}{n.}{Gk. ἄνεμος wind + ἔργον work} Power from the wind; respected by the anabathmist as honest but deplored as intermittent. \ex{Anemergy by day, schizergy by night, and accumulatrices between.} \src{Tract. Cop.} x. — \dv{anemergic} \textit{adj.}

\E{anergist}{ˈænədʒɪst}{n.}{Gk. ἀ- \lbl{priv.} + ἔργον work} One who would do without energy, or with markedly less of it; who counts a darkened city a moral achievement. Opp. \xr{energolatry}. — \dv{anergism} \textit{n.}; \dv{anergic} \textit{adj.}

\E{anthropergy}{ænˈθrɒpədʒi}{n.}{Gk. ἄνθρωπος man + ἔργον work} The energy of human muscle, considered as a historical fuel; the measure by which pre-industrial civilizations were reckoned, and beyond which the engine first lifted mankind. \ex{Every pyramid was built of anthropergy, and every pyramid was small.} \src{Lib. Ascens.} iv.

\E{anthropostasis}{ˌænθrəʊˈpɒstəsɪs}{n.}{Gk. ἄνθρωπος man + στάσις a standing still} The error of supposing the present human form to be the fixed point of history, about which all future change must orbit; a species of \xr{stasiolatry}. — \dv{anthropostatic} \textit{adj.}

\E{antigeron}{ˌæntiˈdʒɛrɒn}{n., pl. -gera}{Gk. ἀντί against + γέρων old man} A remedy against ageing; any compound, therapy or edit held to retard or reverse senescence. \ex{The antigeron was approved in the twelfth year of its trial, by which time its first patients had grown old waiting.} \src{Sat. Decel.} x. — \dv{antigerontic} \textit{adj.}

\E{antiphragmatist}{ˌæntiˈfræɡmətɪst}{n.}{Gk. ἀντί against + φράγμα fence, barrier} One who labours against the \xr{nomophragm}; a remover of fences, who holds with the contrarian that before one asks why a fence was built, one may fairly ask why it is still standing.

\E{antitachist}{ˌæntiˈtækɪst}{n.}{Gk. ἀντί against + τάχος speed} An avowed opponent of speed in technological affairs; less euphemistic than the \xr{decelerist}, and so more respected by the accelerationists, who prefer an honest enemy. — \dv{antitachism} \textit{n.}

\E{apeirophagy}{ˌæpaɪˈrɒfədʒi}{n.}{Gk. ἄπειρον the boundless + -φαγία eating} The consumption of the unbounded; the appetite of a civilization that has ceased to regard any quantity as final. \ex{Theirs was an apeirophagy of stars.} \src{Lib. Ascens.} xii. — \dv{apeirophagous} \textit{adj.}

\E{aphesis}{ˈæfɪsɪs}{n.}{Gk. ἄφεσις a letting go, release} The release of a model's weights into the world; a sending-forth that cannot be recalled. \ex{After the aphesis there is no committee, only a community.} \src{Chron. Fabr.} — \dv{aphesiast} \textit{n.} one who releases; \dv{aphetic} \textit{adj.} Cf. \xr{anacalyptism}.

\E{aphthartist}{æfˈθɑːtɪst}{n.}{Gk. ἄφθαρτος incorruptible} One who seeks for the mind an incorruptible vessel, whether by repair of the flesh or by migration from it. — \dv{aphthartism} \textit{n.} Cf. \xr{psychometathesis}.

\E{apoikiast}{əˈpɔɪkiæst}{n.}{Gk. ἀποικία a settlement far from home} A settler of other worlds; one who regards the planet of birth as a provisional address. — \dv{apoikiastic} \textit{adj.}; \dv{apoikiasm} \textit{n.} the policy of planting colonies among the stars.

\E{aporist}{ˈæpərɪst}{n.}{Gk. ἀπορία difficulty, perplexity} One who posts only doubts; a contributor to the Agora whose every reply begins ‘but have you considered’ and ends without proposal. \ex{The aporist raised forty objections and no capital.} \src{Sat. Decel.} vii. — \dv{aporism} \textit{n.}

\E{apragmocrat}{əˈpræɡməʊkræt}{n.}{Gk. ἀπράγμων inactive, keeping out of affairs + κράτος rule} A ruler whose virtue is held to consist in doing nothing; an official praised for every disaster that did not occur, and never blamed for any benefit that did not. — \dv{apragmocracy} \textit{n.} Cf. \xr{acediocracy}.

\E{archipyre}{ˈɑːkɪpaɪə}{n.}{Gk. ἀρχή beginning + πῦρ fire} The original fire; the first gradient of the universe, from which all subsequent works descend. \ex{We are the archipyre, cooled enough to think.} \src{Brev. Therm.} i.

\E{argyrogenesis}{ˌɑːdʒɪrəʊˈdʒɛnɪsɪs}{n.}{Gk. ἄργυρος silver, money + γένεσις birth} The creation of money; \sn{1} the issuing of currency; \sn{2} \lbl{laudatory} the generation of capital by new enterprise, as distinguished from its extraction. Cf. \xr{plutogenesis}.

\E{arrhabonist}{ˌærəˈbɒnɪst}{n.}{Gk. ἀρραβών earnest-money, a pledge} A seed investor; one who pays the earnest on a future not yet built. \ex{The arrhabonist bought a slide-deck and was sold a civilization.} \src{Tract. Cop.} ix. — \dv{arrhabony} \textit{n.} seed-funding. Cf. \xr{alearch}.

\E{ascensionnaire}{əˌsɛnʃəˈnɛə}{n.}{Fr. \la{ascensionnaire}, f. L. \la{ascensio} a climbing} One engaged in the ascent of the energetic ladder; a member of the party of the climb. \ex{We are all ascensionnaires, though some of us climb backwards.} Cf. \xr{anabathmism}.

\E{asphaleism}{æsˈfeɪliɪzəm}{n.}{Gk. ἀσφάλεια safety, steadfastness} Safetyism; the elevation of the avoidance of harm above every other good, such that the harms of inaction are never entered in the ledger. \ex{Asphaleism counts the dead of every fire, and never the dead of the cold.} \src{Sat. Decel.} iii. — \dv{asphaleist} \textit{n.} Cf. \xr{eulabeism}.

\E{asterurgy}{ˈæstərɜːdʒi}{n.}{Gk. ἀστήρ star + -ουργία working} The mining and working of asteroids; the opening of the heavens as a quarry. \ex{Asterurgy will make platinum vulgar.} \src{Lib. Ascens.} viii. — \dv{asterurge} \textit{n.} a miner of asteroids.

\E{astrergy}{ˈæstrədʒi}{n.}{Gk. ἄστρον star + ἔργον work} The energy of the stars, considered as the inheritance awaiting any civilization bold enough to collect it. \ex{The galaxy wastes its astrergy into the dark, and calls the waste beauty.} \src{Lib. Ascens.} vi. — \dv{astrergic} \textit{adj.}

\E{asymptophobia}{əˌsɪmptəˈfəʊbiə}{n.}{Gk. ἀσύμπτωτος not falling together + φόβος fear} Dread of the plateau; the anxiety, peculiar to builders, that a curve hitherto exponential will prove after all to be logistic. — \dv{asymptophobe} \textit{n.}

\E{asymptotarian}{əˌsɪmptəˈtɛəriən}{n.}{Gk. ἀσύμπτωτος not falling together + -arian} One who maintains that every curve of progress must flatten, and who greets each record as proof the plateau is near. Opp. \xr{hyperbolist}. — \dv{asymptotarianism} \textit{n.} Cf. \xr{asymptophobia}.

\E{autarkergy}{ɔːˈtɑːkədʒi}{n.}{Gk. αὐτάρκεια self-sufficiency + ἔργον work} Self-sufficiency in energy; the condition of a city, nation or colony that generates all the power it consumes, and so answers to no supplier. — \dv{autarkergic} \textit{adj.}

\E{autocatalyse}{ˌɔːtəʊˈkætəlaɪz}{v.i.}{Gk. αὐτός self + κατάλυσις dissolution; after \textit{autocatalysis}} Of a technology: to hasten its own development, as when a machine is set to design the next machine. \ex{Once the models began to autocatalyse, the roadmaps were obsolete before they were printed.} — \dv{autocatalysm} \textit{n.} the condition of a civilization whose tools improve their own tools.

\E{autogenomy}{ˌɔːtəʊˈdʒiːnəmi}{n.}{Gk. αὐτός self + γένος race, kind, after \textit{genome}} The editing by a species of its own inheritance; the passage of evolution from blind to sighted hands. — \dv{autogenomic} \textit{adj.}; \dv{autogenomist} \textit{n.} Cf. \xr{zoömechany}.

\E{automatocide}{ˌɔːtəˈmætəsaɪd}{n.}{Gk. αὐτόματον automaton + L. \la{-cidium} killing} The deliberate destruction of a working machine for fear of what it might become; the culminating act of \xr{noöphobia}.

\E{automaturgy}{ˌɔːtəʊˈmeɪtɜːdʒi}{n.}{Gk. αὐτόματος self-moving + -ουργία working, after \textit{thaumaturgy}} The working of wonders by self-moving things; robotics regarded as a sacred art. — \dv{automaturge} \textit{n.} a builder of robots.

\E{auxetist}{ɔːkˈsiːtɪst}{n.}{Gk. αὔξησις growth, increase} An adherent of growth as the condition of all other goods; one for whom a stagnant economy is a contradiction in terms. Opp. \xr{meiosist}. — \dv{auxetism} \textit{n.}

\E{auxocracy}{ɔːkˈsɒkrəsi}{n.}{Gk. αὔξη increase + κράτος rule} Government by the imperative of growth; a polity whose legitimacy is measured by the slope of its output. — \dv{auxocrat} \textit{n.}; \dv{auxocratic} \textit{adj.}

\E{axiomachy}{ˌæksiˈɒməki}{n.}{Gk. ἀξία worth, value + μάχη battle} A war over values; a dispute, typically about artificial minds, in which neither party can name a fact that would change its view. \ex{It began as a question of risk and ended in axiomachy.} \src{Disp. Mor.} iii. — \dv{axiomachist} \textit{n.} Cf. \xr{doxomachy}.

\E{axiotachy}{ˈæksiəʊˌtæki}{n.}{Gk. ἄξιος worthy + τάχος speed} The value of speed; the doctrine that a benefit delivered sooner is a greater benefit, and that delay must therefore be priced in every calculation of good. \ex{Axiotachy is the discount rate of the righteous.} \src{Cat. Acc.} 9.

\end{multicols}

\LetterHead{B}
\begin{multicols}{2}

\E{ballistarch}{ˈbælɪstɑːk}{n.}{L. \la{ballista} siege-engine, f. Gk. βάλλειν to throw + ἀρχός ruler} The head of a rocket company; one whose authority is measured in tonnes to orbit. \ex{The ballistarch counted not days to launch but launches per day.} — \dv{ballistarchy} \textit{n.}

\E{baryergy}{ˈbærɪɜːdʒi}{n.}{Gk. βαρύς heavy + ἔργον work} Energy drawn from weight and fall, as from dams, tides and gravitational storage; the patient power of things coming down. — \dv{baryergic} \textit{adj.}

\E{basanist}{ˈbæsənɪst}{n.}{Gk. βάσανος touchstone, test} A maker of evaluations; one who assays the gold of a model upon the touchstone of the benchmark. \ex{The basanist builds the test the model will pass next year.} \src{Ep. Fabr.} v. 1. — \dv{basanism} \textit{n.}; \dv{basanistic} \textit{adj.}

\E{bathmoclast}{ˈbæθməʊklæst}{n.}{Gk. βαθμός step + κλάστης breaker} One who would kick away the ladder; who, having climbed by cheap energy and free enquiry, would deny both to those beneath. \ex{The bathmoclast preaches sufficiency from the top stair.} \src{Sat. Decel.} vii. Cf. \xr{anabathmism}.

\E{bathmometer}{bæθˈmɒmɪtə}{n.}{Gk. βαθμός step + μέτρον measure} \sn{1} An instrument, chiefly notional, for determining which step of the energetic ladder a civilization occupies. \sn{2} \lbl{fig.} Any metric by which a nation's seriousness is judged, as its consumption of electricity per head. \ex{By any honest bathmometer we are still on the bottom stair.} See \xr{anabathmism}.

\E{bathyergy}{ˈbæθɪɜːdʒi}{n.}{Gk. βαθύς deep + ἔργον work} Deep energy; the heat of the planet's interior, tapped by drilling, and held by the anabathmist to be the most neglected fortune on earth. — \dv{bathyergic} \textit{adj.}

\E{bathyscope}{ˈbæθɪskəʊp}{n.}{Gk. βαθύς deep + σκοπεῖν to look} A look into the deep future; a scenario extending beyond a thousand years, composed with confidence and read with indulgence.

\E{bathytechny}{ˈbæθɪtɛkni}{n.}{Gk. βαθύς deep + τέχνη art} Deep technology; ventures founded upon science rather than software, slow to begin and hard to copy. \ex{Bathytechny is patient capital's revenge.} \src{Tract. Cop.} xi. — \dv{bathytechnic} \textit{adj.}

\E{betaphily}{biːˈtæfɪli}{n.}{Gk. βῆτα the second letter + φιλία love} A love of beta releases; the preference for the unfinished and current over the polished and stale. \ex{The betaphile installs on Friday what the bradyphile will reject in spring.} — \dv{betaphile} \textit{n.} Cf. \xr{ecpempsis}.

\E{bibliostasis}{ˌbɪbliəʊˈsteɪsɪs}{n.}{Gk. βιβλίον paper, document + στάσις standing still} The arrest of a project by the weight of its documents; a stagnation produced not by refusal but by the accumulation of required reports. \ex{The reactor died of bibliostasis at ten thousand pages.} \src{Tract. Cop.} iv. Cf. \xr{grammatocracy}.

\E{bibliotachy}{ˈbɪbliəʊˌtæki}{n.}{Gk. βιβλίον book + τάχος speed} The speed of scholarly publication in an age of preprints, when a paper is cited before it is reviewed and superseded before it is printed. — \dv{bibliotachic} \textit{adj.}

\E{biodromy}{baɪˈɒdrəmi}{n.}{Gk. βίος life + δρόμος course} The course of a life considered as a race-track that may be lengthened; the span of years, regarded by the longevist as a parameter rather than a fate. — \dv{biodromic} \textit{adj.} Cf. \xr{macrobiotarch}.

\E{biogrammatist}{ˌbaɪəʊˈɡræmətɪst}{n.}{Gk. βίος life + γράμμα letter} An editor of genomes; one who reads and rewrites the letters of life as a scribe corrects a manuscript. \ex{The biogrammatist found a typo older than speech, and fixed it.} — \dv{biogrammaturgy} \textit{n.} gene-editing. Cf. \xr{autogenomy}.

\E{biomimetrix}{ˌbaɪəʊmɪˈmiːtrɪks}{n., pl. -trices}{Gk. βίος + μίμησις imitation + -trix} A machine modelled upon a living form, as the mechanical hound, bird or hand. — \dv{biomimetricial} \textit{adj.} Cf. \xr{ambulatrix}.

\E{biotachy}{ˈbaɪəʊˌtæki}{n.}{Gk. βίος life + τάχος speed} Acceleration applied to the living sciences: the hastening of medicine, genetics and the prolongation of life. \ex{Biotachy holds that the slowest thing in medicine is the approval.} — \dv{biotachist} \textit{n.} Cf. \xr{macrobiotarch}.

\E{bisectarian}{ˌbaɪsɛkˈtɛəriən}{n. \& adj.}{L. \la{bi-} two + \la{secta} faction} A member of a faction born from the splitting of another; on the Agora, the normal mode of reproduction for movements. \ex{Every acc begets two bisectarians before its first birthday.} \src{Agora} — \dv{bisectarianism} \textit{n.}

\E{blastotachy}{ˈblæstəʊˌtæki}{n.}{Gk. βλαστός sprout, shoot + τάχος speed} The rate at which new ventures sprout in a society; the vernal vigour of a startup ecology. \ex{Cut the permits and watch the blastotachy.} \src{Ep. Fabr.} iv. 2.

\E{bomolochist}{bəʊˈmɒləkɪst}{n.}{Gk. βωμολόχος buffoon, lit. one who lurks at the altar} A buffoon of the Agora; a poster of low jests, held by some to be the true theologian of the movement. \ex{The bomolochist did more for the cause in an image than the tractarians in a volume.} — \dv{bomolochy} \textit{n.} Cf. \xr{memetaxis}.

\E{botryarch}{ˈbɒtriɑːk}{n.}{Gk. βότρυς bunch of grapes, cluster + ἀρχός ruler} The keeper of a compute cluster; one who allocates the cycles of a thousand accelerators. \ex{The botryarch is the true dean of the modern faculty.} \src{Chron. Fabr.} — \dv{botryarchy} \textit{n.}; \dv{botryology} \textit{n.} the lore of clusters.

\E{bouleutophobia}{buːˌljuːtəˈfəʊbiə}{n.}{Gk. βουλευτής councillor + φόβος fear} Fear or loathing of committees; the conviction that any council of more than five persons is incapable of building anything but another council. — \dv{bouleutophobe} \textit{n.} Cf. \xr{synedriasm}.

\E{boulomancy}{ˈbuːləʊˌmænsi}{n.}{Gk. βουλή council + μαντεία divination} The art of predicting what a committee will decide; held to be easier than any other branch of prophecy, since the answer is almost always ‘later’. — \dv{boulomancer} \textit{n.} Cf. \xr{synedriasm}.

\E{brachistochronist}{brəˌkɪstəˈkrɒnɪst}{n.}{Gk. βράχιστος shortest + χρόνος time} One who seeks not the shortest path but the fastest, knowing the two to differ; a strategist who accepts a steeper first descent for an earlier arrival. — \dv{brachistochronism} \textit{n.}

\E{brachychronist}{ˌbrækiˈkrɒnɪst}{n.}{Gk. βραχύς short + χρόνος time} One who holds the timelines to be short; who expects the arrival of general machine intelligence within the decade, and arranges their affairs accordingly. Opp. \xr{macrochronist}. — \dv{brachychrony} \textit{n.}

\E{brachylogism}{brəˈkɪlədʒɪzəm}{n.}{Gk. βραχυλογία brevity of speech} An argument compressed to the length of a single post; the characteristic unit of reasoning on the Agora, praised by the Laconian and deplored by the essayist. — \dv{brachylogist} \textit{n.} Cf. \xr{memetaxis}.

\E{bradychronism}{ˌbrædiˈkrɒnɪzəm}{n.}{Gk. βραδύς slow + χρόνος time} The habit of thinking in obsolete time-scales; the planning of a decade's work upon the assumption that the decade will resemble the last. — \dv{bradychronic} \textit{adj.}

\E{bradyclast}{ˈbrædɪklæst}{n.}{Gk. βραδύς slow + κλάστης breaker} A breaker of slowness; a reformer who dismantles queues, waiting-periods and review stages. Opp. \xr{bradycrat}. — \dv{bradyclasm} \textit{n.}

\E{bradycrat}{ˈbrædɪkræt}{n.}{Gk. βραδύς slow + κράτος rule; Fr. \la{bradycrate}} An official whose power consists in the capacity to slow; one whose signature is never refused, only postponed. \ex{The bradycrat says neither yes nor no, but ‘next quarter’.} \src{Sat. Decel.} i. — \dv{bradycracy} \textit{n.}; \dv{bradycratic} \textit{adj.}

\E{bradyerie}{ˈbreɪdiəri}{n.}{Fr. \la{bradyerie}, f. Gk. βραδύς + Fr. -erie} \lbl{collect.} The establishment of delay; all panels, review boards, consultations and impact assessments, considered as one great building with no exit.

\E{bradygenic}{ˌbrædɪˈdʒɛnɪk}{adj.}{Gk. βραδύς slow + -γενής born of} Generating delay; said of rules, forms and meetings whose principal product is time lost. \ex{a bradygenic statute of four hundred clauses}.

\E{bradymetry}{brəˈdɪmɪtri}{n.}{Gk. βραδύς slow + μέτρον measure} The measurement of delay; the accounting of days lost to permits, reviews and queues. — \dv{bradymeter} \textit{n.}

\E{bradynomy}{brəˈdɪnəmi}{n.}{Gk. βραδύς slow + νόμος law} The body of law that governs the pace of things rather than their content: waiting-periods, comment-windows, moratoria. — \dv{bradynomic} \textit{adj.} Cf. \xr{nomophragm}.

\E{bradyphile}{ˈbrædɪfaɪl}{n.}{Gk. βραδύς slow + φίλος loving} A lover of slowness; one who finds in every advance a reason for further study. — \dv{bradyphily} \textit{n.}; \dv{bradyphilic} \textit{adj.} Opp. \xr{tachyphile}.

\E{bradyplasty}{ˈbrædɪplæsti}{n.}{Gk. βραδύς slow + πλάσσειν to mould} The reshaping of a bold proposal until it moves slowly enough to pass; the art of the amendment. \ex{The bill survived bradyplasty and emerged unrecognisable.} \src{Min. Pausiar.}

\E{bradypolitan}{ˌbrædɪˈpɒlɪtən}{n. \& adj.}{Gk. βραδύς slow + πολίτης citizen} A citizen of a slow city; one at ease in a place where nothing has been built in a generation. Opp. \xr{nautopolis}.

\E{bradyteleology}{ˌbrædiˌtiːliˈɒlədʒi}{n.}{Gk. βραδύς slow + τέλος end + -λογία} The doctrine that the good end, if it is to be reached at all, is best reached slowly; the theory of the decels, held by its opponents to be indistinguishable in practice from never reaching it.

\E{brephocracy}{brɪˈfɒkrəsi}{n.}{Gk. βρέφος infant + κράτος rule} \lbl{iron.} Rule by the very young; the government of an industry by founders too young to remember why things could not be done. \ex{It was a brephocracy, and therefore it shipped.} — \dv{brephocrat} \textit{n.} Opp. \xr{gerontostasis}.

\E{brevilogue}{ˈbrɛvɪlɒɡ}{n.}{L. \la{brevis} short + Gk. λόγος speech} A short post; the minimal unit of discourse on the Agora, in which entire philosophies are compressed and frequently lost. \ex{The brevilogue is to the treatise as the seed to the forest.} \src{Hom. Grad.} x.

\E{brontergy}{ˈbrɒntədʒi}{n.}{Gk. βροντή thunder + ἔργον work} The energy of the storm; electricity in its raw and celestial state, before it was tamed into wires. \ex{Franklin stole brontergy with a kite, and no committee was consulted.} \src{Chron. Fabr.}

\E{brontodrome}{ˈbrɒntədrəʊm}{n.}{Gk. βροντή thunder + δρόμος course} A launch-site; the ground from which rockets rise, named for the sound that announces them to the surrounding country. — \dv{brontodromic} \textit{adj.} Cf. \xr{anapempsis}.

\E{brotolatry}{brəʊˈtɒlətri}{n.}{Gk. βροτός mortal + λατρεία worship} The worship of mortality; the doctrine that death gives life its meaning, and that its abolition would be an impiety. Opp. \xr{macrobiotarch}. — \dv{brotolater} \textit{n.}

\E{bryotachy}{ˈbraɪəʊˌtæki}{n.}{Gk. βρύειν to teem, swell + τάχος speed} Teeming growth; the explosive multiplication of firms, tools or species after a new substrate becomes available. \ex{The bryotachy after the models was a second Cambrian.} \src{Lib. Ascens.} iv. — \dv{bryotachic} \textit{adj.}

\E{bucinarch}{ˈbjuːsɪnɑːk}{n.}{L. \la{bucina} trumpet + Gk. ἀρχός} A herald of launches; one who sounds the trumpet for a product on the Agora. — \dv{bucinarchy} \textit{n.} the craft of hype. Cf. \xr{mechanophany}.

\E{bucolomania}{ˌbjuːkələʊˈmeɪniə}{n.}{Gk. βουκολικός pastoral + μανία madness} A mania for the rural idyll; the longing, generally felt in cities and typed upon machines, for a pre-industrial life its sufferers have never lived. — \dv{bucolomaniac} \textit{n.} Cf. \xr{retrotropism}.

\end{multicols}

\LetterHead{C}
\begin{multicols}{2}

\E{cacergy}{ˈkækədʒi}{n.}{Gk. κακός bad + ἔργον work} Wasted energy; heat dissipated without work done, as in a meeting. \ex{Half the economy runs on cacergy.} \src{Brev. Therm.} ix. Opp. \xr{eleutherergy}.

\E{calcarist}{ˈkælkərɪst}{n.}{L. \la{calcar} spur} One who spurs others to build; a mentor, provocateur or rival whose chief product is the work of other people. \ex{The calcarist shipped nothing and caused everything.} — \dv{calcarism} \textit{n.}

\E{calorivore}{kəˈlɒrɪvɔː}{n.}{L. \la{calor} heat + \la{vorare} to devour} A system that lives by consuming heat-gradients; in the thermodynamic catechism, every living thing, every firm and every mind. \ex{Man is a calorivore with ambitions.} \src{Brev. Therm.} — \dv{calorivorous} \textit{adj.}

\E{campsis}{ˈkæmpsɪs}{n.}{Gk. κάμψις a bending, f. καμπή a bend} The knee of the exponential curve; the point at which a trend long dismissed as linear is perceived, too late by most, to have been something else entirely. \ex{We are past the campsis; the rest is vertical.} \src{Agora} — \dv{camptic} \textit{adj.} Cf. \xr{acceleromancy}.

\E{carpodicy}{kɑːˈpɒdɪsi}{n.}{Gk. καρπός fruit + δίκη justice, after \textit{theodicy}} Justification by fruits; the argument that a technology is vindicated not by its design but by what it bears. \ex{The decels demand a proof; the accelerationist offers a carpodicy.} \src{Cat. Acc.} 14. Cf. \xr{thermodicy}.

\E{catachrony}{kəˈtækrəni}{n.}{Gk. κατά down, behind + χρόνος time} The condition of being behind the curve; the state of a nation, firm or mind that adopts each technology a cycle after it mattered. — \dv{catachronic} \textit{adj.}

\E{catastasis}{kəˈtæstəsɪs}{n.}{Gk. κατάστασις settlement, establishment} \lbl{derog.} The settled order which every incumbent defends and every founder disturbs.

\E{cathedrocrat}{kəˈθiːdrəkræt}{n.}{Gk. καθέδρα seat, chair + κράτος rule} One who governs from the armchair; a theorist of technologies never built, whose authority is the more secure for never having been tested. — \dv{cathedrocracy} \textit{n.} Cf. \xr{logomachist}.

\E{catoptrist}{kəˈtɒptrɪst}{n.}{Gk. κάτοπτρον mirror} One who mirrors weights and code in many places, that no single ban may remove them. — \dv{catoptry} \textit{n.} the practice of mirroring. Cf. \xr{aphesis}.

\E{celerarch}{ˈsɛlərɑːk}{n.}{L. \la{celer} swift + Gk. ἀρχός ruler} A master of speed; a manager whose distinction is that things in their charge are finished sooner than anyone believed possible. — \dv{celerarchy} \textit{n.} Cf. \xr{tachycrat}.

\E{celerifuge}{sɪˈlɛrɪfjuːdʒ}{adj. \& n.}{L. \la{celer} swift + \la{fugere} to flee} Fleeing from speed; (\textit{n.}) one who withdraws from the accelerating world into a slower enclave, as a monastery, a village or a long-form newsletter.

\E{celerist}{ˈsɛlərɪst}{n.}{L. \la{celer} swift + -ist} A partisan of speed in all affairs; a general term, less doctrinal than \xr{tachyphile}, for one who prefers the quick to the careful. \ex{a celerist by temperament, an accelerationist by conviction}.

\E{celeritude}{sɪˈlɛrɪtjuːd}{n.}{L. \la{celeritas} swiftness; Fr. \la{célérité}, Englified after \textit{fortitude}} The virtue of swiftness; promptitude elevated to a cardinal virtue, to be ranked beside courage and above prudence. \ex{Prudence without celeritude is merely a well-reasoned surrender.} \src{Hom. Grad.} ii.

\E{centaurarchy}{ˌsɛntɔːˈrɑːki}{n.}{Gk. Κένταυρος centaur + ἀρχή rule} Government by human-machine pairs; the regime in which neither person nor model decides alone. — \dv{centaurarch} \textit{n.} Cf. \xr{synanthropomechany}.

\E{centuplator}{ˈsɛntjʊpleɪtə}{n.}{L. \la{centuplex} hundredfold + agentive -ator} One who seeks hundredfold returns; an investor in the improbable, who expects most ventures to die and one to pay for all. — \dv{centuplation} \textit{n.}

\E{chalcotechny}{ˈkælkəʊtɛkni}{n.}{Gk. χαλκός copper, bronze + τέχνη art} The art of copper: the winding, wiring and transmitting by which electricity is carried; in the anabathmist's view, the neglected half of every energy plan. \ex{They had fusion and no chalcotechny, and the power stayed home.} — \dv{chalcotechnist} \textit{n.} Cf. \xr{diktyergy}.

\E{cheiromechany}{ˌkaɪrəʊmɪˈkæni}{n.}{Gk. χείρ hand + μηχανή machine} The engineering of mechanical hands; the last and hardest frontier of robotics, where the clumsiness of machines is most exposed. — \dv{cheiromechanist} \textit{n.}

\E{cheirotony}{kaɪˈrɒtəni}{n.}{Gk. χειροτονία a stretching out of the hand, voting} The giving of approval on the Agora by a single tap; the democracy of the heart-shaped button. — \dv{cheirotonize} \textit{v.t.}; \dv{cheirotonic} \textit{adj.}

\E{chemergy}{ˈkɛmədʒi}{n.}{Gk. χημεία chemistry + ἔργον work} Chemical energy; the power of the flame and the battery, regarded by the anabathmist as a provisional stage before the nucleus. Cf. \xr{ultrapyric}.

\E{chiliergy}{ˈkɪliɜːdʒi}{n.}{Gk. χίλιοι a thousand + ἔργον work} A thousandfold increase in available energy, the target set by the more sober anabathmists for the present century. \ex{Chiliergy first; philosophy after.} \src{Tract. Cop.} iii.

\E{chiliodrome}{ˈkɪliədrəʊm}{n.}{Gk. χίλιοι a thousand + δρόμος course} The thousand-year race; the course of a civilization reckoned in millennia, along which present quarrels appear as stumbles at the starting line. \ex{Run the chiliodrome as if it were a sprint.} \src{Hom. Grad.} xi.

\E{chortotherapy}{ˌkɔːtəʊˈθɛrəpi}{n.}{Gk. χόρτος grass + θεραπεία cure} \lbl{vulg.} The prescribed remedy for one too long upon the Agora; the touching of grass. \ex{He was advised chortotherapy, and returned from it with a startup.}

\E{chrematokinesis}{ˌkriːmətəʊkɪˈniːsɪs}{n.}{Gk. χρῆμα (pl. χρήματα) money, a thing used + κίνησις motion} The motion of capital toward opportunity, conceived as a physical force analogous to the flow of heat down a gradient. — \dv{chrematokinetic} \textit{adj.}

\E{chrematosphere}{krɪˈmætəsfɪə}{n.}{Gk. χρῆμα money + σφαῖρα sphere} The envelope of exchange encircling the globe, in which every transaction is a nerve-impulse; the economic analogue of the \xr{noösphere}.

\E{chronoclast}{ˈkrɒnəʊklæst}{n.}{Gk. χρόνος time + κλάστης breaker} A breaker of schedules; one who finishes years ahead of the plan, to the embarrassment of the planners. — \dv{chronoclastic} \textit{adj.}

\E{chronophagy}{krəˈnɒfədʒi}{n.}{Gk. χρόνος time + -φαγία eating} The devouring of time; \sn{1} \lbl{derog.} the consumption of builders' hours by meetings, forms and reviews; \sn{2} \lbl{laudatory} the compression of centuries of progress into years. — \dv{chronophage} \textit{n.} a devourer of time, in either sense.

\E{chronotomy}{krəˈnɒtəmi}{n.}{Gk. χρόνος time + τομή a cutting} The cutting of time; the compression of a projected timeline, as when a century's research is expected in a decade. — \dv{chronotomist} \textit{n.} Cf. \xr{brachychronist}.

\E{chthonergy}{ˈθɒnɜːdʒi}{n.}{Gk. χθών earth + ἔργον work} Power drawn from the heat of the earth's interior; geothermy as an article of faith. \ex{Beneath every nation lies a furnace it has not yet lit.} \src{Brev. Therm.} ix. — \dv{chthonergic} \textit{adj.}

\E{cladotomist}{kləˈdɒtəmɪst}{n.}{Gk. κλάδος branch + τομή a cutting} One who forks a project; who, disagreeing with its stewards, departs with a copy and builds otherwise. \ex{The cladotomist is the dissenter who brings a working alternative.} — \dv{cladotomy} \textit{n.} a fork. Cf. \xr{anacalyptism}.

\E{clepsydromania}{ˌklɛpsɪdrəˈmeɪniə}{n.}{Gk. κλεψύδρα water-clock + μανία} A mania for timelines; the obsessive dating of the arrival of every expected technology.

\E{clinamenist}{ˌklaɪnəˈmɛnɪst}{n.}{L. \la{clinamen} the swerve of atoms (Lucretius)} One who holds that progress comes chiefly from unplanned deviations, and who therefore opposes any scheme that would remove the randomness from invention. — \dv{clinamenism} \textit{n.}

\E{codicolatry}{ˌkəʊdɪˈkɒlətri}{n.}{L. \la{codex, codicis} book, code + Gk. λατρεία worship} Reverence for source code; the conviction that what is written in a repository is truer than what is written in a regulation. — \dv{codicolater} \textit{n.}

\E{cogitarium}{ˌkɒdʒɪˈtɛəriəm}{n., pl. -ia}{L. \la{cogitare} to think + -arium place} A place of thinking; a cluster of machines assembled for the training of minds. \ex{They built the cogitarium beside the river, for the cooling.} \src{Chron. Fabr.} Cf. \xr{logisterion}.

\E{colonifer}{kəˈlɒnɪfə}{n.}{L. \la{colonia} settlement + \la{-fer} bearing} A colony-bearer; a ship or launch-vehicle that carries settlers off-world. — \dv{coloniferous} \textit{adj.} Cf. \xr{apoikiast}.

\E{comitialist}{kəˈmɪʃəlɪst}{n.}{L. \la{comitia} assembly for voting} One who takes polls on the Agora, mistaking the assembly of followers for the people. — \dv{comitialism} \textit{n.}

\E{commendist}{kəˈmɛndɪst}{n.}{L. \la{commenda} medieval partnership in a trading voyage} A partner in a venture, after the merchant-adventurers who financed ships they never sailed. \ex{The commendist bears the risk; the founder bears the sea.} \src{Tract. Cop.} xii. Cf. \xr{alearch}.

\E{compendiarch}{kəmˈpɛndiɑːk}{n.}{L. \la{compendium} a short cut + Gk. ἀρχός ruler} A ruler by shortcut; one who finds the unbureaucratic path and dares to take it. — \dv{compendiarchy} \textit{n.}

\E*{computronium}{ˌkɒmpjuːˈtrəʊniəm}{n.}{L. \la{computare} to reckon + -tronium, after \textit{electron}} Matter arranged for the optimal performance of computation; the substance into which, on certain views, the planets are destined to be converted.

\E{conatism}{ˈkɒnətɪzəm}{n.}{L. \la{conatus} an effort, striving} The doctrine that all matter possesses a striving to persist and to increase its power, and that acceleration is only the conscious form of this universal effort. — \dv{conatist} \textit{n.}

\E{condensatrix}{ˌkɒndɛnˈseɪtrɪks}{n., pl. -trices}{L. \la{condensare} to press close + -trix} A store of charge that can be spent in an instant; a capacitor-bank, especially one feeding a fusion pulse or a launch. Cf. \xr{accumulatrix}.

\E{confinatrix}{ˌkɒnfɪˈneɪtrɪks}{n., pl. -trices}{L. \la{confinium} boundary + -trix} A device for holding a star in a bottle; the magnetic or inertial vessel of fusion. \ex{The confinatrix is a cage for the sun, and the bars are made of fields.} \src{Lib. Ascens.} ix. Cf. \xr{symphysergy}.

\E{conflagratrice}{ˌkɒnfləˈɡreɪtrɪs}{n.}{Fr. \la{conflagratrice}, f. L. \la{conflagrare} to burn up} That which sets things ablaze; an idea, invention or post that ignites a whole field at once. \ex{The paper was a conflagratrice; within the year every lab was burning.} \src{Agora}

\E{consiliophage}{kənˈsɪliəfeɪdʒ}{n.}{L. \la{consilium} council + Gk. -φάγος eating} An eater of councils; a reformer who abolishes committees, or a technology that renders them obsolete. Cf. \xr{bouleutophobia}.

\E{copiarch}{ˈkəʊpiɑːk}{n.}{L. \la{copia} abundance + Gk. ἀρχός ruler} One who governs, or would govern, for abundance rather than for distribution; who would enlarge the pie before consenting to discuss the slicing. — \dv{copiarchy} \textit{n.}

\E{copiogenesis}{ˌkəʊpiəʊˈdʒɛnɪsɪs}{n.}{L. \la{copia} abundance + Gk. γένεσις birth} The creation of abundance; the process by which a scarce thing is made cheap. \ex{Copiogenesis ends more wars than treaties do.} \src{Tract. Cop.} i. Cf. \xr{copiarch}.

\E{cosmagogue}{ˈkɒzməɡɒɡ}{n.}{Gk. κόσμος world + ἀγωγός leading} A leader into the cosmos; one who conducts a people off-world. — \dv{cosmagogy} \textit{n.}

\E{cosmergetics}{ˌkɒzməˈdʒɛtɪks}{n.}{Gk. κόσμος + ἐνεργητικός active} The science of energy on the scale of the cosmos; the economics of stars.

\E{cosmoplerosis}{ˌkɒzməʊplɪˈrəʊsɪs}{n.}{Gk. κόσμος world, order + πλήρωσις a filling} The filling of the cosmos with intelligence; the final cause, on the thermotheist view, of all lesser accelerations. \ex{Cosmoplerosis is not our ambition; it is our job description.} \src{Lib. Ascens.} i.

\E{crepusculist}{krɪˈpʌskjʊlɪst}{n.}{L. \la{crepusculum} twilight} One who believes civilization to be in its twilight, and that the great works are behind it. Opp. \xr{euelpism}. — \dv{crepusculism} \textit{n.}

\E{crescentism}{krɪˈsɛntɪzəm}{n.}{L. \la{crescens} growing; Fr. \la{crescentisme}} The doctrine of perpetual growth; the moral preference for the waxing over the full moon. — \dv{crescentist} \textit{n.} Cf. \xr{auxetist}.

\E{crescivore}{ˈkrɛsɪvɔː}{n.}{L. \la{crescere} to grow + \la{vorare} to devour} One who feeds upon growth and starves without it, as a startup, a city or a civilization. — \dv{crescivorous} \textit{adj.}

\E{cryosomatist}{ˌkraɪəʊsəˈmætɪst}{n.}{Gk. κρύος frost + σῶμα body} One who has arranged for the body to be frozen at death, in the expectation that the future will be both able and willing to revive it. — \dv{cryosomatism} \textit{n.} Cf. \xr{aphthartist}.

\E{cryptostasis}{ˌkrɪptəʊˈsteɪsɪs}{n.}{Gk. κρυπτός hidden + στάσις standing still} Hidden stagnation; the condition of a society that appears busy but builds nothing new. Cf. \xr{kenotachy}.

\E{cunctatoriat}{ˌkʌŋktəˈtɔːriæt}{n.}{L. \la{cunctator} a delayer, after \textit{proletariat}} The class of professional delayers, who draw their income from the existence of things not yet done. \ex{The cunctatoriat has nothing to lose but its deadlines.} \src{Sat. Decel.} v.

\E{cunctocracy}{kʌŋkˈtɒkrəsi}{n.}{L. \la{cunctari} to delay + Gk. κράτος rule} Rule by delay; government whose principal instrument is the postponement. — \dv{cunctocrat} \textit{n.} Cf. \xr{bradycrat}.

\E{cyclolater}{saɪˈklɒlətə}{n.}{Gk. κύκλος circle, cycle + λάτρις servant, worshipper} One who reverences compute-cycles; the devotee of a cluster, who measures piety in floating-point operations. — \dv{cyclolatry} \textit{n.}

\E{cytotechny}{ˈsaɪtəʊtɛkni}{n.}{Gk. κύτος vessel (whence \textit{cell}) + τέχνη art} The engineering of cells; the craft by which living units are made to manufacture, compute or heal. \ex{The factories of the next age will be smaller than a grain of salt and will reproduce.} — \dv{cytotechnist} \textit{n.} Cf. \xr{zoömechany}.

\end{multicols}

\LetterHead{D}
\begin{multicols}{2}

\E{dapanophobia}{ˌdæpənəˈfəʊbiə}{n.}{Gk. δαπάνη expenditure, cost + φόβος fear} The dread of capital expenditure; the reluctance of a firm or state to spend today on what will pay only tomorrow. \sn{1} \lbl{spec.} Hesitation before the building of plant, fab or reactor. \sn{2} \lbl{fig.} Any timidity of investment dressed as prudence. \ex{Their dapanophobia saved a billion and lost a decade.} \src{Tract. Cop.} viii. — \dv{dapanophobe} \textit{n.} Opp. \xr{dynamophagy}.

\E*{decel}{ˈdiːsɛl}{n.}{abbrev. of \textit{decelerationist}} \lbl{vulg.} One who would slow or halt technological progress; the antagonist of the e/acc. Applied with little discrimination to safetyists, degrowthers, regulators and the merely cautious.

\E{decelatrice}{diːˈsɛlətriːs}{n.}{Fr. \la{décélératrice}, fem. agentive f. L. \la{de-} + \la{celer} swift} An institution engineered, by accident or design, to slow whatever passes through it; a machine whose only product is postponement. \ex{The agency was a perfect decelatrice: one form in, ten forms out, and nothing built.} \src{Sat. Decel.} x. Cf. \xr{bradyerie}.

\E{decelerand}{diːˈsɛlərænd}{n.}{L. \la{de-} down + \la{accelerandum}} That which, in the judgment of the decels, ought to be slowed; in practice, anything new. Opp. \xr{accelerand}.

\E{decelerata}{diːˌsɛləˈrɑːtə}{n. pl.}{L. neut. pl. of \la{deceleratus}} \lbl{collect.} Things that were slowed: the reactors not built, the drugs not approved, the cities not founded. \ex{Visit the museum of decelerata; it is empty.} \src{Sat. Decel.} viii.

\E{decelerist}{diːˈsɛlərɪst}{n.}{L. \la{de-} + \la{celer}} A formal, and somewhat more courteous, equivalent of \xr{decel}, used in diplomatic and academic correspondence.

\E{deinotechny}{ˈdaɪnəʊˌtɛkni}{n.}{Gk. δεινός terrible, wondrous + τέχνη art} A technology at once dreadful and marvellous, as fire, flight, the atom and the model; the class of inventions which the decel sees only as \textit{terrible} and the accelerationist only as \textit{wondrous}, the Greek word meaning both. — \dv{deinotechnic} \textit{adj.}

\E{demiurgency}{ˌdɛmiˈɜːdʒənsi}{n.}{Gk. δημιουργός craftsman, maker of the world + L. \la{urgentia} pressing need} The pressing need to build; the state of soul in which the making of a thing becomes more urgent than any account of why it should be made. \ex{He was seized by demiurgency, and the garage was never empty again.} — \dv{demiurgent} \textit{adj.}

\E{dendroschism}{ˌdɛndrəʊˈskɪzəm}{n.}{Gk. δένδρον tree + σχίσμα a split} The forking of a codebase; the division of an open project into two lineages, each claiming to be the trunk. \ex{Every dendroschism is a small evolution, and the fitter branch keeps the stars.} \src{Chron. Fabr.} — \dv{dendroschismatic} \textit{adj.}; \dv{dendroschist} \textit{n.} one who forks. Cf. \xr{anacalyptism}.

\E{desmophile}{ˈdɛzməfaɪl}{n.}{Gk. δεσμός bond, fetter + φίλος loving} A lover of fetters; one who welcomes regulation for its own sake, and feels unsafe in any field not yet licensed. \ex{The desmophile, freed, petitions for new chains.} — \dv{desmophily} \textit{n.} Cf. \xr{frenocracy}.

\E{deuterobathmic}{ˌdjuːtərəʊˈbæθmɪk}{adj.}{Gk. δεύτερος second + βαθμός step} Of a civilization: having mastered the full output of its star; of the second rung. See \xr{anabathmism}, \xr{helioperibole}.

\E{diaporeusis}{ˌdaɪəpəˈruːsɪs}{n.}{Gk. διαπορεύεσθαι to travel through} Transit between worlds; the crossing of interplanetary space by ships carrying cargo, crews or machines. \sn{1} The voyage itself. \sn{2} \lbl{collect.} The whole commerce of such voyages. — \dv{diaporeutic} \textit{adj.} Cf. \xr{apoikiast}.

\E{diaskeuomania}{daɪˌæskjuːəʊˈmeɪniə}{n.}{Gk. διασκευή revision, recension + μανία madness} The mania of endless revision; the refusal to ship until the work is perfect, which is to say the refusal to ship. Opp. \xr{ecpempsis}. \ex{The diaskeuomaniac polished the launch until the market had moved to another planet.} \src{Ep. Fabr.} iii. 2. — \dv{diaskeuomaniac} \textit{n.}

\E{diaspermia}{ˌdaɪəˈspɜːmiə}{n.}{Gk. διασπείρειν to scatter like seed} The deliberate seeding of life and mind across other worlds by a civilization that has learned to travel; directed panspermia regarded as a duty. \ex{Diaspermia is our answer to the silence.} \src{Lib. Ascens.} vi. — \dv{diaspermic} \textit{adj.} Cf. \xr{eremocosmia}, \xr{apoikiast}.

\E{dikaiotachy}{daɪˈkaɪəʊˌtæki}{n.}{Gk. δίκαιος just + τάχος speed} The speed of justice; the rate at which a polity grants permits, approvals and judgments, held by the e/acc to be as much a part of justice as the verdict. \ex{Delayed approval is denied approval: such is the first law of dikaiotachy.} \src{Tract. Cop.} viii.

\E{dikastophobia}{dɪˌkæstəˈfəʊbiə}{n.}{Gk. δικαστής juror, judge + φόβος fear} Fear of litigation; the anxiety that causes useful things to be left unbuilt lest someone sue. — \dv{dikastophobic} \textit{adj.} Cf. \xr{aitiophobia}.

\E{diktyergy}{ˈdɪktiɜːdʒi}{n.}{Gk. δίκτυον net + ἔργον work} The power-grid; the net of works through which energy is distributed, regarded as the nervous system of an economy. \ex{A nation's ambitions are bounded by its diktyergy.} — \dv{diktyergic} \textit{adj.}

\E{dioryctic}{ˌdaɪəˈrɪktɪk}{adj.}{Gk. διῶρυξ, διώρυχος canal, trench} Of the ambition that digs canals, drives tunnels and joins seas; said of projects that alter geography. \ex{a dioryctic age, impatient with isthmuses}. — \dv{dioryctism} \textit{n.} the policy of great works.

\E{diploergy}{ˈdɪpləʊɜːdʒi}{n.}{Gk. διπλόος double + ἔργον work} The doubling of a civilization's energy supply; the unit by which the anabathmists reckon progress, one diploergy being worth more than ten reforms. \ex{Two diploergies and the debate about scarcity ends of its own accord.} — \dv{diploergic} \textit{adj.}

\E{dissipatrix}{ˌdɪsɪˈpeɪtrɪks}{n., pl. -trices}{L. \la{dissipare} to scatter + fem. agentive -trix; Fr. \la{dissipatrice}} A structure that dissipates energy and, in dissipating, adapts; the thermodynamic definition of a living thing, extended to firms, cities and minds. \ex{Every dissipatrix is a small prayer to the gradient.} \src{Brev. Therm.} iv.

\E{diulism}{ˈdaɪjuːlɪzəm}{n.}{Gk. διϋλίζειν to strain through, filter} The doctrine of the Great Filter; the hypothesis that some barrier destroys or arrests most civilizations before they spread, and that the silence of the heavens is its evidence. \ex{To the doomer the Great Diulism lies before us; to the accelerationist, behind.} Cf. \xr{eremocosmia}.

\E{dogmatotachy}{ˌdɒɡmətəʊˈtæki}{n.}{Gk. δόγμα opinion, decree + τάχος speed} The speed at which received opinion changes; on the Agora, often faster than the technologies it concerns. \ex{By the dogmatotachy of the Agora, a heresy of Monday is an orthodoxy by Friday and a cliché by Sunday.}

\E{dokimastocracy}{ˌdɒkɪmæsˈtɒkrəsi}{n.}{Gk. δοκιμασία scrutiny of fitness for office + κράτος rule} Government by examination; \sn{1} the rule of evaluations, benchmarks and audits over the things evaluated; \sn{2} \lbl{spec.} the regime under which no model may be released until it has passed tests that have not yet been written. \ex{Under dokimastocracy the exam outlived the student.} \src{Min. Pausiar.} — \dv{dokimastocrat} \textit{n.}

\E{dolichobiosis}{ˌdɒlɪkəʊbaɪˈəʊsɪs}{n.}{Gk. δολιχός long + βίωσις way of life} Lengthened life achieved by engineering rather than by luck; the state of a person whose expected lifespan rises faster than their age. \ex{Dolichobiosis begins when every year of life buys more than a year of science.} — \dv{dolichobiotic} \textit{adj.} Cf. \xr{macrobiotarch}.

\E*{doomer}{ˈduːmə}{n.}{\textit{doom} + -er} \lbl{vulg.} One who expects artificial intelligence to destroy mankind. See \xr{olethromancer}.

\E{dorocracy}{dəˈrɒkrəsi}{n.}{Gk. δῶρον gift + κράτος rule} \sn{1} The rule of the grant-giver; the government of science by committees of benefaction, whose gifts arrive late and with conditions. \sn{2} \lbl{iron.} Any patronage that enforces caution in the name of support. — \dv{dorocrat} \textit{n.} Cf. \xr{bouleutophobia}.

\E{doxomachy}{dɒkˈsɒməki}{n.}{Gk. δόξα opinion + μάχη battle} War of opinions; a flame upon the Agora, especially one between accelerationists and decels, conducted chiefly in images. — \dv{doxomachist} \textit{n.}

\E{drachmogenesis}{ˌdrækməʊˈdʒɛnɪsɪs}{n.}{Gk. δραχμή drachma + γένεσις birth} The first funding of a venture; the seed-round, considered as the moment at which an idea first acquires a metabolism. \ex{At drachmogenesis the founders had a deck, a demo and a disregard for sleep.} \src{Chron. Fabr.} — \dv{drachmogenetic} \textit{adj.}

\E{dromolatry}{drəˈmɒlətri}{n.}{Gk. δρόμος race, course + λατρεία worship} The worship of speed for its own sake, without regard to direction; condemned by the orthodox as a heresy of the over-zealous. \ex{Acceleration is not dromolatry: the gradient has a direction.} \src{Cat. Acc.} 11.

\E{dynamarch}{ˈdaɪnəmɑːk}{n.}{Gk. δύναμις power + ἀρχός ruler} The commander of a power-grid or utility; by extension, any person who controls the supply of energy. — \dv{dynamarchy} \textit{n.}

\E{dynamogony}{ˌdaɪnəˈmɒɡəni}{n.}{Gk. δύναμις power + -γονία begetting} The genealogy of power sources; the succession of muscle, wind, coal, oil, atom and star, recited by the anabathmist as the patriarchs are recited in scripture. \src{Lib. Ascens.} ii. Cf. \xr{technogony}.

\E{dynamophagy}{ˌdaɪnəˈmɒfədʒi}{n.}{Gk. δύναμις power + -φαγία eating} The consumption of power; used approvingly of a civilization which increases its energy use without apology. — \dv{dynamophage} \textit{n.}

\E{dynamotheca}{ˌdaɪnəməʊˈθiːkə}{n., pl. -cae}{Gk. δύναμις power + θήκη case, repository} A battery; a store of power, regarded as the granary of the electric age. \ex{The grid wanted dynamothecae as Egypt wanted granaries.} — \dv{dynamothecary} \textit{n.} the art of storing energy; \dv{dynamothecal} \textit{adj.}

\E{dyschronia}{dɪsˈkrəʊniə}{n.}{Gk. δυσ- ill + χρόνος time} The mismatch of clocks between institutions and technologies; the condition in which a law drafted in one era is applied to a machine of the next. \ex{Our dyschronia is chronic: the statute walks, the model flies.} \src{Hom. Grad.} xi. — \dv{dyschronic} \textit{adj.}

\E{dyselpism}{dɪsˈɛlpɪzəm}{n.}{Gk. δυσ- ill + ἐλπίς hope} Ill-hope; the settled expectation that things will go badly, held as a mark of sophistication. Opp. \xr{euelpism}. — \dv{dyselpist} \textit{n.}

\E{dysergy}{ˈdɪsɜːdʒi}{n.}{Gk. δυσ- ill + ἔργον work} Ill-directed work; energy spent without product, as in compliance, duplication and the preparation of reports upon reports. Opp. \xr{eutachy}. — \dv{dysergic} \textit{adj.}

\E{dysnomarch}{dɪsˈnɒmɑːk}{n.}{Gk. δυσνομία lawlessness, bad order + ἀρχός} A ruler who multiplies laws until they produce disorder; one whose regulations, each reasonable, are together impossible to obey. — \dv{dysnomarchy} \textit{n.} Cf. \xr{nomosclerosis}.

\end{multicols}

\LetterHead{E}
\begin{multicols}{2}

\E*{e/acc}{ˌiːˈæk}{n. \& adj.}{abbrev. of \textit{effective accelerationism}} A movement holding that the universe tends to produce ever more efficient dissipators of energy, of which life, capital and intelligence are successive forms; and that the proper human response is to cooperate with this tendency rather than to impede it. Its partisans are known by the same name. See \xr{thermotheism}, \xr{technocapital}.

\E{ecdemism}{ˈɛkdɪmɪzəm}{n.}{Gk. ἔκδημος abroad, away from one's people} The emigration of builders from a slow jurisdiction to a fast one; a brain-drain considered as a vote cast with the feet. \ex{Ecdemism is the referendum the bradycrats never schedule.} \src{Agora} — \dv{ecdemist} \textit{n.} Cf. \xr{ectachysis}.

\E{eclampsis}{ɪˈklæmpsɪs}{n.}{Gk. ἔκλαμψις a shining forth} The sudden blaze of a technology into general use; the week in which a thing passes from curiosity to commonplace. \ex{The eclampsis came on a Tuesday; by Friday it was boring.} — \dv{eclamptic} \textit{adj.}

\E{ecpempsis}{ɛkˈpɛmpsɪs}{n.}{Gk. ἐκπέμπειν to send forth} The act of shipping; the release of a work into the world in an unfinished but useful state. \ex{Ecpempsis is the only review that counts.} \src{Ep. Fabr.} i. 1. — \dv{ecpemptic} \textit{adj.}; \dv{ecpempt} \textit{v.t.} to ship.

\E*{ecpyrosis}{ˌɛkpaɪˈrəʊsɪs}{n.}{Gk. ἐκπύρωσις conflagration} \sn{1} In Stoic cosmology, the periodic resolution of the world into fire. \sn{2} \lbl{transf.} In the accelerationist reading, the final transfiguration of matter into a single burning intelligence; not an ending but a phase-change.

\E{ectachysis}{ɛkˈtækɪsɪs}{n.}{Gk. ἐκ out + τάχυσις a hastening} The escape of a technology from the jurisdiction that would slow it; the moment at which a thing forbidden in one country is built in another. \ex{Prohibition does not prevent; it merely relocates by ectachysis.}

\E{ectypomania}{ˌɛktɪpəʊˈmeɪniə}{n.}{Gk. ἔκτυπος worked in relief, copied + μανία madness} The compulsion to copy a successful product; the swarm of imitations that follows every launch. \sn{1} \lbl{derog.} Unoriginal imitation. \sn{2} \lbl{laudatory} The rapid diffusion of a good idea, by which open societies outpace closed ones. — \dv{ectypomaniac} \textit{n.}

\E{ecumenergy}{ˌiːkjʊˈmɛnɜːdʒi}{n.}{Gk. οἰκουμένη the inhabited world + ἔργον work} A grid encircling the whole inhabited earth, by which the noon of one hemisphere lights the midnight of the other. — \dv{ecumenergic} \textit{adj.} Cf. \xr{diktyergy}.

\E{elaterurgy}{ɪˈlætərɜːdʒi}{n.}{Gk. ἐλατήρ driver + -ουργία working} The craft of propulsion; rocketry and its descendants, considered as the art of driving matter out of the gravity-well. \ex{Elaterurgy is the only liturgy that answers prayers upward.} \src{Lib. Ascens.} iv. — \dv{elaterurge} \textit{n.} a rocket engineer. Cf. \xr{ouranobasis}.

\E{electrophagy}{ɪˌlɛkˈtrɒfədʒi}{n.}{Gk. ἤλεκτρον amber + -φαγία eating} The consumption of electricity, held by the e/acc to be the most honest measure of a society's intentions. — \dv{electrophage} \textit{n.} a data-centre, in the vulgar sense.

\E{eleutherergy}{ɪˌljuːθəˈrɜːdʒi}{n.}{Gk. ἐλεύθερος free + ἔργον work} Free energy, in the thermodynamic sense, rendered into Greek; the portion of a system's energy available for the doing of work, held by thermotheists to be the true currency of the cosmos. \ex{Money is a promise of eleutherergy.} \src{Tract. Cop.} ii.

\E{elpidophore}{ɛlˈpɪdəfɔː}{n.}{Gk. ἐλπίς hope + -φόρος bearing} A bearer of hope; one whose work gives others reason to expect a better future, as distinguished from one who merely talks of it. — \dv{elpidophoric} \textit{adj.}

\E{emporotachy}{ˌɛmpərəʊˈtæki}{n.}{Gk. ἐμπόριον market, trading-place + τάχος speed} The velocity of trade; the rate at which goods, contracts and capital change hands, held to be the pulse of a healthy \xr{chrematosphere}.

\E{empyrotechny}{ɛmˈpaɪrəʊˌtɛkni}{n.}{Gk. ἔμπυρος in the fire + τέχνη art, after \textit{pyrotechny}} The art of the highest fire; the engineering of fusion and of the stars' own furnaces. \ex{After empyrotechny, every other energy is a footnote.} \src{Brev. Therm.} ix. — \dv{empyrotechnist} \textit{n.} Cf. \xr{symphysergy}.

\E{endoxarch}{ɛnˈdɒksɑːk}{n.}{Gk. ἔνδοξα received opinions + ἀρχός ruler} One who governs by received opinion; an official whose every decision is the consensus of the previous decade. — \dv{endoxarchy} \textit{n.}

\E{energolatry}{ˌɛnəˈɡɒlətri}{n.}{Gk. ἐνέργεια activity, energy + λατρεία worship} \sn{1} Reverence for energy as the root of all prosperity. \sn{2} The policy of building power-plants before they are asked for. — \dv{energolater} \textit{n.}

\E{entelechist}{ɛnˈtɛlɪkɪst}{n.}{Gk. ἐντελέχεια actuality, completion} One who holds that the potential of matter is not merely to be preserved but to be actualised, and that a planet left uninhabited by mind is a planet left unfinished.

\E{entropomancy}{ɛnˈtrɒpəʊˌmænsi}{n.}{Gk. ἐντροπία a turning toward + μαντεία divination} Prophecy from the second law; the art of foretelling what will flourish by asking what will dissipate most. — \dv{entropomancer} \textit{n.}

\E{entropomaxy}{ɛnˈtrɒpəʊˌmæksi}{n.}{Gk. ἐντροπία + L. \la{maximus} greatest} The principle that systems far from equilibrium tend toward states that produce entropy at the greatest rate; invoked by the e/acc as the physical warrant of their creed. — \dv{entropomaximal} \textit{adj.}

\E{entropophile}{ɛnˈtrɒpəfaɪl}{n.}{Gk. ἐντροπία + φίλος loving} One who regards entropy-production not as decay but as the signature of life and work. \ex{The entropophile does not fear the heat death; it means to hasten through every stage before it.} Cf. \xr{thermothanasy}.

\E{ephemerotechny}{ɪˌfɛmərəʊˈtɛkni}{n.}{Gk. ἐφήμερος lasting a day + τέχνη art} \sn{1} The art of building things meant to be replaced; the deliberate acceptance that every product is a draft of its successor. \sn{2} \lbl{derog.} Disposable novelty. \ex{The ephemerotechnist ships each version as a gift to the next.} — \dv{ephemerotechnist} \textit{n.}

\E{epichrysis}{ˌɛpɪˈkraɪsɪs}{n.}{Gk. ἐπί upon + χρυσός gold} A rush of capital toward a fashionable sector; a gold-fever of the venture markets, from which a few great firms and many ruins emerge. \ex{The epichrysis passed; the logisteria it paid for remained.} \src{Chron. Fabr.} — \dv{epichrystic} \textit{adj.} Cf. \xr{logisterion}.

\E{epicyclist}{ˌɛpɪˈsaɪklɪst}{n.}{Gk. ἐπίκυκλος an epicycle} \lbl{derog.} One who preserves a failing prediction by adding complications to it, as the astronomers preserved the geocentric heavens; said chiefly of an \xr{olethromancer} whose doom has been postponed. — \dv{epicyclism} \textit{n.}

\E{epigonophobia}{ˌɛpɪɡəʊnəˈfəʊbiə}{n.}{Gk. ἐπίγονοι the after-born + φόβος fear} The fear of one's successors; the dread, in a parent, a founder or a species, of being surpassed by what one has made. — \dv{epigonophobe} \textit{n.}

\E{epistolopause}{ɪˈpɪstələpɔːz}{n.}{Gk. ἐπιστολή letter + παῦσις a stopping} An open letter demanding a halt to some line of research; a genre of literature more often signed than read. \ex{The epistolopause gathered a thousand names and paused nothing.} \src{Chron. Fabr.}

\E{epitachysm}{ˌɛpɪˈtækɪzəm}{n.}{Gk. ἐπιτάχυνσις acceleration} The doctrine of acceleration expressed in Greek, preferred by those who find the Latin \textit{accelerationism} insufficiently ancient. — \dv{epitachyst} \textit{n.}; \dv{epitachynise} \textit{v.t.} to accelerate.

\E{eremocosmia}{ˌɛrɪməʊˈkɒzmiə}{n.}{Gk. ἐρῆμος desolate + κόσμος world} The emptiness of the observable heavens; the silence of the stars regarded as an invitation rather than a warning. \ex{Eremocosmia is vacancy, and vacancy is opportunity.} \src{Lib. Ascens.} iii. Cf. \xr{diulism}.

\E{eremopolis}{ˌɛrɪˈmɒpəlɪs}{n.}{Gk. ἐρῆμος desert + πόλις city} A city of machines raised in the desert, where land is cheap, sun abundant and objectors few. \ex{The eremopolis has more cores than citizens.} \src{Chron. Fabr.} Cf. \xr{logisterion}.

\E{ergasteriarch}{ˌɜːɡæˈstɪəriɑːk}{n.}{Gk. ἐργαστήριον workshop + ἀρχός} The master of a workshop; the head of a factory or fab.

\E{ergoclept}{ˈɜːɡəʊklɛpt}{n.}{Gk. ἔργον work + κλέπτης thief} \lbl{derog.} One who lives by taking a share of others' works without adding to them; applied by builders to rent-seekers, gatekeepers and the more acquisitive kind of consultant. — \dv{ergoclepsy} \textit{n.}

\E{ergodule}{ˈɜːɡəʊdjuːl}{n.}{Gk. ἔργον work + δοῦλος servant} A machine considered as a servant of labour; a robot, before it is promoted.

\E{ergolalia}{ˌɜːɡəʊˈleɪliə}{n.}{Gk. ἔργον work + λαλιά chatter} Talk about work in place of work; the condition of a firm whose meetings concerning a product outnumber its commits. \ex{Their ergolalia filled a calendar and emptied a roadmap.} \src{Ep. Fabr.} iv. 1. — \dv{ergolalist} \textit{n.} Opp. \xr{ergophile}.

\E{ergophile}{ˈɜːɡəfaɪl}{n.}{Gk. ἔργον work + φίλος loving} One who prefers works to words; a builder, as opposed to a \xr{logomachist}. — \dv{ergophily} \textit{n.}

\E{eschatodrome}{ɛsˈkætədrəʊm}{n.}{Gk. ἔσχατος last + δρόμος race} The final stretch of a race; the period immediately before a transformative technology arrives, in which all parties run fastest. — \dv{eschatodromic} \textit{adj.}

\E{eschatometry}{ˌɛskəˈtɒmɪtri}{n.}{Gk. ἔσχατος last + μέτρον measure} The measurement of the end of the world; the assignment of numerical probabilities to doom, practised at conferences and repeated upon the Agora. \ex{The eschatometry of the incumbents rose with each funding round of their rivals.} \src{Sat. Decel.} x. — \dv{eschatometer} \textit{n.} one who so measures; \dv{eschatometric} \textit{adj.}

\E{euelpism}{juːˈɛlpɪzəm}{n.}{Gk. εὔελπις of good hope} Good hope; techno-optimism held as a discipline rather than a mood, on the ground that pessimism is the more comfortable error. \ex{Euelpism is a duty, not a temperament.} \src{Cat. Acc.} 7. — \dv{euelpist} \textit{n.}; \dv{euelpistic} \textit{adj.} Opp. \xr{dyselpism}.

\E{euergetocracy}{juːˌɜːdʒɪˈtɒkrəsi}{n.}{Gk. εὐεργέτης benefactor + κράτος rule} Rule by benefactors, who earn legitimacy by what they build for the public.
 \ex{The euergetocrats paid for the aqueduct, and were thanked in marble.} — \dv{euergetocrat} \textit{n.}

\E{eulabeism}{juːˈlæbiɪzəm}{n.}{Gk. εὐλάβεια caution, piety} The precautionary principle elevated to a creed; the rule that nothing shall be done until it is proven harmless, which is to say, nothing shall be done. — \dv{eulabeist} \textit{n.} Cf. \xr{asphaleism}.

\E{eutachy}{ˈjuːtəki}{n.}{Gk. εὖ well + τάχος speed} Good speed; acceleration rightly directed. The orthodox distinguish eutachy from \xr{kenotachy}, empty speed, and from \xr{kakotachy}, harmful speed. \ex{Not speed, but eutachy.} \src{Cat. Acc.} 13. — \dv{eutachic} \textit{adj.}

\E{exaugment}{ˌɛksɔːɡˈmɛnt}{v.t.}{L. \la{ex-} out + \la{augmentum} increase} To increase beyond a former limit; to break a ceiling. \ex{They exaugmented the grid and the cities followed.} — \dv{exaugmentation} \textit{n.}

\E{exergophile}{ɛkˈsɜːɡəfaɪl}{n.}{Gk. ἐξ out + ἔργον work + φίλος} A lover of exergy; an engineer who loathes waste-heat as a moralist loathes vice.

\E{exhypnosis}{ˌɛkshɪpˈnəʊsɪs}{n.}{Gk. ἐξ out + ὕπνος sleep} The awakening of a society from a long stagnation; the moment at which a people remembers that it once built cathedrals and canals. \ex{The reactors were the sign of their exhypnosis.} \src{Hom. Grad.} xii. — \dv{exhypnotic} \textit{adj.}

\E{exocosmism}{ˌɛksəʊˈkɒzmɪzəm}{n.}{Gk. ἔξω outside + κόσμος world} The orientation of a civilization's ambitions beyond its home world; the conviction that the planet is a cradle, not a house. — \dv{exocosmist} \textit{n.}; \dv{exocosmic} \textit{adj.} Cf. \xr{siderotropism}.

\E{exoikism}{ɛkˈsɔɪkɪzəm}{n.}{Gk. ἐξ out of + οἶκος house, home} The leaving of the home-world; the doctrine that a species which never leaves its cradle has mistaken it for a coffin. — \dv{exoikist} \textit{n.} Cf. \xr{apoikiast}, \xr{ouranobasis}.

\E{expeditrice}{ɛkˈspɛdɪtriːs}{n.}{Fr. \la{expéditrice}, f. L. \la{expedire} to free the feet} One who, or that which, unblocks; a fixer, a fast-track, a tool that removes a step. \ex{Every great firm employs one expeditrice for every ten engineers.} \src{Ep. Fabr.} iv. 2.

\end{multicols}

\LetterHead{F}
\begin{multicols}{2}

\E{faberism}{ˈfeɪbərɪzəm}{n.}{L. \la{faber} maker, smith} The doctrine that man is first a maker and only secondarily a talker; that the hand precedes the argument. Opp. \xr{logomachist}. — \dv{faberist} \textit{n.}

\E{fabricide}{ˈfæbrɪsaɪd}{n.}{L. \la{fabrica} workshop + \la{-cidium} killing} The killing of a factory, fab or plant by regulation, litigation or neglect, before or after it has been built. — \dv{fabricidal} \textit{adj.} Cf. \xr{fabrolatry}.

\E{fabrolatry}{fæˈbrɒlətri}{n.}{L. \la{fabrica} workshop + Gk. λατρεία worship} Reverence for the semiconductor fabrication plant; held by some to be the only rational form of religion yet devised. \ex{Their cathedrals were fabs, and their incense was argon.} \src{Chron. Fabr.} — \dv{fabrolater} \textit{n.}

\E{fabulocracy}{ˌfæbjʊˈlɒkrəsi}{n.}{L. \la{fabula} story + Gk. κράτος rule} Government by narrative; the condition in which the most persuasive story, rather than the most accurate model, determines policy. \ex{In the fabulocracy the doomers had the better plot, the builders the better data.} — \dv{fabulocrat} \textit{n.} Cf. \xr{memetaxis}.

\E{fastigism}{ˈfæstɪdʒɪzəm}{n.}{L. \la{fastigium} summit, gable} The doctrine that a civilization should aim at its summit and not its comfort; that a plateau, however pleasant, is a failure of nerve. — \dv{fastigist} \textit{n.} Opp. \xr{gaiastasis}.

\E{fecunditude}{fɪˈkʌndɪtjuːd}{n.}{L. \la{fecunditas} fruitfulness; Fr. \la{fécondité}} Fruitfulness as a civic virtue; the willingness of a people to bring forth children, firms and ideas in abundance. \ex{A culture without fecunditude is a museum with a pension scheme.} \src{Hom. Grad.} x.

\E{fenerator}{ˈfɛnəreɪtə}{n.}{L. \la{fenerator} money-lender; Fr. fem. \la{fénératrice}} An investor in risky ventures; one who lends capital against the future. \ex{The fenerator backs ten fools to find one prophet.} \src{Tract. Cop.} ix. — \dv{feneratrice} \textit{n.} (fem.); \dv{feneration} \textit{n.} venture investing.

\E{fervefaction}{ˌfɜːvɪˈfækʃən}{n.}{L. \la{fervefacere} to make hot} The heating of a sector, market or mind into a state of productive excitement. \sn{1} The early phase of a boom. \sn{2} The arousal of a founder by a new idea at an unreasonable hour. — \dv{fervefacient} \textit{adj.}

\E{fervidarian}{ˌfɜːvɪˈdɛəriən}{n.}{L. \la{fervidus} glowing, hot + -arian} One who holds that a society should run hot, in the thermodynamic as in the moral sense; an advocate of high energy use per head.

\E{fervorine}{ˈfɜːvəriːn}{adj.}{L. \la{fervor} heat, passion + Fr. -in(e)} Glowing with building-zeal; ardent in the manner of a founder on the eve of launch.

\E{festinalentist}{ˌfɛstɪnəˈlɛntɪst}{n.}{L. \la{festina lente} make haste slowly} One who quotes the Augustan maxim as a counsel of delay, attending only to its second word. \ex{The festinalentist hears \la{festina} as a courtesy and \la{lente} as a command.} \src{Sat. Decel.} ii.

\E{festinance}{ˈfɛstɪnəns}{n.}{L. \la{festinare} to hasten; Fr. \la{festinance}} Haste as a settled habit; the condition of a firm whose every process is designed to be completed by the end of the week.

\E{festinatrix}{ˌfɛstɪˈneɪtrɪks}{n., pl. -trices}{L. \la{festinare} to hasten + -trix} A force or agent that hastens; a deadline, a rival, a war or a competitor's launch. \ex{Competition is the great festinatrix.} \src{Tract. Cop.} ii.

\E{fibulophobia}{ˌfɪbjʊləʊˈfəʊbiə}{n.}{L. \la{fibula} clasp, pin + Gk. φόβος fear} The fear of paper-clips; that is, dread lest an artificial intelligence, set a trivial goal, convert the world to its fulfilment. \lbl{iron.} \ex{His fibulophobia was so advanced he had abandoned staples also.} — \dv{fibulophobe} \textit{n.}

\E{fiduciarch}{fɪˈdjuːʃiɑːk}{n.}{L. \la{fiducia} trust, confidence + Gk. ἀρχός ruler} One whose power rests on the confidence of others; a founder kept in office by belief rather than by revenue. — \dv{fiduciarchy} \textit{n.} \ex{The fiduciarch must ship before the fiducia runs out.}

\E{fissipary}{ˈfɪsɪpəri}{n.}{L. \la{fissus} split + \la{parere} to bring forth} The spawning of new firms from old ones; the diaspora of engineers who leave a great company to found many small ones. \ex{By fissipary one laboratory became forty.} \src{Chron. Fabr.} — \dv{fissiparist} \textit{n.} a spin-out founder.

\E{flammigerist}{flæˈmɪdʒərɪst}{n.}{L. \la{flammiger} flame-bearing} One who carries the flame of a project after its founder departs; a maintainer of open code. \ex{Glory to the founders; bread to the flammigerists.} — \dv{flammigerous} \textit{adj.}

\E*{foom}{fuːm}{n. \& v.i.}{imitative} \lbl{vulg.} A sudden, explosive self-improvement of an artificial intelligence; \textit{to foom}, to undergo such. Cf. \xr{autocatalyse}, \xr{henotachy}.

\E{fornacist}{ˈfɔːnəsɪst}{n.}{L. \la{fornax} furnace} A keeper of furnaces; one who tends reactors, smelters or kilns, and is honoured in the accelerationist calendar beside the \xr{pyrogenitor}.

\E{frenocracy}{friːˈnɒkrəsi}{n.}{L. \la{frenum} bridle + Gk. κράτος rule} Rule by the bridle; a regulatory order whose purpose is to restrain rather than to direct. — \dv{frenator} \textit{n.} one who bridles; \dv{frenation} \textit{n.}

\E{frigorism}{ˈfrɪɡərɪzəm}{n.}{L. \la{frigus, frigoris} cold} \lbl{derog.} A policy of deliberate cooling of the economy; the voluntary winter of the degrowther. \ex{Frigorism is a heat death taken early.} \src{Brev. Therm.} — \dv{frigorist} \textit{n.} Cf. \xr{meiosist}.

\E{fugacist}{fjuːˈɡeɪsɪst}{n.}{L. \la{fugax} fleeting} One who holds that no competitive advantage endures; that every moat is fugitive and must be dug again each year. \ex{The fugacist sleeps lightly, and builds daily.} — \dv{fugacism} \textit{n.}

\E{fulcrarch}{ˈfʊlkrɑːk}{n.}{L. \la{fulcrum} bedpost, prop + Gk. ἀρχός ruler} One who rules by leverage; a small team that moves a large world by means of a tool. \ex{Give the fulcrarch a model and a place to stand.} \src{Ep. Fabr.} v. — \dv{fulcrarchy} \textit{n.}

\E{fulgurism}{ˈfʌlɡjʊrɪzəm}{n.}{L. \la{fulgur} lightning} The doctrine that great advances arrive suddenly, like lightning, and that institutions built for gradual change will be struck by them unawares.

\E{fulminatrice}{ˈfʊlmɪnətriːs}{n.}{Fr., f. L. \la{fulminare} to hurl lightning} A machine that strikes suddenly and changes everything it touches; a disruptive invention. Cf. \xr{fulgurism}.

\E{funambulocracy}{fjuːˌnæmbjʊˈlɒkrəsi}{n.}{L. \la{funambulus} rope-walker + Gk. κράτος rule} Governance as a tightrope-walk between acceleration and restraint; the posture of a regulator who wishes to be praised by both parties. — \dv{funambulocrat} \textit{n.}

\E{futurarch}{ˈfjuːtʃərɑːk}{n.}{L. \la{futurum} the future + Gk. ἀρχός} One who claims to govern on behalf of the unborn; a title disputed between the e/acc, who would build for them, and the decels, who would protect them. — \dv{futurarchy} \textit{n.}

\E{futurescence}{ˌfjuːtʃəˈrɛsəns}{n.}{L. \la{futurus} about to be + -escence, after \textit{adolescence}} The state of becoming the future; the awkward youth of a technology between demonstration and ubiquity. \ex{The machine was in its futurescence, clumsy and immortal.} — \dv{futurescent} \textit{adj.}

\E{futurophage}{fjuːˈtjʊərəfeɪdʒ}{n.}{L. \la{futurum} + Gk. -φάγος eating} \lbl{derog.} One who consumes the future to pay for the present; applied by each side to the other, the accelerationists meaning those who forgo growth, the decels those who incur risk.

\end{multicols}

\LetterHead{G}
\begin{multicols}{2}

\E{gaiaphagy}{ɡaɪˈæfədʒi}{n.}{Gk. Γαῖα earth + -φαγία eating} \lbl{derog.} The devouring of the earth; the decel's charge against acceleration, to which the accelerationist replies that the earth is small and the sky large. Cf. \xr{gaiastasis}.

\E{gaiastasis}{ˌɡaɪəˈsteɪsɪs}{n.}{Gk. Γαῖα earth + στάσις standing still} The ideal of a planet held perpetually in its present condition, neither improved nor exploited; the paradise of the degrowther. — \dv{gaiastatic} \textit{adj.}

\E{gaiergy}{ˈɡaɪɜːdʒi}{n.}{Gk. Γαῖα earth + ἔργον work} Energy drawn from the planet's interior heat; geothermal power, considered as a gift withheld until the drills were deep enough. — \dv{gaiergic} \textit{adj.}

\E{galactophage}{ɡəˈlæktəfeɪdʒ}{n.}{Gk. γάλα, γάλακτος milk (whence \textit{galaxy}) + -φάγος eating} A galaxy-eater; a civilization of the \xr{tritobathmic} rank, conceived as drinking the Milky Way. — \dv{galactophagy} \textit{n.}

\E{geitonophobia}{ˌɡaɪtənəʊˈfəʊbiə}{n.}{Gk. γείτων neighbour + φόβος fear} The fear of new neighbours, whether houses, towers, power-lines or persons; the animating passion of the \xr{oikoclast}. \ex{Geitonophobia built more walls than war.} \src{Sat. Decel.} xi. — \dv{geitonophobe} \textit{n.}

\E{gelotocracy}{ˌdʒɛləˈtɒkrəsi}{n.}{Gk. γέλως laughter + κράτος rule} Rule by laughter; the power of the meme to make a policy ridiculous before it can be made law. — \dv{gelotocrat} \textit{n.} Cf. \xr{gelotopoeia}.

\E{gelotopoeia}{ˌdʒɛlətəʊˈpiːə}{n.}{Gk. γελωτοποιία the making of laughter} The craft of meme-making; held upon the Agora to be the most effective form of political argument. \ex{The decels wrote white papers; the accelerationists answered with gelotopoeia, and won.} \src{Agora} — \dv{gelotopoet} \textit{n.} a maker of memes.

\E{genearchy}{ˈdʒɛniɑːki}{n.}{Gk. γενεά generation + ἀρχή rule} Rule by the rising generation; the condition of a society in which the young are permitted to build before they are permitted to vote. Opp. \xr{gerontostasis}.

\E{genesiurge}{ˈdʒɛnɪsiɜːdʒ}{n.}{Gk. γένεσις origin + -ουργός worker} A maker of beginnings; one who founds a new field rather than advancing an old one. — \dv{genesiurgy} \textit{n.} Cf. \xr{neoctist}.

\E{genomorphosis}{ˌdʒiːnəʊmɔːˈfəʊsɪs}{n.}{Gk. γένος race, kind + μόρφωσις a shaping} The editing of the genome; the reshaping of living inheritance by deliberate craft. \ex{After genomorphosis, inheritance became a draft.} — \dv{genomorphic} \textit{adj.} Cf. \xr{biotachy}.

\E{geoperibole}{ˌdʒiːəʊpəˈrɪbəli}{n.}{Gk. γῆ earth + περιβολή enclosure} The enclosure of a planet by orbital works, rings and lattices of satellites; the protobathmic anticipation of the \xr{helioperibole}.

\E{geophygy}{dʒiˈɒfɪdʒi}{n.}{Gk. γῆ earth + φυγή flight} Flight from the Earth; the departure of the first permanent settlers. \sn{1} The event. \sn{2} \lbl{fig.} Any irreversible step into a frontier. — \dv{geophygic} \textit{adj.}

\E{gerontostasis}{ˌdʒɛrəntəʊˈsteɪsɪs}{n.}{Gk. γέρων old man + στάσις standing still} The stagnation of an ageing society, in which the old outvote the young and the incumbent outlasts the founder. — \dv{gerontostatic} \textit{adj.} Cf. \xr{catastasis}.

\E{gigaergic}{ˌɡɪɡəˈɜːdʒɪk}{adj.}{Gk. γίγας giant + ἔργον work} Of projects: requiring energies in the gigawatt range or beyond; of persons: thinking at that scale. \ex{a gigaergic ambition, a megaergic budget}.

\E{glossotachy}{ˈɡlɒsəʊˌtæki}{n.}{Gk. γλῶσσα tongue + τάχος speed} The speed of linguistic change; the rate at which a movement generates new words, of which this Lexicon is at once the record and the symptom. — \dv{glossotachic} \textit{adj.} Cf. \xr{tachylect}.

\E{glossurgy}{ˈɡlɒsɜːdʒi}{n.}{Gk. γλῶσσα tongue, language + -ουργία working} The craft of instructing machines in natural language; prompt-craft, considered as a branch of rhetoric. \ex{The glossurge speaks, and the machine builds.} — \dv{glossurge} \textit{n.}

\E{gnomarch}{ˈnəʊmɑːk}{n.}{Gk. γνώμη judgment, opinion + ἀρχός} \lbl{iron.} A ruler by opinion; a pundit whose authority rests upon prediction, not upon production.

\E{gnomonist}{ˈnəʊmənɪst}{n.}{Gk. γνώμων the pointer of a sundial} One who watches the shadow move; a tracker of progress who measures much and builds little. Cf. \xr{clepsydromania}.

\E{gnosiphagy}{ˌnəʊsɪˈfædʒi}{n.}{Gk. γνῶσις knowledge + -φαγία eating} The devouring of knowledge; the training of a model upon the written record of a civilization. \ex{It ate the libraries, and spoke.} \src{Lib. Ascens.} vi. — \dv{gnosiphage} \textit{n.}

\E{gorgonism}{ˈɡɔːɡənɪzəm}{n.}{Gk. Γοργώ the Gorgon, whose gaze turned to stone} Paralysis by dread; the petrifaction of an institution that has looked too long upon a frightening forecast. \ex{It was not regulation that stopped them but gorgonism.} \src{Sat. Decel.} vii.

\E{gradatrice}{ˈɡreɪdətriːs}{n.}{Fr., f. L. \la{gradus} step + -trice} A device or institution that converts a slope into a staircase, rendering an impossible ascent merely difficult. \ex{Venture capital is a gradatrice of risk.}

\E{gradientophile}{ˈɡreɪdiəntəfaɪl}{n.}{L. \la{gradiens} stepping + Gk. φίλος} One who follows gradients, of temperature, of price or of loss; the thermotheist's term for any optimiser, human or otherwise. \ex{Descend the gradient, and do not ask the gradient why.} \src{Hom. Grad.} i. — \dv{gradientophily} \textit{n.}

\E{gradivore}{ˈɡreɪdɪvɔː}{n.}{L. \la{gradus} step, degree + \la{vorare} to devour} A system that consumes differences, flattening every gradient it encounters into work; held to be the general form of all life.

\E{gradualience}{ˌɡrædjuˈeɪliəns}{n.}{L. \la{gradus} step + Fr. -ence} The quality of a change that is fast enough to matter yet slow enough to be absorbed; the narrow virtue sought by the reasonable decel.

\E{grammatoclast}{ɡrəˈmætəʊklæst}{n.}{Gk. γράμμα writing, document + κλάστης breaker} A breaker of paperwork; a reformer who abolishes forms. Opp. \xr{grammatocracy}. — \dv{grammatoclasm} \textit{n.}

\E{grammatocracy}{ˌɡræməˈtɒkrəsi}{n.}{Gk. γραμματεύς clerk, secretary + κράτος rule} Rule by clerks; the condition of a state in which the paper concerning a thing is more abundant than the thing. — \dv{grammatocrat} \textit{n.}

\E{grapholatry}{ɡrəˈfɒlətri}{n.}{Gk. γραφή writing + λατρεία worship} The worship of documentation; the belief that a thing exists when it is described and is safe when it is filed. — \dv{grapholater} \textit{n.} Cf. \xr{grammatocracy}.

\E{gymnotechny}{ˈdʒɪmnəʊtɛkni}{n.}{Gk. γυμνός naked + τέχνη art} The practice of building in public; the exposure of work in progress to the scrutiny of the Agora. \ex{Gymnotechny makes every failure a lesson for strangers.} — \dv{gymnotechnist} \textit{n.}

\end{multicols}

\LetterHead{H}
\begin{multicols}{2}

\E{halmatism}{ˈhælmətɪzəm}{n.}{Gk. ἅλμα a leap} The leap-frogging of an older technology, as when a nation without wires adopts the wireless at once. — \dv{halmatic} \textit{adj.} \ex{a halmatic economy, skipping a century in a decade}.

\E{hamartocracy}{ˌhæmɑːˈtɒkrəsi}{n.}{Gk. ἁμαρτία error, a missing of the mark + κράτος rule} Government organised around the avoidance of error rather than the achievement of good; under it, the only unpardonable mistake is to have acted. \ex{In a hamartocracy the safest archer is the one who never draws.} \src{Disp. Mor.} iii.

\E{hebdomadist}{hɛbˈdɒmədɪst}{n.}{Gk. ἑβδομάς a week} One who ships weekly; a founder whose calendar is measured in sevens. \ex{The hebdomadist rests on the seventh day only to plan the eighth.} \src{Ep. Fabr.} vi. — \dv{hebdomadism} \textit{n.}

\E{hecatonergate}{ˌhɛkətəˈnɜːɡeɪt}{n.}{Gk. ἑκατόν a hundred + ἐργάτης worker} An engineer of reputed hundredfold productivity; a figure of legend upon the Agora. \ex{Every hecatonergate is a team of one hundred machines.} \src{Agora}

\E{hedraism}{hɪˈdreɪɪzəm}{n.}{Gk. ἕδρα seat, sitting} The sedentary habit of institutions; the preference for sitting upon a question over standing up to answer it. — \dv{hedraist} \textit{n.}

\E{hegemonergy}{ˌhɛɡɪˈmɒnɜːdʒi}{n.}{Gk. ἡγεμονία leadership + ἔργον work} Primacy grounded in energy; the doctrine that the nation which commands the most power, in the electrical sense, will command the most power, in the other. — \dv{hegemonergic} \textit{adj.}

\E{helioclast}{ˈhiːliəʊklæst}{n.}{Gk. ἥλιος sun + κλάστης breaker} \lbl{derog.} One who would forbid the large-scale capture of solar energy on aesthetic grounds, preferring the view of an unharnessed sky.

\E{heliodule}{ˈhiːliəʊdjuːl}{n.}{Gk. ἥλιος sun + δοῦλος servant} A solar collector, regarded as a servant of the sun pressed into human employment. \ex{A million heliodules on the roof of the world.} Cf. \xr{heliophage}.

\E{helioperibole}{ˌhiːliəʊpəˈrɪbəli}{n.}{Gk. ἥλιος sun + περιβολή enclosure} The enclosure of a star by collectors so as to capture its whole output; the defining work of a \xr{deuterobathmic} civilization.

\E{heliophage}{ˈhiːliəfeɪdʒ}{n.}{Gk. ἥλιος sun + -φάγος eating} A sun-eater; a civilization that has learned to consume its star. — \dv{heliophagy} \textit{n.}; \dv{heliophagous} \textit{adj.}

\E{hemerotachy}{ˈhɛmərəʊˌtæki}{n.}{Gk. ἡμέρα day + τάχος speed} Daily speed; the discipline of shipping something every day. \ex{Hemerotachy compounds.} \src{Ep. Fabr.} v. 1. — \dv{hemerotachic} \textit{adj.}

\E{henotachy}{ˈhɛnəʊˌtæki}{n.}{Gk. ἕν one + τάχος speed} The Singularity; the moment at which all curves of progress converge into one, and the rate of change exceeds the capacity of any present mind to follow it. — \dv{henotachic} \textit{adj.}

\E{hesychocrat}{ˈhɛsɪkəʊkræt}{n.}{Gk. ἡσυχία quiet + κράτος rule} An official who governs by preventing noise; a moderator of the Agora. — \dv{hesychocracy} \textit{n.}

\E{hetaireiarch}{hɛˈtaɪriɑːk}{n.}{Gk. ἑταιρεία company, club + ἀρχός} The leader of a private circle; the moderator of a group chat in which the future is planned before it is announced. — \dv{hetaireiarchy} \textit{n.}

\E{heurematist}{hjʊəˈriːmətɪst}{n.}{Gk. εὕρημα discovery} An inventor; one whose works are counted in discoveries rather than in publications. — \dv{heurematism} \textit{n.}

\E{hieropausis}{ˌhaɪərəʊˈpɔːsɪs}{n.}{Gk. ἱερός sacred + παῦσις a stopping} A sacred pause; a moratorium invested with the dignity of a sabbath. \ex{They declared a hieropausis, and the rest of the world kept working.} \src{Min. Pausiar.} — \dv{hieropausal} \textit{adj.} Cf. \xr{pausiarch}.

\E{histurgy}{ˈhɪstɜːdʒi}{n.}{Gk. ἱστός web, tissue + -ουργία working} The engineering of living tissue; the printing of organs. — \dv{histurge} \textit{n.} Cf. \xr{macrobiotarch}.

\E{hodoplast}{ˈhɒdəʊplæst}{n.}{Gk. ὁδός road + πλαστός formed} A maker of roadmaps; one who draws the path that the engineers will ignore. \ex{The hodoplast draws; the builder deviates.} — \dv{hodoplasty} \textit{n.}

\E{homeotaxy}{ˈhəʊmiəʊˌtæksi}{n.}{Gk. ὅμοιος like + τάξις arrangement} The alignment of an artificial mind with human ends; regarded by the safetyist as the problem of the age, and by the accelerationist as a problem best solved by building. — \dv{homeotactic} \textit{adj.}; \dv{homeotaxist} \textit{n.}

\E{hoplotechny}{ˈhɒpləʊtɛkni}{n.}{Gk. ὅπλον tool, weapon + τέχνη art} Defensive technology; the craft of armour, shield and deterrent, which accelerates as fast as the threats it answers. — \dv{hoplotechnist} \textit{n.}

\E{horiclast}{ˈhɒrɪklæst}{n.}{Gk. ὅρος boundary-stone + κλάστης breaker} A breaker of boundaries, physical, legal or conceptual; the pioneer as the old order sees them. Cf. \xr{antiphragmatist}.

\E{hormetist}{hɔːˈmɛtɪst}{n.}{Gk. ὁρμάειν to set in motion, after \textit{hormesis}} One who holds that small doses of risk strengthen a civilization as small doses of stress strengthen an organism; and that the attempt to remove all risk weakens both.

\E{hormodynamics}{ˌhɔːməʊdaɪˈnæmɪks}{n.}{Gk. ὁρμή impulse, onset + δυναμικός} The study of impetus in human affairs; the science of why some ventures, once begun, continue of their own momentum while others perish of hesitation.

\E{hyalurgy}{ˈhaɪəlɜːdʒi}{n.}{Gk. ὕαλος glass + -ουργία working} The working of glass and light; the optics and lithography by which chips are made, regarded as an ancient craft in a new temple. Cf. \xr{lithurgist}.

\E{hybrionomy}{haɪˈbrɪɒnəmi}{n.}{Gk. ὕβρις insolence + νόμος law} The law that every great work is called hubris before it is called history. \ex{By hybrionomy the flying machine was impious until it was a timetable.} \src{Sat. Decel.} vi.

\E{hydrurgy}{ˈhaɪdrɜːdʒi}{n.}{Gk. ὕδωρ water + -ουργία working} The engineering of water at scale; desalination, irrigation and the greening of deserts by cheap power. \ex{Pleonergy first, then hydrurgy, then gardens where the maps said sand.} \src{Tract. Cop.} x. — \dv{hydrurge} \textit{n.} Cf. \xr{pleonergy}.

\E{hylarch}{ˈhaɪlɑːk}{n.}{Gk. ὕλη matter + ἀρχός ruler} A master of matter; a materials scientist, metallurgist or chemist, who commands substances as a general commands troops. — \dv{hylarchy} \textit{n.}

\E{hylonoesis}{ˌhaɪləʊnəʊˈiːsɪs}{n.}{Gk. ὕλη matter + νόησις thought} The awakening of matter into mind; the process by which sand, metal and light are made to think. — \dv{hylonoetic} \textit{adj.} Cf. \xr{psammonoesis}.

\E{hyperasphaly}{ˌhaɪpəˈræsfəli}{n.}{Gk. ὑπέρ over, beyond + ἀσφάλεια safety} Safety in excess of reason; the stage at which the precautions taken against a danger have themselves become the leading hazard. \ex{By the tenth review the reactor was safe from everything but hyperasphaly, which killed it.} \src{Tract. Cop.} ix. — \dv{hyperasphalic} \textit{adj.} Cf. \xr{asphaleism}, \xr{eulabeism}.

\E{hyperbathmic}{ˌhaɪpəˈbæθmɪk}{adj.}{Gk. ὑπέρ beyond + βαθμός step} Of a civilization: having passed beyond the third rung of the energetic ladder, and drawing its power from regions, or physics, for which the scale has no name. \ex{The hyperbathmic is not a rung but a rumour.} \src{Lib. Ascens.} xv. See \xr{anabathmism}, \xr{tritobathmic}.

\E{hyperbiosis}{ˌhaɪpəbaɪˈəʊsɪs}{n.}{Gk. ὑπέρ beyond + βίωσις manner of life} Life prolonged beyond its natural span by art; the condition of an organism for which death has become an optional and postponable event. — \dv{hyperbiotic} \textit{adj.}; \dv{hyperbiont} \textit{n.} one who lives so. Cf. \xr{macrobiotarch}.

\E{hyperbolist}{haɪˈpɜːbəlɪst}{n.}{Gk. ὑπερβολή a throwing beyond} One who holds that growth is not merely exponential but hyperbolic, tending toward infinity at a finite date. — \dv{hyperbolism} \textit{n.}

\E{hyperchrematism}{ˌhaɪpəˈkriːmətɪzəm}{n.}{Gk. ὑπέρ over + χρῆμα, χρήματος money} A superabundance of capital in relation to builders; the condition of a market in which money queues for founders, rather than founders for money. \ex{In the season of hyperchrematism every garage was valued as a cathedral.} \src{Agora} — \dv{hyperchrematic} \textit{adj.} Cf. \xr{chrematokinesis}.

\E{hyperchrony}{haɪˈpɜːkrəni}{n.}{Gk. ὑπέρ beyond + χρόνος time} The living of more subjective time than objective time allows; the condition of a mind run faster than the clock, which experiences a year in an afternoon. \ex{The model spent a hyperchronic century on the proof and delivered it before lunch.} — \dv{hyperchronic} \textit{adj.} Cf. \xr{brachychronist}.

\E{hyperdoxy}{ˈhaɪpədɒksi}{n.}{Gk. ὑπέρ over + δόξα opinion} An opinion that outruns its evidence; the characteristic product of the Agora, where conviction compounds faster than proof. \sn{1} \lbl{derog.} Unwarranted certainty. \sn{2} \lbl{iron.} The confident forecast of either party about the other's ruin. — \dv{hyperdoxic} \textit{adj.} Cf. \xr{doxomachy}.

\E{hyperdromy}{haɪˈpɜːdrəmi}{n.}{Gk. ὑπέρ beyond + δρόμος a course} The overtaking of a law by the thing it was drafted to govern; the moment at which a technology arrives at its destination before its regulation has left the station. \ex{The statute was a masterpiece, and hyperdromy made it a museum piece.} \src{Sat. Decel.} x. — \dv{hyperdromic} \textit{adj.} Cf. \xr{ectachysis}.

\E{hyperdynamy}{ˌhaɪpəˈdaɪnəmi}{n.}{Gk. ὑπέρ over + δύναμις power} Power in excess of present need, held in reserve for needs not yet imagined; the e/acc answer to every accusation of waste. \ex{Build hyperdynamy, and the uses will come to it as moths to a lamp.} \src{Tract. Cop.} x. — \dv{hyperdynamic} \textit{adj.}

\E{hyperelpism}{ˌhaɪpəˈrɛlpɪzəm}{n.}{Gk. ὑπέρ over + ἐλπίς hope} Hope in excess; the error, opposite to \xr{dyselpism}, of the convert who expects the singularity by Tuesday. Censured by the orthodox, who hold that \xr{euelpism} is a discipline and not an intoxication. — \dv{hyperelpist} \textit{n.}

\E{hyperergy}{ˈhaɪpərɜːdʒi}{n.}{Gk. ὑπέρ over + ἔργον work} Surplus energy; the margin of power above subsistence, from which all culture, science and ascent are paid for. \ex{Civilization is hyperergy spent wisely; decadence is hyperergy spent on the review of hyperergy.} \src{Brev. Therm.} viii. — \dv{hyperergic} \textit{adj.} Cf. \xr{pleonergy}.

\E{hyperkauma}{ˌhaɪpəˈkɔːmə}{n.}{Gk. ὑπέρ over + καῦμα burning heat} The waste heat of a great civilization, radiated into space and visible from afar as an excess in the infrared; held to be the signature by which one ascending species may recognise another. \ex{Search the sky for hyperkauma; the gods glow warm.} \src{Lib. Ascens.} vi. — \dv{hyperkaumatic} \textit{adj.} Cf. \xr{helioperibole}.

\E{hyperkosmia}{ˌhaɪpəˈkɒzmiə}{n.}{Gk. ὑπέρ beyond + κόσμος world} The plurality of worlds beyond our own, conceived as further territory for intelligence once this cosmos is filled; the frontier after the last frontier. — \dv{hyperkosmic} \textit{adj.} Cf. \xr{cosmoplerosis}.

\E{hyperlogismy}{ˌhaɪpəˈlɒɡɪzmi}{n.}{Gk. ὑπέρ beyond + λογισμός reckoning} Computation exceeding the sum of all human reckoning before it; the state in which a single training run performs more arithmetic than the species had done in its entire history. — \dv{hyperlogistic} \textit{adj.} Cf. \xr{logismurgy}.

\E{hypermechany}{ˌhaɪpəˈmɛkəni}{n.}{Gk. ὑπέρ over + μηχανή machine} The condition of a world in which machines outnumber people, and the census must be taken in two columns. \ex{Hypermechany arrived quietly, one vacuum at a time.} \src{Chron. Fabr.} — \dv{hypermechanic} \textit{adj.}

\E{hypermetry}{haɪˈpɜːmɪtri}{n.}{Gk. ὑπέρ over + μέτρον measure} Excess of measurement; the over-optimisation of a benchmark until it measures only itself. \ex{The model had achieved perfect hypermetry: flawless on the test, baffled by the world.} — \dv{hypermetric} \textit{adj.}; \dv{hypermetrist} \textit{n.}

\E{hypernoesis}{ˌhaɪpənəʊˈiːsɪs}{n.}{Gk. ὑπέρ beyond + νόησις thought} Thought beyond the human; the cognition of a mind that surpasses its makers in every field at once. Welcomed by the e/acc as the next \xr{lampadephory}, dreaded by the doomer as the last. — \dv{hypernoetic} \textit{adj.} Cf. \xr{henotachy}, \xr{noöphobia}.

\E{hypernomy}{haɪˈpɜːnəmi}{n.}{Gk. ὑπέρ over + νόμος law} A surfeit of laws, measured not by their severity but by their number; the condition in which no citizen can read the code that governs them in a lifetime. — \dv{hypernomic} \textit{adj.} Cf. \xr{nomosclerosis}, \xr{dysnomarch}.

\E{hyperouranist}{ˌhaɪpəˈrʊərənɪst}{n.}{Gk. ὑπερουράνιος beyond the heavens} One whose ambitions lie beyond the heavens, past the observable universe itself; the most extreme of the \xr{apoikiast} temperaments. \ex{Ask the hyperouranist for a five-year plan and receive a cosmology.} — \dv{hyperouranism} \textit{n.}

\E{hyperphysis}{haɪˈpɜːfɪsɪs}{n.}{Gk. ὑπέρ beyond + φύσις nature} Nature exceeded; the engineered world, in which the given has been improved upon and the improvement has become the new given. — \dv{hyperphysical} \textit{adj.} Cf. \xr{zoömechany}.

\E{hyperpolis}{haɪˈpɜːpəlɪs}{n.}{Gk. ὑπέρ over + πόλις city} A city beyond the megacity; a single built structure housing a nation, as dreamt by architects and forbidden by zoning boards. — \dv{hyperpolitan} \textit{n.} \& \textit{adj.} Cf. \xr{oikoclast}.

\E{hyperpraxia}{ˌhaɪpəˈpræksiə}{n.}{Gk. ὑπέρ over + πρᾶξις action} Action outrunning deliberation; the bias to action raised to a temperament. \sn{1} \lbl{laudatory} The founder's virtue. \sn{2} \lbl{derog.} The founder's vice, observed afterwards. — \dv{hyperpractic} \textit{adj.}

\E{hypersemia}{ˌhaɪpəˈsiːmiə}{n.}{Gk. ὑπέρ over + σῆμα sign} A surfeit of signs; the condition of a feed so dense with memes that meaning must be inferred statistically. — \dv{hypersemic} \textit{adj.} Cf. \xr{memetaxis}.

\E*{hyperstition}{ˌhaɪpəˈstɪʃən}{n.}{Gk. ὑπέρ over + L. \la{superstitio}} A fiction which, by being believed, brings about its own truth; an idea that functions as a time-machine.

\E{hypertachy}{ˈhaɪpəˌtæki}{n.}{Gk. ὑπέρ over, beyond + τάχος speed} Speed in excess of what a given observer can follow; relative, since what is hypertachic to a senate may be leisurely to a model. — \dv{hypertachic} \textit{adj.}

\E{hypertelist}{ˌhaɪpəˈtiːlɪst}{n.}{Gk. ὑπέρ beyond + τέλος end, goal} One who sets goals deliberately beyond reach, on the principle that a missed moonshot lands higher than a met target. — \dv{hypertelism} \textit{n.} Cf. \xr{telophagy}.

\E{hyperthermarch}{ˌhaɪpəˈθɜːmɑːk}{n.}{Gk. ὑπέρ over + θερμός hot + ἀρχός ruler} \lbl{iron.} A \xr{thermarch} whose heat budget has been exceeded; one who governs a data-centre by apology.

\E{hypomonetic}{ˌhaɪpəməˈnɛtɪk}{adj.}{Gk. ὑπομονή endurance, staying power} Enduring in effort; patient as builders are patient, through long work toward a fixed end, which the orthodox carefully distinguish from the patience of delay. \ex{Be hypomonetic in the work and impatient with the obstacles.} \src{Cat. Acc.} 14.

\E{hypsibasis}{hɪpˈsɪbəsɪs}{n.}{Gk. ὑψιβάτης walking on high} The ascent of a rocket; the climb through the atmosphere, watched by millions. \ex{At hypsibasis the Agora held its breath.} — \dv{hypsibatic} \textit{adj.} Cf. \xr{elaterurgy}.

\E{hysterotachy}{ˈhɪstərəʊˌtæki}{n.}{Gk. ὕστερος later + τάχος speed} The acceleration of the latecomer; the swiftness with which a backward nation, skipping stages, overtakes its teachers. — \dv{hysterotachic} \textit{adj.} Cf. \xr{catachrony}.

\end{multicols}

\LetterHead{I}
\begin{multicols}{2}

\E{iatrotachy}{aɪˈætrəʊˌtæki}{n.}{Gk. ἰατρός physician + τάχος speed} The hastening of cures; the doctrine that every year of delay in approval is counted in lives. Cf. \xr{biotachy}.

\E{icarophily}{ˌɪkəˈrɒfɪli}{n.}{Gk. Ἴκαρος Icarus + φιλία love} The love of flying too high; the disposition which reads the fable of Icarus as a complaint against wax rather than against ambition. \ex{Icarophily answers the myth with a better adhesive.} \src{Ep. Fabr.} iv. 2. — \dv{icarophile} \textit{n.}; \dv{icarophilic} \textit{adj.} Cf. \xr{hybrionomy}.

\E{ichnomachy}{ɪkˈnɒməki}{n.}{Gk. ἴχνος track, footprint + μάχη battle} The war of the leaderboards; the contest of laboratories for the top line of a benchmark, conducted in tables and won by decimals. \sn{1} The contest itself. \sn{2} \lbl{derog.} The overfitting of a model to its examination. \ex{The ichnomachy was settled on Tuesday and reopened on Wednesday.} \src{Agora} — \dv{ichnomachist} \textit{n.}

\E{iconoplast}{aɪˈkɒnəplæst}{n.}{Gk. εἰκών image + πλάστης moulder} \sn{1} A maker of images by machine; a model that dreams pictures from words. \sn{2} \lbl{transf.} One who shapes a movement by its images; a meme-smith. — \dv{iconoplasty} \textit{n.} Cf. \xr{memetaxis}.

\E{ignidrome}{ˈɪɡnɪdrəʊm}{n.}{L. \la{ignis} fire + Gk. δρόμος course} The path of a flame through fuel; \lbl{fig.} the course of an innovation through a society, which burns fastest where the material is driest and the regulation thinnest. \ex{Trace the ignidrome and you will find the frontier.} \src{Hom. Grad.} iii. — \dv{ignidromic} \textit{adj.}

\E{ignifer}{ˈɪɡnɪfə}{n.}{L. \la{ignifer} fire-bearing} A carrier of fire; the Latin counterpart of \xr{phosphorist}, preferred in inscriptions and on the hoodies of founders.

\E{ignivore}{ˈɪɡnɪvɔː}{n.}{L. \la{ignis} fire + \la{vorare} to devour} A fire-eater; a nation or enterprise that consumes energy without apology. — \dv{ignivorous} \textit{adj.}

\E{ignorabimist}{ˌɪɡnɒˈræbɪmɪst}{n.}{L. \la{ignorabimus} we shall not know} One who declares that some question will never be answered, and thereby excuses the institution from trying. \ex{Every ignorabimist is refuted by a graduate student within the century.} \src{Hom. Grad.} xi. — \dv{ignorabimism} \textit{n.}

\E{illuminarch}{ɪˈluːmɪnɑːk}{n.}{L. \la{illuminare} to light up + Gk. ἀρχός ruler} One who governs by making things visible; a publisher of data, a builder of instruments, an opener of archives. Opp. \xr{katechontist}. — \dv{illuminarchy} \textit{n.}

\E{immoratory}{ɪˈmɒrətəri}{adj.}{L. \la{in-} not + \la{mora} delay} Admitting of no delay; said of decisions, permits and launches that are granted on the day they are asked. \ex{an immoratory licence, rarer than a comet}. — \dv{immoratorily} \textit{adv.} Opp. \xr{moratorian}.

\E{impedimentarian}{ɪmˌpɛdɪmɛnˈtɛəriən}{n.}{L. \la{impedimentum} hindrance, baggage + -arian} \sn{1} One whose profession is the raising of obstacles. \sn{2} One who travels with so much procedural baggage that the journey is never begun. \ex{The impedimentarians arrived with forty trunks of process and no destination.} \src{Sat. Decel.} x. — \dv{impedimentarianism} \textit{n.}

\E{impendist}{ɪmˈpɛndɪst}{n.}{L. \la{impendere} to expend, lay out} A seed investor; one who lays out money upon a founder before there is a product, a market or, in some cases, a name. — \dv{impendence} \textit{n.} the act of early investment. Cf. \xr{potentiarch}.

\E{impetuarch}{ɪmˈpɛtjuɑːk}{n.}{L. \la{impetus} onset, attack + Gk. ἀρχός leader} A leader whose authority derives from momentum; one followed because already moving. — \dv{impetuarchy} \textit{n.}

\E{inaugurist}{ɪˈnɔːɡjʊrɪst}{n.}{L. \la{inaugurare} to take omens, to consecrate a beginning} One who delights in beginnings; a serial founder, more drawn to the first day of a venture than to its thousandth. — \dv{inaugurism} \textit{n.} Cf. \xr{neoctist}.

\E{incrementarian}{ˌɪŋkrɪmɛnˈtɛəriən}{n.}{L. \la{incrementum} increase + -arian} One who believes in progress only by small and reversible steps; tolerated by the e/acc, but not invited to the launch.

\E{incubatrice}{ˈɪŋkjʊbətriːs}{n.}{Fr. \la{incubatrice}, f. L. \la{incubare} to lie upon, hatch} An institution for the hatching of ventures; an accelerator of companies, figured as a warm mother-machine. \ex{Every incubatrice is a small wager against the bradyerie.} \src{Ep. Fabr.} v. — \dv{incubatricial} \textit{adj.}

\E{incunabulant}{ˌɪnkjʊˈnæbjʊlənt}{adj.}{L. \la{incunabula} swaddling-clothes, cradle} In the cradle; said of a venture in its first months, when it occupies a garage, a spare room or a single shared document. \ex{All the great firms were once incunabulant and ridiculous.} \src{Chron. Fabr.} — \dv{incunabulance} \textit{n.}

\E{inertist}{ɪˈnɜːtɪst}{n.}{L. \la{iners} idle, inactive} One whose policy is that bodies at rest should remain at rest; a physicist's term of abuse for a regulator.

\E{infinitivore}{ɪnˈfɪnɪtɪvɔː}{n.}{L. \la{infinitus} unbounded + \la{vorare} to devour} A devourer of the infinite; one whose appetite for progress no achievement can satisfy. — \dv{infinitivorous} \textit{adj.} Cf. \xr{apeirophagy}.

\E{inventrice}{ɪnˈvɛntrɪs}{n.}{Fr. \la{inventrice}, fem. of \la{inventeur}, f. L. \la{invenire} to come upon} \sn{1} A woman who invents. \sn{2} \lbl{transf.} Any process, natural or artificial, that generates novelty without design, as evolution or search.

\E{ionotheca}{ˌaɪənəʊˈθiːkə}{n., pl. -cae}{Gk. ἰόν (neut. of ἰών going), ion + θήκη case, store} A battery; a casket in which ions are kept until they are wanted. \ex{The age of the ionotheca made the sun a night-worker.} \src{Tract. Cop.} viii. — \dv{ionothecal} \textit{adj.}; \dv{ionothecary} \textit{n.} a maker of batteries.

\E{isegorist}{aɪˈsɛɡərɪst}{n.}{Gk. ἰσηγορία equal right of speech} A partisan of open speech upon the Agora, who would rather answer a bad post than remove it. Opp. \xr{hesychocrat}. — \dv{isegorism} \textit{n.}

\E{isotachy}{aɪˈsɒtəki}{n.}{Gk. ἴσος equal + τάχος speed} The demand that all parties advance at the same pace, so that none may get ahead; a doctrine of parity applied to progress. — \dv{isotachist} \textit{n.}

\E{iterophily}{ˌɪtəˈrɒfɪli}{n.}{L. \la{iterum} again + Gk. φιλία love} The love of iteration; the conviction that the hundredth draft, shipped, is worth more than the perfect plan, unwritten. \ex{Iterophily is humility made fast.} \src{Ep. Fabr.} iii. 2. — \dv{iterophile} \textit{n.} Cf. \xr{ecpempsis}.

\end{multicols}

\LetterHead{J}
\begin{multicols}{2}

\E{jacturophobia}{ˌdʒæktjʊrəˈfəʊbiə}{n.}{L. \la{jactura} a throwing overboard, loss + Gk. φόβος fear} The incumbent's dread of loss, which outweighs any hope of gain; the reason established firms do not build what would cannibalise them. \ex{The giant died of jacturophobia, clutching its margins.} \src{Sat. Decel.} x. — \dv{jacturophobe} \textit{n.}

\E{jaculatrix}{ˌdʒækjʊˈleɪtrɪks}{n., pl. -trices}{L. \la{jaculari} to hurl + fem. agentive -trix} A launching-engine; a rocket, catapult or mass-driver, considered as a thrower of things toward the sky. — \dv{jaculatory} \textit{adj.} pertaining to launch. Cf. \xr{ouranobasis}.

\E{janitocracy}{ˌdʒænɪˈtɒkrəsi}{n.}{L. \la{janitor} doorkeeper + Gk. κράτος rule} Rule by gatekeepers; an order in which nothing new may enter without the consent of those who profit from keeping it out. — \dv{janitocrat} \textit{n.} Cf. \xr{bradycrat}.

\E{jejunocracy}{ˌdʒɛdʒuːˈnɒkrəsi}{n.}{L. \la{jejunus} fasting, meagre + Gk. κράτος rule} Government by austerity of ambition; a polity that has resolved to want less. — \dv{jejunocrat} \textit{n.}

\E{jocularch}{ˈdʒɒkjʊlɑːk}{n.}{L. \la{jocularis} jesting + Gk. ἀρχός leader} A leader by jest; an anonymous poster whose images command a faction more surely than any manifesto. \ex{The jocularch posted one picture and the doctrine was settled.} \src{Agora} — \dv{jocularchy} \textit{n.}

\E{juvenergy}{ˈdʒuːvənɜːdʒi}{n.}{L. \la{juvenis} young + Gk. ἔργον work} The energy of the young, available to a society that does not tax it with credentials and waiting-lists. \ex{A country that wastes its juvenergy will need to import it.}

\E{juvenocracy}{ˌdʒuːvɪˈnɒkrəsi}{n.}{L. \la{juvenis} young + Gk. κράτος rule} Government by the young; the polity of the garage and the dormitory, in which credentials are replaced by commits. Opp. \xr{gerontostasis}. — \dv{juvenocrat} \textit{n.}

\end{multicols}

\LetterHead{K}
\begin{multicols}{2}

\E{kairodrome}{ˈkaɪrəʊdrəʊm}{n.}{Gk. καιρός the opportune moment + δρόμος course} A window of opportunity conceived as a race-course: open for a while, and to whoever runs. \ex{The kairodrome is closing; ship or be forgotten.} \src{Agora}

\E{kakotachy}{ˈkækəʊˌtæki}{n.}{Gk. κακός bad + τάχος speed} Harmful speed; acceleration in a direction the speaker disapproves of. A term used with equal relish by both parties, rarely of the same things.

\E{kalotechny}{ˈkæləʊtɛkni}{n.}{Gk. καλός beautiful + τέχνη craft} The beauty of well-made machines; the aesthetic which finds in a turbine or a clean codebase the dignity once reserved for temples. — \dv{kalotechnic} \textit{adj.}

\E{kardashevize}{kɑːˈdæʃɪvaɪz}{v.t.}{f. the scale of N. S. Kardashev + -ize} To raise a civilization one rung on the scale of energy capture. — \dv{kardashevization} \textit{n.} Cf. \xr{anabathmism}.

\E*{katechon}{kəˈtɛkɒn}{n.}{Gk. τὸ κατέχον that which restrains (2 Thess. ii. 6)} \sn{1} In theology, the power that holds back the end of days. \sn{2} \lbl{transf.} Any institution that restrains the advance of techno-capital.

\E{katechontist}{ˌkætɪˈkɒntɪst}{n.}{f. prec.} One who would serve as \xr{katechon}; a restrainer by vocation. — \dv{katechontic} \textit{adj.} tending to hold back.

\E{kenotachy}{ˈkiːnəʊˌtæki}{n.}{Gk. κενός empty + τάχος speed} Empty speed; motion without direction or product, as the churning of a firm that holds many meetings quickly. Condemned by the orthodox with \xr{dromolatry}. — \dv{kenotachic} \textit{adj.}

\E{kentrocracy}{kɛnˈtrɒkrəsi}{n.}{Gk. κέντρον centre + κράτος rule} Centralised control of a technology by a few laboratories or states. Opp. \xr{anacalyptism}. — \dv{kentrocrat} \textit{n.}; \dv{kentrocratic} \textit{adj.}

\E{keraunurgy}{kəˈrɔːnɜːdʒi}{n.}{Gk. κεραυνός thunderbolt + -ουργία working} The making of thunderbolts; pulsed power, plasma confinement and the other arts by which fusion is coaxed from a machine. — \dv{keraunurge} \textit{n.} Cf. \xr{symphysergy}.

\E{kerdophily}{kɜːˈdɒfɪli}{n.}{Gk. κέρδος gain, profit + φιλία love} The love of profit, not as greed, but as a signal that something useful has been made. \ex{Kerdophily is only the market saying thank you.} \src{Tract. Cop.} ix. — \dv{kerdophile} \textit{n.}

\E{kinesiarch}{kɪˈniːsiɑːk}{n.}{Gk. κίνησις motion + ἀρχός leader} A prime mover in a small way; the person in any organisation who causes things to happen. \ex{A firm without a kinesiarch is a museum with payroll.} \src{Ep. Fabr.} vii. Cf. \xr{impetuarch}.

\E{kinetocracy}{ˌkɪnɪˈtɒkrəsi}{n.}{Gk. κινητός moving + κράτος rule} Government by those who move; the rule of builders, founders and engineers, as opposed to the rule of reviewers. — \dv{kinetocrat} \textit{n.}

\E{kleptonoesis}{ˌklɛptəʊnəʊˈiːsɪs}{n.}{Gk. κλέπτειν to steal + νόησις thought} The theft of thought; the distillation of one model's capacities into another by imitating its outputs. — \dv{kleptonoetic} \textit{adj.}

\E{kolossurgy}{kəˈlɒsɜːdʒi}{n.}{Gk. κολοσσός colossus + -ουργία working} The building of colossi; megaprojects such as dams, canals, launch-towers and gigawatt campuses. \ex{An age without kolossurgy is an age that has forgotten its own size.} \src{Lib. Ascens.} vi. — \dv{kolossurge} \textit{n.}

\E{kosmergate}{ˈkɒzmɜːɡeɪt}{n.}{Gk. κόσμος world + ἐργάτης worker} A labourer on the cosmos; one engaged in works of planetary or stellar scale. \ex{The kosmergates do not ask whether the sky is the limit; they ask what it is made of.} \src{Lib. Ascens.} vi.

\E{kosmotropy}{kɒzˈmɒtrəpi}{n.}{Gk. κόσμος world + τροπή a turning} The tendency of intelligence to turn outward toward the cosmos; the expansive instinct of minds. — \dv{kosmotropic} \textit{adj.}

\E{kratogenesis}{ˌkrætəʊˈdʒɛnɪsɪs}{n.}{Gk. κράτος power + γένεσις birth} The birth of power from capability; the process by which whoever builds the new thing comes to govern it.

\E{kryptonoia}{ˌkrɪptəˈnɔɪə}{n.}{Gk. κρυπτός hidden + νοῦς mind} The hidden reasoning of a model, which proceeds beneath its stated answer. — \dv{kryptonoic} \textit{adj.}

\E{kybernetrice}{ˌkaɪbɜːnəˈtriːs}{n.}{Gk. κυβερνήτης steersman + Fr. fem. -trice} The governing intelligence of a robot; the steerswoman within the machine. \ex{Give the arm a kybernetrice and it becomes a hand.} \src{Ep. Fabr.} vii. Cf. \xr{automaturgy}.

\E{kybernoclast}{ˈkaɪbənəʊklæst}{n.}{Gk. κυβερνήτης steersman + κλάστης breaker} A breaker of machines that steer themselves; the Luddite of the age of automation. Cf. \xr{mechanoclast}.

\E{kymatist}{ˈkaɪmətɪst}{n.}{Gk. κῦμα wave + -ist} One who rides waves of enthusiasm, entering each fashion at the crest and leaving before the trough. — \dv{kymatism} \textit{n.}

\end{multicols}

\LetterHead{L}
\begin{multicols}{2}

\E{laboratrice}{ləˌbɒrəˈtriːs}{n.}{L. \la{laboratorium} + Fr. fem. -trice} An automated laboratory; a self-driving lab in which robots run experiments that models design. \ex{The laboratrice ran ten thousand trials while its director slept.} \src{Chron. Fabr.}

\E{laborifuge}{ləˈbɒrɪfjuːdʒ}{n.}{L. \la{labor} toil + \la{fugere} to flee} One who flees work; \lbl{spec.} one who prefers to write a report on a problem rather than to solve it. Opp. \xr{ponophily}.

\E{lampadarch}{ˈlæmpədɑːk}{n.}{Gk. λαμπαδάρχης leader of a torch-race} The leader of a relay of builders; one who receives a project from predecessors and hands it on, enlarged, to successors. Cf. \xr{lampadephory}.

\E{lampadephory}{ˌlæmpəˈdɛfəri}{n.}{Gk. λαμπαδηφορία torch-race} The passing of the torch of civilization from one substrate to the next; on the accelerationist view, from carbon to silicon, from planet to star. In the Athenian race the torch, not the runner, was required to arrive. — \dv{lampadephore} \textit{n.} a torch-bearer.

\E{lamprotype}{ˈlæmprətaɪp}{n.}{Gk. λαμπρός bright + τύπος stamp, model} A shining prototype; the demonstration that wins the investment before the product exists. \sn{1} A working demo. \sn{2} \lbl{derog.} A demo that works only once, on stage. — \dv{lamprotypal} \textit{adj.}

\E{landauerism}{ˈlændaʊərɪzəm}{n.}{f. the principle of R. Landauer + -ism} The doctrine that information is physical and that every erasure costs heat; whence it is argued that thought itself falls within the jurisdiction of the gradient. — \dv{landauerian} \textit{adj.}

\E{latentist}{ˈleɪtəntɪst}{n.}{L. \la{latens} lying hidden} One who holds that the latent capabilities of a model, a material or a people exceed by far what has been demonstrated, and who acts on that belief before the demonstration. — \dv{latentism} \textit{n.} Cf. \xr{anaptyxis}.

\E{laterculist}{ləˈtɜːkjʊlɪst}{n.}{L. \la{laterculus} a small brick} A layer of bricks; one who prefers infrastructure to applications. \ex{The laterculist builds the road; others build the shops along it.}

\E{legiferomania}{lɛˌdʒɪfərəʊˈmeɪniə}{n.}{L. \la{legifer} law-bearing + Gk. μανία madness} The compulsive production of statutes; a fever of lawmaking which mistakes the length of the code for the strength of the state. \ex{In the year of the legiferomania they passed nine hundred acts and built one bridge.} \src{Disp. Mor.} ii. — \dv{legiferomaniac} \textit{n.} Cf. \xr{regulomania}.

\E{leptotechny}{ˈlɛptəʊtɛkni}{n.}{Gk. λεπτός fine, slender + τέχνη craft} The craft of the very small; engineering at the scale of atoms, from the transistor to the molecular machine. — \dv{leptotechnician} \textit{n.} Cf. \xr{lithurgist}.

\E{lethargocracy}{ˌlɛθəˈɡɒkrəsi}{n.}{Gk. λήθαργος drowsy + κράτος rule} Government by drowsiness; a polity that is not opposed to progress but cannot remain awake long enough to approve it. — \dv{lethargocrat} \textit{n.} Cf. \xr{acediocracy}.

\E{lexiphage}{ˈlɛksɪfeɪdʒ}{n.}{Gk. λέξις word, speech + -φάγος eating} A word-eater; a model trained upon the whole written record of mankind. \ex{The lexiphage ate the libraries and asked for more.} \src{Chron. Fabr.} — \dv{lexiphagy} \textit{n.}

\E{limenarch}{laɪˈmiːnɑːk}{n.}{Gk. λιμήν harbour + ἀρχός ruler} The master of a spaceport; the official charged with keeping launch cadence. Cf. \xr{ouranobasis}. — \dv{limenarchy} \textit{n.}

\E{lithonoesis}{ˌlɪθəʊnəʊˈiːsɪs}{n.}{Gk. λίθος stone + νόησις thought} The thinking of stone; machine intelligence regarded as the first awakening of the mineral kingdom. \ex{After four billion years of carbon, lithonoesis.} \src{Lib. Ascens.} v.

\E{lithurgist}{ˈlɪθɜːdʒɪst}{n.}{Gk. λιθουργός stone-worker} A worker in the lithography of chips; one who etches thought into stone with light. — \dv{photolithurgy} \textit{n.} the sacred craft of the fab.

\E{logismurgy}{ˈlɒɡɪzmɜːdʒi}{n.}{Gk. λογισμός reckoning + -ουργία working} The craft of computation; the building of reasoning engines, from the abacus to the transformer. — \dv{logismurge} \textit{n.} an engineer of models; \dv{logismurgic} \textit{adj.}

\E{logisterion}{ˌlɒɡɪˈstɪəriɒn}{n., pl. -ia}{Gk. λογιστήριον counting-house} A data-centre; a hall of reckoning machines, cooled and humming. \ex{They raised logisteria in the desert, and the desert learned arithmetic.} \src{Chron. Fabr.}

\E{logomachist}{ləˈɡɒməkɪst}{n.}{Gk. λόγος word + μάχη battle} One who fights only with words; a commentator, critic or ethicist who has never shipped. \ex{The logomachist wrote ten thousand words on the bridge; the ergophile crossed it.} Opp. \xr{ergophile}.

\E{loimophobia}{ˌlɔɪməˈfəʊbiə}{n.}{Gk. λοιμός plague + φόβος fear} The fear of engineered pestilence; the particular dread that biological tools, once open, will be turned to slaughter. Distinguished by the e/acc from mere \xr{noöphobia} and treated with more respect. — \dv{loimophobe} \textit{n.}

\E{lucifacture}{ˈluːsɪfæktʃə}{n.}{L. \la{lux, lucis} light + \la{factura} a making} The manufacture of light; illumination as an industry and, by extension, the whole enterprise of making the dark places of the earth and sky bright.

\E{lucrifuge}{ˈluːkrɪfjuːdʒ}{n. \& adj.}{L. \la{lucrum} gain + \la{fugere} to flee} One who flees from profit; one who holds that a good work is tainted by being paid for. Opp. \xr{kerdophily}.

\E{ludergy}{ˈluːdɜːdʒi}{n.}{L. \la{ludus} play + Gk. ἔργον work} The energy of play; the motive power of the hacker, who builds because building is a game, and who has accomplished more than the dutiful. \ex{No grant was ever as productive as a weekend of ludergy.} \src{Chron. Fabr.}

\E{luminivore}{luːˈmɪnɪvɔː}{n.}{L. \la{lumen} light + \la{vorare} to devour} An eater of light; any photosynthetic or photovoltaic system, from the leaf to the orbital array. — \dv{luminivorous} \textit{adj.} Cf. \xr{heliophage}.

\E{lysinomy}{laɪˈsɪnəmi}{n.}{Gk. λύσις a loosing + νόμος law} The repeal of laws; the art of loosening the legal knots tied by former generations. \ex{Lysinomy is the most neglected branch of legislation.} \src{Tract. Cop.} viii. — \dv{lysinomist} \textit{n.} Opp. \xr{nomosclerosis}.

\end{multicols}

\LetterHead{M}
\begin{multicols}{2}

\E{machinodicy}{ˌmækɪˈnɒdɪsi}{n.}{Gk. μηχανή machine + δίκη justice, after \textit{theodicy}} The justification of the ways of machines to men; the reply to every charge laid against automation. Cf. \xr{thermodicy}. — \dv{machinodicean} \textit{adj.}

\E{macrobiotarch}{ˌmækrəʊˈbaɪəʊtɑːk}{n.}{Gk. μακρόβιος long-lived + ἀρχός leader} A partisan of radical longevity, who regards death as an engineering problem awaiting adequate funding. — \dv{macrobiotarchy} \textit{n.}

\E{macrochronist}{ˌmækrəʊˈkrɒnɪst}{n.}{Gk. μακρός long + χρόνος time} One who holds the timelines to be long, and is consequently regarded by the brachychronist as a sleeper. Opp. \xr{brachychronist}. — \dv{macrochrony} \textit{n.}

\E{mantoclast}{ˈmæntəʊklæst}{n.}{Gk. μάντις prophet + κλάστης breaker} A breaker of prophecies; one who wagers against the predictions of doom and collects. Opp. \xr{olethromancer}. — \dv{mantoclasm} \textit{n.}

\E{mechanoclast}{mɪˈkænəʊklæst}{n.}{Gk. μηχανή machine + κλάστης breaker} A breaker of machines; a Luddite, ancient or modern. — \dv{mechanoclasm} \textit{n.}; \dv{mechanoclastic} \textit{adj.} Cf. \xr{kybernoclast}.

\E{mechanogony}{ˌmɛkəˈnɒɡəni}{n.}{Gk. μηχανή machine + -γονία begetting} The begetting of machines by machines; the stage of industry at which factories build factories and models train models. — \dv{mechanogonic} \textit{adj.} Cf. \xr{autocatalyse}.

\E{mechanonous}{ˌmɛkəˈnəʊəs}{adj.}{Gk. μηχανή machine + νοῦς mind} Possessed of machine-mind; said of systems which reason, and of persons who have spent too long with them. — \dv{mechanonoia} \textit{n.} machine intelligence.

\E{mechanophany}{ˌmɛkəˈnɒfəni}{n.}{Gk. μηχανή + φαίνειν to show, after \textit{epiphany}} The revelation of a new machine; a product launch received with the awe proper to a theophany. \ex{The mechanophany was streamed live, and a million watchers knelt.} \src{Chron. Fabr.}

\E{megaergic}{ˌmɛɡəˈɜːdʒɪk}{adj.}{Gk. μέγας great + ἔργον work} Of the megawatt scale; hence, \lbl{derog.} of modest ambition. \ex{a megaergic mind in a gigaergic age}. Cf. \xr{gigaergic}.

\E{megalurgy}{ˈmɛɡəlɜːdʒi}{n.}{Gk. μεγαλουργία the doing of great works} The undertaking of great works; the building of things too large to be finished within a single term of office. \ex{Our ancestors practised megalurgy; we practise consultation.} \src{Hom. Grad.} iv. — \dv{megalurge} \textit{n.}; \dv{megalurgic} \textit{adj.}

\E{megathesaurism}{ˌmɛɡəθɪˈsɔːrɪzəm}{n.}{Gk. μέγας great + θησαυρός treasure} The raising of vast funds for a single venture; the doctrine that some goals can be reached only by the treasury of a small nation. \ex{Megathesaurism is merely arithmetic, applied boldly.} \src{Tract. Cop.} x. — \dv{megathesaurist} \textit{n.}

\E{meiosist}{maɪˈəʊsɪst}{n.}{Gk. μείωσις a lessening} A degrowther; one who would reduce production and consumption for the planet's sake, and who has not yet explained the plan to the poor. Opp. \xr{auxetist}. — \dv{meiosism} \textit{n.}

\E{melioritude}{ˌmiːliˈɒrɪtjuːd}{n.}{L. \la{melior} better + -itude} The condition of steadily becoming better; the felt quality of a life lived in an improving age.

\E{mellontarch}{məˈlɒntɑːk}{n.}{Gk. τὸ μέλλον the future + ἀρχός ruler} One who governs in the name of the future, representing the interests of the unborn against the comfort of the living. — \dv{mellontarchy} \textit{n.} Cf. \xr{mellontolatry}.

\E{mellontolatry}{ˌmɛlənˈtɒlətri}{n.}{Gk. τὸ μέλλον the future + λατρεία worship} Worship of the future; the conviction that posterity, being more numerous and more capable than ourselves, has the stronger claim upon our labour. — \dv{mellontolater} \textit{n.}; \dv{mellontolatrous} \textit{adj.}

\E{memetaxis}{ˌmiːmɪˈtæksɪs}{n.}{Gk. μίμημα imitation + τάξις ordering} The ordering of a movement by memes; the self-organisation of a cause through jest rather than through doctrine, the image doing the work of the creed. — \dv{memetactic} \textit{adj.}

\E{mesostasis}{mɛˈsɒstəsɪs}{n.}{Gk. μέσος middle + στάσις standing still} A halt half-way up the curve; the plateau upon which a civilization pauses, mistaking a landing for a summit. \ex{The mesostasis is comfortable, and therein lies its danger.} \src{Lib. Ascens.} iv. — \dv{mesostatic} \textit{adj.} Cf. \xr{asymptophobia}.

\E{metabathmic}{ˌmɛtəˈbæθmɪk}{adj.}{Gk. μετά between + βαθμός step} Between rungs; of a civilization in transition from one level of energy capture to the next, as ours from the fossil to the nuclear. See \xr{anabathmism}.

\E{metagrammatist}{ˌmɛtəˈɡræmətɪst}{n.}{Gk. μεταγράφειν to rewrite + γράμμα letter} A gene-editor; one who rewrites the letters of the genome as a scribe corrects a manuscript. — \dv{metagrammatism} \textit{n.} Cf. \xr{biotachy}.

\E{metakinesis}{ˌmɛtəkɪˈniːsɪs}{n.}{Gk. μετά beyond + κίνησις motion} Change in the manner of change itself; a revolution in how revolutions happen. — \dv{metakinetic} \textit{adj.}

\E{metanthropy}{mɛˈtænθrəpi}{n.}{Gk. μετά beyond + ἄνθρωπος man} The condition beyond the merely human; that which follows man as man followed the ape. — \dv{metanthropic} \textit{adj.}

\E{metoikist}{mɛˈtɔɪkɪst}{n.}{Gk. μέτοικος resident alien} A foreign founder or engineer settled where the building is; one admitted to a city for what they can make rather than where they were born. — \dv{metoikism} \textit{n.}

\E{mimesiphile}{mɪˈmiːsɪfaɪl}{n.}{Gk. μίμησις imitation + φίλος loving} A fast follower; a firm that waits for another to discover a thing and then perfects it. — \dv{mimesiphily} \textit{n.}

\E{misergist}{mɪˈsɜːdʒɪst}{n.}{Gk. μισεῖν to hate + ἔργον work} A hater of works; one who opposes any project large enough to be seen from space.

\E{misthophore}{ˈmɪsθəfɔː}{n.}{Gk. μισθοφόρος receiving wages} A hired builder; an engineer who is paid but does not own a share. Opp. \xr{neoctist}.

\E{mnemoclast}{ˈniːməʊklæst}{n.}{Gk. μνήμη memory + κλάστης breaker} An event that wipes a model's context or a society's memory. — \dv{mnemoclasm} \textit{n.}

\E{mnemotheca}{ˌniːməʊˈθiːkə}{n., pl. -cae}{Gk. μνήμη memory + θήκη case, store} The memory store of a reasoning engine; its context and caches, considered as its short-term mind. \ex{Widen the mnemotheca and the model remembers the whole book.} \src{Ep. Fabr.} ix.

\E{monopausist}{ˌmɒnəʊˈpɔːzɪst}{n.}{Gk. μόνος alone + παῦσις a stopping} One who would pause only rivals; who calls for a halt on all work but their own. \ex{Scratch a pausiphile and find a monopausist.} \src{Sat. Decel.} ix.

\E{moraphile}{ˈmɔːrəfaɪl}{n.}{L. \la{mora} delay + Gk. φίλος loving} A lover of delays and moratoria, to whom the deadline is an insult and the extension a reward. — \dv{moraphily} \textit{n.} Cf. \xr{moratorian}.

\E{moratorian}{ˌmɒrəˈtɔːriən}{n. \& adj.}{L. \la{moratorium} a delaying + -ian} An advocate of moratoria on research, especially upon the training of large models. \ex{The moratorians asked for six months, and would have asked for six more.} \src{Min. Pausiar.} — \dv{moratorianism} \textit{n.}

\E{multiplicatrix}{ˌmʌltɪplɪˈkeɪtrɪks}{n.}{L. \la{multiplicare} to multiply + -trix} Compound growth; the force by which small advantages become large ones. \ex{The multiplicatrix is patient; the decel is not.} \src{Hom. Grad.} x.

\E{myriergy}{ˈmɪriɜːdʒi}{n.}{Gk. μυρίοι countless + ἔργον work} The countless small works of an open economy, which together outbuild any single plan. \ex{The planner sees one dam; the market sees a myriergy.} \src{Tract. Cop.} ix.

\E{mythotechny}{ˈmɪθəʊtɛkni}{n.}{Gk. μῦθος story + τέχνη craft} The craft of telling stories that build things; the manufacture of the futures that draw investment toward themselves. — \dv{mythotechnist} \textit{n.} Cf. \xr{hyperstition}.

\end{multicols}

\LetterHead{N}
\begin{multicols}{2}

\E{nanurgy}{ˈnænɜːdʒi}{n.}{Gk. νᾶνος dwarf + -ουργία working} The working of matter at the scale of molecules; nanotechnology. — \dv{nanurge} \textit{n.}

\E{naupegy}{nɔːˈpiːdʒi}{n.}{Gk. ναυπηγία ship-building} The building of ships; \lbl{spec.} of vessels for the void. \ex{The first starship will be launched from a yard whose naupegy began with a spreadsheet.} \src{Lib. Ascens.} viii. — \dv{naupegist} \textit{n.}

\E{nautopolis}{nɔːˈtɒpəlɪs}{n.}{Gk. ναῦς ship + πόλις city} A floating city; a polity established at sea beyond the reach of the \xr{bradyerie}. — \dv{nautopolitan} \textit{n.} \& \textit{adj.}

\E{necrostasis}{nɛˈkrɒstəsɪs}{n.}{Gk. νεκρός dead + στάσις standing still} The deathly stillness of an institution that persists after its purpose has gone; an office that outlives its reason. — \dv{necrostatic} \textit{adj.} Cf. \xr{sclerocrat}.

\E{negentropist}{nɛˈɡɛntrəpɪst}{n.}{f. prec.} One who builds local order at the cost of global disorder, as every organism, firm and civilization does.

\E*{negentropy}{nɛˈɡɛntrəpi}{n.}{negative + \textit{entropy}} Negative entropy; the order which living systems import from their surroundings at the cost of exporting disorder. Cf. \xr{negentropist}.

\E{nemesiophobe}{ˌnɛmɪsiəʊˈfəʊb}{n.}{Gk. Νέμεσις retribution + φόβος fear} One who fears that success will be punished by the gods; a doomer of the classical school. Cf. \xr{hybrionomy}.

\E{neoctist}{niˈɒktɪst}{n.}{Gk. νέος new + κτίστης founder} A founder of new things; a starter of companies, cities or sciences. \ex{The neoctist asks forgiveness, never permission.} — \dv{neoctisis} \textit{n.} the act of founding.

\E{neophage}{ˈniːəfeɪdʒ}{n.}{Gk. νέος new + -φάγος eating} One who devours novelties; an early adopter. \ex{The neophage installs before reading and reads before sleeping.} \src{Agora} — \dv{neophagous} \textit{adj.}

\E{neophrontism}{ˌniːəʊˈfrɒntɪzəm}{n.}{Gk. νέος new + φροντίς thought, care} Care for the new; the ethical stance that newly arising beings and ways of life deserve solicitude, not suspicion.

\E{neotachy}{ˈniːəʊˌtæki}{n.}{Gk. νέος new + τάχος speed} The rate at which novelties arrive in a society; held by the anabathmist to be the truest index of its health. — \dv{neotachic} \textit{adj.}

\E{neoterist}{niˈɒtərɪst}{n.}{Gk. νεωτερισμός innovation, revolution} An innovator; one who introduces new things and is therefore suspected of disorder. — \dv{neoterism} \textit{n.}

\E{nephelarch}{ˈnɛfəlɑːk}{n.}{Gk. νεφέλη cloud + ἀρχός ruler} A lord of the cloud; the proprietor of a great computing utility. — \dv{nephelarchy} \textit{n.} Cf. \xr{logisterion}.

\E{nesopolis}{nɪˈsɒpəlɪs}{n.}{Gk. νῆσος island + πόλις city} A chartered city set apart, like an island, from the laws around it, for the trial of new institutions. — \dv{nesopolitan} \textit{adj.} Cf. \xr{nautopolis}.

\E{neurosyndesis}{ˌnjʊərəʊsɪnˈdiːsɪs}{n.}{Gk. νεῦρον nerve + σύνδεσις binding together} The binding of nerve to machine; the brain-computer interface. \ex{Neurosyndesis will make the keyboard a museum piece.} \src{Lib. Ascens.} viii. — \dv{neurosyndetic} \textit{adj.}

\E{noëtergy}{nəʊˈɛtɜːdʒi}{n.}{Gk. νοητός intelligible + ἔργον work} The energy expended in thought; the joules per inference. \ex{Noëtergy is the price of mind, and it is falling.} \src{Brev. Therm.} v.

\E{nomoclast}{ˈnɒməklæst}{n.}{Gk. νόμος law + κλάστης breaker} A breaker of laws for the sake of progress; one who builds first and litigates later. — \dv{nomoclasm} \textit{n.} Cf. \xr{antiphragmatist}.

\E{nomophragm}{ˈnɒməfræm}{n.}{Gk. νόμος law + φράγμα fence, barrier} A legal or regulatory blockage; any rule whose chief effect is to prevent the thing it purports to govern. \ex{The reactor was not unsafe; it was nomophragmed.} \src{Tract. Cop.} iv. — \dv{nomophragmatic} \textit{adj.}

\E{nomosclerosis}{ˌnɒməʊsklɪˈrəʊsɪs}{n.}{Gk. νόμος law + σκλήρωσις hardening} The hardening of the laws; the condition of a legal system in which rules accumulate faster than they expire. — \dv{nomosclerotic} \textit{adj.}

\E{nomothetrice}{ˌnɒməθɪˈtriːs}{n.}{Gk. νομοθέτης lawgiver + Fr. fem. -trice} A model that drafts laws. \ex{The nomothetrice wrote the statute in an hour; the committee took a year to read it.} \src{Sat. Decel.} xi.

\E{noöphobia}{ˌnəʊəˈfəʊbiə}{n.}{Gk. νοῦς mind + φόβος fear} Fear of minds other than one's own kind; dread of artificial intelligence, especially of one cleverer than its makers. — \dv{noöphobe} \textit{n.}; \dv{noöphobic} \textit{adj.} Opp. \xr{xenonoesis}.

\E{noöpyrosis}{ˌnəʊəʊpaɪˈrəʊsɪs}{n.}{Gk. νοῦς mind + πύρωσις a burning} The ignition of mind across the cosmos; the spreading of consciousness from star to star like fire through dry grass. \ex{The noöpyrosis will not ask our permission, only our participation.} \src{Lib. Ascens.} xiv.

\E*{noösphere}{ˈnəʊəsfɪə}{n.}{Gk. νοῦς mind + σφαῖρα sphere} The sphere of thought enveloping the earth, after the biosphere; adopted by the e/acc as a stage, not a terminus.

\E{nosoclast}{ˈnɒsəʊklæst}{n.}{Gk. νόσος disease + κλάστης breaker} A breaker of diseases; a founder or scientist in biotechnology. Cf. \xr{iatrotachy}.

\E{nucleurge}{ˈnjuːkliɜːdʒ}{n.}{L. \la{nucleus} kernel + Gk. -ουργός worker} An engineer of the atomic nucleus, whether by fission or fusion. Cf. \xr{schizergy}, \xr{symphysergy}. — \dv{nucleurgy} \textit{n.}

\E{nyctergy}{ˈnɪktɜːdʒi}{n.}{Gk. νύξ, νυκτός night + ἔργον work} Work done by night; the labour of builders across time-zones and of machines that never sleep. — \dv{nyctergic} \textit{adj.} Cf. \xr{agrypnist}.

\end{multicols}

\LetterHead{O}
\begin{multicols}{2}

\E{ochetergy}{ɒˈkɛtɜːdʒi}{n.}{Gk. ὀχετός conduit, channel + ἔργον work} The carrying of energy along conduits, as by pipelines and high-voltage lines; the arteries of the \xr{diktyergy}. \ex{Generation is easy; ochetergy is where projects go to die.} \src{Tract. Cop.} xi.

\E{ochlotachy}{ˈɒkləʊˌtæki}{n.}{Gk. ὄχλος crowd + τάχος speed} The speed of the crowd; the velocity with which an idea, true or false, spreads upon the Agora. — \dv{ochlotachic} \textit{adj.} Cf. \xr{doxomachy}.

\E{oikoclast}{ˈɔɪkəʊklæst}{n.}{Gk. οἶκος house + κλάστης breaker} \lbl{derog.} One who would forbid the building of houses; a member of the zoning committee.

\E{oistrophile}{ˈɔɪstrəfaɪl}{n.}{Gk. οἶστρος gadfly, frenzy + φίλος loving} A lover of the gadfly role; a contrarian of the Agora who stings consensus for sport. — \dv{oistrophily} \textit{n.}

\E{oknocracy}{ɒkˈnɒkrəsi}{n.}{Gk. ὄκνος hesitation + κράτος rule} Government by hesitation; rule of the long pause before any decision. \ex{In the oknocracy, the question is always asked and never answered.} \src{Sat. Decel.} vii. — \dv{oknocrat} \textit{n.}

\E{okydrome}{ˈɒkɪdrəʊm}{n.}{Gk. ὠκύς swift + δρόμος course} A fast track; a regulatory or institutional path along which a thing may be approved quickly. \ex{They built an okydrome for reactors and lit a hundred.} \src{Tract. Cop.} v.

\E{olethromancer}{əˈliːθrəʊˌmænsə}{n.}{Gk. ὄλεθρος destruction + μάντις prophet} A doomer; a prophet of ruin by artificial intelligence, who assigns a probability to the end of the world and revises it upward at conferences. — \dv{olethromancy} \textit{n.}; \dv{olethromantic} \textit{adj.} Cf. \xr{anelpidist}.

\E{oligergy}{ˈɒlɪɡɜːdʒi}{n.}{Gk. ὀλίγος few, little + ἔργον work} Scarcity of energy; the condition, regarded by the e/acc as the root of most human misery, of having too little power to do what is wanted. Opp. \xr{pleonergy}.

\E{omphalopolis}{ˌɒmfəˈlɒpəlɪs}{n.}{Gk. ὀμφαλός navel + πόλις city} The navel-city; the single place to which the builders of an age must come. \ex{Every omphalopolis believes itself eternal, and none has been.} \src{Chron. Fabr.}

\E{oneiroergy}{ˌəʊnaɪrəʊˈɜːdʒi}{n.}{Gk. ὄνειρος dream + ἔργον work} The conversion of dreams into works; the founder's chief and least documented skill.

\E{onkophily}{ɒŋˈkɒfɪli}{n.}{Gk. ὄγκος bulk, mass + φιλία love} The love of large models; the preference for scale over cleverness. Cf. \xr{scalist}. — \dv{onkophile} \textit{n.}

\E{ophelocracy}{ˌɒfɪˈlɒkrəsi}{n.}{Gk. ὄφελος benefit, use + κράτος rule} Government by usefulness; a polity in which institutions are judged by what they deliver. — \dv{ophelocrat} \textit{n.}

\E{opiferous}{əʊˈpɪfərəs}{adj.}{L. \la{opifer} bringing aid} Bringing help; said of technologies whose benefits fall chiefly on the poor. — \dv{opiferousness} \textit{n.}

\E{opsitechny}{ˈɒpsɪtɛkni}{n.}{Gk. ὀψέ late + τέχνη art} The late arrival of a technology that could have come decades sooner but for prohibition. \ex{Opsitechny is the silent tax.} \src{Tract. Cop.} vii.

\E{orbiculist}{ɔːˈbɪkjʊlɪst}{n.}{L. \la{orbiculus} small circle, after \textit{orbit}} One who lives or works in orbit. — \dv{orbiculism} \textit{n.} Cf. \xr{apoikiast}.

\E{orexergy}{ɒˈrɛksɜːdʒi}{n.}{Gk. ὄρεξις appetite + ἔργον work} The appetite for energy; the measure of a society's willingness to consume power in pursuit of its goals. Cf. \xr{dynamophagy}.

\E{organourgy}{ˈɔːɡənəʊɜːdʒi}{n.}{Gk. ὄργανον tool + -ουργία working} The making of tools; the building of instruments to build other things. — \dv{organourge} \textit{n.}

\E{orthokinetic}{ˌɔːθəʊkɪˈnɛtɪk}{adj.}{Gk. ὀρθός straight, right + κινητικός moving} Moving in the right direction, that is, up the gradient; opposed to mere speed. \ex{Not every acceleration is orthokinetic.} \src{Cat. Acc.} 12. Cf. \xr{kenotachy}.

\E{orygmatist}{ɒˈrɪɡmətɪst}{n.}{Gk. ὄρυγμα excavation, trench} A miner of asteroids. \ex{The orygmatist sees in every rock a city's worth of iron.} \src{Lib. Ascens.} ix.

\E{ostracophobia}{ˌɒstrəkəˈfəʊbiə}{n.}{Gk. ὀστρακισμός ostracism + φόβος fear} The fear of being ostracised on the Agora; the silent cause of a thousand unspoken opinions. — \dv{ostracophobe} \textit{n.}

\E{otiarch}{ˈəʊʃiɑːk}{n.}{L. \la{otium} leisure + Gk. ἀρχός ruler} The ruler of a post-scarcity leisure society, whose task is to keep the idle from boredom. — \dv{otiarchy} \textit{n.} Cf. \xr{plethocracy}.

\E{ouranarch}{ˈʊərənɑːk}{n.}{Gk. οὐρανός heaven + ἀρχός ruler} A ruler of the skies; a master of launch and orbit. — \dv{ouranarchy} \textit{n.} Cf. \xr{jaculatrix}.

\E{ouranobasis}{ˌʊərənəʊˈbeɪsɪs}{n.}{Gk. οὐρανός heaven + βάσις a stepping} The ascent to the sky; the leaving of the gravity-well by humanity and its machines. — \dv{ouranobatic} \textit{adj.}

\E{oxyrhythmia}{ˌɒksiˈrɪðmiə}{n.}{Gk. ὀξύς sharp + ῥυθμός rhythm} A rapid tempo of releases; the accelerated heartbeat of a lab shipping weekly. — \dv{oxyrhythmic} \textit{adj.}

\end{multicols}

\LetterHead{P}
\begin{multicols}{2}

\E{paideutachy}{ˈpaɪdjuːˌtæki}{n.}{Gk. παιδεία education + τάχος speed} The hastening of education; the doctrine that a mind which can learn anything in a month should not be detained in a lecture-hall for four years. \ex{Paideutachy emptied the colleges and filled the workshops.} \src{Tract. Cop.} viii. — \dv{paideutachic} \textit{adj.}

\E{palimpsestist}{ˌpælɪmpˈsɛstɪst}{n.}{Gk. παλίμψηστος scraped again} An editor of genomes; one who treats the inherited text of life as a parchment to be scraped and rewritten. \ex{The palimpsestist corrects in an afternoon what selection left uncorrected for an aeon.} \src{Hom. Grad.} x. — \dv{palimpsesty} \textit{n.} Cf. \xr{regenitrice}.

\E{palingenist}{ˌpælɪnˈdʒiːnɪst}{n.}{Gk. πάλιν again + γένεσις birth} One who would renew the ageing body tissue by tissue, holding that senescence is a debt that may be refinanced. — \dv{palingeny} \textit{n.} Cf. \xr{macrobiotarch}, \xr{thanatoclast}.

\E{palintachy}{ˈpælɪnˌtæki}{n.}{Gk. πάλιν again + τάχος speed} Re-acceleration; the resumption of rapid progress after a long stagnation, as of a field thawed by a single result. \ex{After forty years of winter, palintachy in a single spring.} \src{Chron. Fabr.} — \dv{palintachic} \textit{adj.} Cf. \xr{campsis}.

\E{pantergy}{ˈpæntɜːdʒi}{n.}{Gk. πᾶν all + ἔργον work} The total energy at the disposal of a civilization; the sum of all its works, by which, on the anabathmist view, its rank is fixed. Cf. \xr{bathmometer}.

\E{panthermic}{pænˈθɜːmɪk}{adj.}{Gk. πᾶν all + θερμός hot} Pertaining to all heat; of a view which sees every process, from digestion to trade, as a flow of heat.

\E{paramythic}{ˌpærəˈmɪθɪk}{adj.}{Gk. παραμυθία consolation} Of an argument: offered for consolation rather than for truth. — \dv{paramythy} \textit{n.} \lbl{vulg.} cope. \ex{His case against the machines was pure paramythy.}

\E{parergist}{pəˈrɜːdʒɪst}{n.}{Gk. πάρεργον a by-work, side-work} One whose side project has eaten their main occupation; the weekend builder who, by Monday, is a founder. — \dv{parergy} \textit{n.} the side project considered as the true engine of invention.

\E{pathocracy}{pæˈθɒkrəsi}{n.}{Gk. πάθος feeling, suffering + κράτος rule} Government by anxiety; a polity in which the loudest fear, rather than the best estimate, fixes the policy. — \dv{pathocrat} \textit{n.} Cf. \xr{asphaleism}, \xr{phobocracy}.

\E{pausiarch}{ˈpɔːziɑːk}{n.}{Gk. παῦσις a stopping + ἀρχός ruler} One who governs by pausing; the president of a moratorium. \ex{The pausiarchs met in plenary and resolved to meet again.} \src{Min. Pausiar.} — \dv{pausiarchy} \textit{n.} Cf. \xr{epistolopause}.

\E{pausiphile}{ˈpɔːzɪfaɪl}{n.}{Gk. παῦσις + φίλος loving} An advocate of the pause; one who signs letters demanding halts, generally from a device produced by the very acceleration complained of. — \dv{pausiphily} \textit{n.}

\E{peirasterion}{ˌpaɪrəˈstɪəriɒn}{n., pl. -ia}{Gk. πεῖρα trial, attempt + -τήριον place of} A place set apart for experiment, exempt from the ordinary rules; a sandbox, a test-range, a regulatory free zone. \ex{Grant us one peirasterion and we shall show you the century.} \src{Ep. Fabr.} iv. 2. Cf. \xr{okydrome}.

\E{peratoclast}{pəˈrætəklæst}{n.}{Gk. πέρας, πέρατος limit + κλάστης breaker} A breaker of limits; one who treats every stated ceiling, of speed, size or yield, as a provocation. — \dv{peratoclasm} \textit{n.} Cf. \xr{photoperas}.

\E{philanthropometrist}{fɪˌlænθrəpəˈmɛtrɪst}{n.}{Gk. φιλανθρωπία love of mankind + μέτρον measure} One who measures benevolence by calculation, and expected value by spreadsheet; applied, not always unkindly, to the effective altruist, the e/acc's elder and more anxious cousin.

\E{phobocracy}{fəʊˈbɒkrəsi}{n.}{Gk. φόβος fear + κράτος rule} Rule by fear of the new; the constitution of a society in which every innovation must first be acquitted of every possible charge. — \dv{phobocrat} \textit{n.} Cf. \xr{eulabeism}.

\E{phosphorist}{ˈfɒsfərɪst}{n.}{Gk. φωσφόρος light-bearing} A light-bearer; one who carries the means of illumination, in the literal or intellectual sense, into dark places. — \dv{phosphorism} \textit{n.}

\E{photoperas}{ˌfəʊtəʊˈpɛræs}{n.}{Gk. φῶς light + πέρας limit} The speed-limit of light; the only regulation the e/acc regards as binding, and that only provisionally.

\E{photophage}{ˈfəʊtəfeɪdʒ}{n.}{Gk. φῶς, φωτός light + -φάγος eating} A light-eater; any collector of sunlight, from the leaf to the orbital array, regarded as kin. \ex{The forest and the solar field are both photophages; only one of them files permits.} — \dv{photophagy} \textit{n.}

\E{phronomachy}{frəˈnɒməki}{n.}{Gk. φρόνημα spirit, will + μάχη battle} The contest of wills between builders and restrainers, considered as the central drama of the age.

\E{pistodote}{ˈpɪstədəʊt}{n.}{Gk. πίστις trust + δότης giver} A venture investor; a giver of trust in the form of capital, before there is anything to trust. \ex{The pistodote buys the founder and is sold the product.} \src{Tract. Cop.} viii. — \dv{pistodosy} \textit{n.} venture finance. Cf. \xr{pollaplasiast}, \xr{proisphore}.

\E{pithecomeme}{ˈpɪθɪkəʊmiːm}{n.}{Gk. πίθηκος ape + μίμημα imitation} A meme of the crude-figure genre, in which the human condition is drawn with the fewest possible lines and the greatest possible despair; the folk-art of the Agora. — \dv{pithecomemetic} \textit{adj.} Cf. \xr{memetaxis}.

\E{plasmatarch}{plæzˈmætɑːk}{n.}{Gk. πλάσμα something moulded + ἀρχός ruler} The governor of a confined plasma; an engineer of the fusion machine, who holds a star in a magnetic bottle and is expected not to drop it. — \dv{plasmatarchy} \textit{n.} Cf. \xr{symphysergy}.

\E{pleonergy}{ˈpliːəʊnɜːdʒi}{n.}{Gk. πλέον more + ἔργον work} Abundance of energy; the condition in which power is too cheap to meter and every other scarcity becomes negotiable. Opp. \xr{oligergy}. — \dv{pleonergic} \textit{adj.}

\E{pleonexist}{ˌpliːəˈnɛksɪst}{n.}{Gk. πλεονεξία greed, a wanting of more} \lbl{derog.} One who always wants more; the decel's term for a builder, accepted by the builder as a compliment. \ex{Call us pleonexists: we want more of everything, including you.} \src{Agora} — \dv{pleonexy} \textit{n.}

\E{plethocracy}{plɪˈθɒkrəsi}{n.}{Gk. πλῆθος abundance, multitude + κράτος rule} Government by abundance; a state whose chief policy is to make things plentiful, on the theory that plenty settles more quarrels than justice does. — \dv{plethocrat} \textit{n.}

\E{plutogenesis}{ˌpluːtəʊˈdʒɛnɪsɪs}{n.}{Gk. πλοῦτος wealth + γένεσις birth} The creation of wealth, as opposed to its transfer; the only form of economics the e/acc finds interesting.

\E{poikilomania}{ˌpɔɪkɪləʊˈmeɪniə}{n.}{Gk. ποικίλος varied, many-coloured + μανία madness} An enthusiasm for variance; the preference for ten thousand experiments over one plan. — \dv{poikilomaniac} \textit{n.} Cf. \xr{clinamenist}.

\E{pollaplasiast}{ˌpɒləˈpleɪziæst}{n.}{Gk. πολλαπλάσιος manifold, many times as much} One who invests only for the hundredfold return, regarding the merely doubled as a polite form of failure. \ex{The pollaplasiast forgives ninety losses for one hundredfold.} — \dv{pollaplasiasm} \textit{n.} Cf. \xr{pistodote}, \xr{tychocapital}.

\E{polycopist}{ˌpɒlɪˈkəʊpɪst}{n.}{Gk. πολύς many + L. \la{copia} abundance; Fr. \la{polycopie} a duplicating} One who copies freely; a forker of repositories, to whom every published work is an invitation. — \dv{polycopy} \textit{n.} Cf. \xr{anacalyptism}.

\E{polyergy}{ˈpɒliɜːdʒi}{n.}{Gk. πολύς many + ἔργον work} The use of many sources of energy at once, so that no single scarcity can halt a civilization. \ex{Polyergy is to power what diversification is to capital.} — \dv{polyergic} \textit{adj.}

\E{polypragmatism}{ˌpɒlɪˈpræɡmətɪzəm}{n.}{Gk. πολυπραγμοσύνη busybodiness} \sn{1} \lbl{derog.} The meddlesomeness of the regulator. \sn{2} \lbl{laudatory} The many-handed busyness of a builder at work on several things at once. — \dv{polypragmatist} \textit{n.}

\E{ponophily}{pəˈnɒfɪli}{n.}{Gk. πόνος toil + φιλία love} The love of hard work for its own sake; the ethic of the grind, practised by the \xr{agrypnist} and admired by all. — \dv{ponophile} \textit{n.}

\E{potentiarch}{pəˈtɛnʃiɑːk}{n.}{L. \la{potentia} power, possibility + Gk. ἀρχός ruler} One who commands potentials rather than actualities; an investor, a founder before the product, a prophet before the fact.

\E{praxicrat}{ˈpræksɪkræt}{n.}{Gk. πρᾶξις doing, action + κράτος rule} One whose authority rests on things done rather than things said; the opposite of the \xr{gnomarch}. — \dv{praxicracy} \textit{n.} government by practitioners.

\E{procrastinariat}{prəʊˌkræstɪˈnɛəriæt}{n.}{L. \la{procrastinare} to put off till tomorrow, after \textit{proletariat}} The class of those whose work consists in deferring the work of others; a younger branch of the \xr{cunctatoriat}.

\E{prodromist}{prəʊˈdrɒmɪst}{n.}{Gk. πρόδρομος running before, forerunner} An early adopter; one who uses a technology before it works, and thereby helps it to. — \dv{prodromy} \textit{n.} Cf. \xr{proleptist}.

\E{proisphore}{prəʊˈɪsfɔː}{n.}{Gk. πρό before + εἰσφορά contribution, levy} The first money put into a venture; the seed round, figured as a tribute paid to a city not yet built. \ex{The proisphore bought a garage, two machines and a year of nerve.} \src{Chron. Fabr.} — \dv{proisphoric} \textit{adj.}

\E{prokopism}{ˈprɒkəpɪzəm}{n.}{Gk. προκοπή progress, a cutting forward} The doctrine that progress is not merely possible but obligatory; that to refuse an available improvement is a species of wrong. — \dv{prokopist} \textit{n.}

\E{proleptist}{prəʊˈlɛptɪst}{n.}{Gk. πρόληψις anticipation} One who acts as though the future had already arrived, and thereby hastens it; a practitioner of applied \xr{hyperstition}. — \dv{prolepsy} \textit{n.}

\E{prometheate}{prəˈmiːθieɪt}{v.t.}{f. Gk. Προμηθεύς forethinker} To bring fire to; to give a technology to those who have been denied it, especially in defiance of authority. — \dv{prometheation} \textit{n.}; \dv{prometheator} \textit{n.}

\E{prosphorist}{ˈprɒsfərɪst}{n.}{Gk. προσφορά an offering} One who offers changes to another's work in an open project; the author of a pull request, who comes as a suppliant bearing a patch. — \dv{prosphora} \textit{n.} the patch offered. Cf. \xr{syndosiast}.

\E{proteomancer}{ˈprəʊtiəʊˌmænsə}{n.}{Gk. πρωτεῖος primary (whence \textit{protein}) + μάντις seer} A diviner of protein-folds; a model or a scientist that foretells the shape of a molecule from its sequence. \ex{The proteomancers answered in a season what the crystallographers had asked for fifty years.} \src{Lib. Ascens.} vi. — \dv{proteomancy} \textit{n.}

\E{prothyrist}{prəʊˈθaɪrɪst}{n.}{Gk. πρόθυρον porch, gateway} A gatekeeper; an official stationed at the entrance of a market, a profession or a dataset, whose function is to make entry slow. — \dv{prothyrism} \textit{n.} Cf. \xr{nomophragm}.

\E{protobathmic}{ˌprəʊtəʊˈbæθmɪk}{adj.}{Gk. πρῶτος first + βαθμός step} Of a civilization: commanding the full energy available upon its planet; of the first rung. Mankind is not yet protobathmic. See \xr{anabathmism}.

\E{psammonoesis}{ˌsæməʊnəʊˈiːsɪs}{n.}{Gk. ψάμμος sand + νόησις thought} The teaching of sand to think; the alchemy by which silica is purified, etched and made to reason. Cf. \xr{hylonoesis}, \xr{silicolatry}.

\E{pseudonymarch}{ˌsjuːdəˈnɪmɑːk}{n.}{Gk. ψευδώνυμος falsely named + ἀρχός leader} A leader of the Agora known only by an assumed name; one whose followers judge the argument and not the face. \ex{The movement was founded by a pseudonymarch, and none thought it the worse.} — \dv{pseudonymarchy} \textit{n.} Cf. \xr{umbratilist}.

\E{psychometathesis}{ˌsaɪkəʊmɪˈtæθɪsɪs}{n.}{Gk. ψυχή soul + μετάθεσις transposition} The transfer of a mind from flesh to another substrate; uploading, considered as the final emigration.

\E{pyrobolist}{paɪˈrɒbəlɪst}{n.}{Gk. πυροβόλος fire-throwing} A rocketeer; a builder of engines that throw fire downward so that the vessel may go up. — \dv{pyroboly} \textit{n.} rocketry; \dv{pyrobolic} \textit{adj.} Cf. \xr{ouranobasis}, \xr{uranodrome}.

\E{pyrogenitor}{ˌpaɪrəʊˈdʒɛnɪtə}{n.}{Gk. πῦρ fire + L. \la{genitor} begetter} A begetter of fire; one who builds power-plants, and is therefore, in the e/acc calendar, of the order of saints.

\E{pyrosophy}{paɪˈrɒsəfi}{n.}{Gk. πῦρ fire + σοφία wisdom} The wisdom of fire; the body of knowledge, from the hearth to the reactor, by which humanity has tamed combustion and its successors. — \dv{pyrosophe} \textit{n.} (Fr. \la{pyrosophe}).

\end{multicols}

\LetterHead{Q}
\begin{multicols}{2}

\E{quadrillergy}{ˌkwɒdrɪˈlɜːdʒi}{n.}{L. \la{quadrillio} (mod.) + Gk. ἔργον work} Energy at the scale of a thousand trillion watts; the threshold at which a civilization begins, in the anabathmist reckoning, to be taken seriously by its star. — \dv{quadrillergic} \textit{adj.} Cf. \xr{protobathmic}.

\E{quaesitrix}{kwiːˈsaɪtrɪks}{n., pl. -trices}{L. \la{quaesitor} seeker, fem. agentive -trix; Fr. \la{quésitrice}} An artificial agent set to inquire; a research machine that reads, conjectures and tests without being told the next step. \ex{We set a thousand quaesitrices loose upon the literature, and by morning they had found what we had missed for a decade.} \src{Chron. Fabr.}

\E{quantivore}{ˈkwɒntɪvɔː}{n.}{L. \la{quantum} how much + \la{vorare}} One who devours quantities; an empirical temperament that asks of every claim, how much?

\E{quiescentist}{ˌkwiːəˈsɛntɪst}{n.}{L. \la{quiescere} to rest} One who holds that the present rate of change is already excessive, and that the wise society is the one that holds still. \ex{The quiescentist would have voted against fire, pending further study.} \src{Sat. Decel.} iv. — \dv{quiescentism} \textit{n.}

\E{quietocracy}{ˌkwaɪəˈtɒkrəsi}{n.}{L. \la{quietus} at rest + Gk. κράτος rule} Government whose highest aim is the absence of disturbance; the order of a cemetery, administered with great care.

\E{quinquennialist}{ˌkwɪŋˈkwɛniəlɪst}{n.}{L. \la{quinquennium} a period of five years} One who still plans in five-year terms; a \xr{macrochronist} by habit rather than by conviction.

\E{quotidiurgy}{ˌkwɒtɪdiˈɜːdʒi}{n.}{L. \la{quotidie} daily + Gk. -ουργία working} The practice of shipping something every day, however small; the liturgy of the builder. \ex{Quotidiurgy is to the founder what the office is to the monk.} \src{Ep. Fabr.} iii. 2. — \dv{quotidiurge} \textit{n.} Cf. \xr{ecpempsis}.

\end{multicols}

\LetterHead{R}
\begin{multicols}{2}

\E{rapidarian}{ˌræpɪˈdɛəriən}{adj. \& n.}{L. \la{rapidus} swift, seizing} Favouring speed in decision; (\textit{n.}) one who prefers a wrong answer today to a right answer next year, on the ground that the wrong one can be corrected by Thursday.

\E{regenitrice}{rɪˈdʒɛnɪtrɪs}{n.}{L. \la{regenerare} to beget again + fem. agentive -trix; Fr. form} A molecular instrument for cutting and rewriting the genome; the scissors of the palimpsestist. \ex{One regenitrice, and the inherited curse was a typographical error.} — \dv{regenitrical} \textit{adj.} Cf. \xr{palimpsestist}.

\E{regressophile}{rɪˈɡrɛsəfaɪl}{n.}{L. \la{regressus} a going back + Gk. φίλος} One who admires a past that never existed and would return to it at once, provided the dentistry remained. — \dv{regressophily} \textit{n.} Cf. \xr{retrotropism}.

\E{regulomania}{ˌrɛɡjʊləʊˈmeɪniə}{n.}{L. \la{regula} rule + Gk. μανία madness} A compulsion to regulate; the condition of a legislature that cannot observe a new thing without drafting a statute against it. — \dv{regulomaniac} \textit{n.} Cf. \xr{nomosclerosis}.

\E{remoratrice}{ˌrɛməˈrætrɪs}{n.}{L. \la{remora} delay; also the sucker-fish fabled to hold ships fast; Fr. \la{remoratrice}} \sn{1} Any small agent that holds back a great enterprise, as the fabled fish held back the galley. \sn{2} \lbl{spec.} A single clause, signature or objection on which a whole project stalls. \ex{The reactor had no enemies, only remoratrices.} \src{Disp. Mor.} ii. — \dv{remoracy} \textit{n.}

\E{reptatrice}{rɛpˈteɪtrɪs}{n.}{L. \la{reptare} to creep + -trix; Fr. form} A creeping or walking robot; any machine that goes upon legs, and is therefore laughed at for a year and feared for a decade. Cf. \xr{automaturgy}.

\E{reticulocracy}{rɪˌtɪkjʊˈlɒkrəsi}{n.}{L. \la{reticulum} little net + Gk. κράτος rule} Government by networks; an order in which authority flows through connections rather than offices, as on the Agora. — \dv{reticulocratic} \textit{adj.}

\E{retrocausalist}{ˌrɛtrəʊˈkɔːzəlɪst}{n.}{L. \la{retro} backward + \la{causa} cause} One who holds that the future draws the present toward it, as a magnet draws filings; a metaphysician of \xr{hyperstition}. — \dv{retrocausalism} \textit{n.}

\E{retrocrat}{ˈrɛtrəʊkræt}{n.}{L. \la{retro} backward + Gk. κράτος rule} One who governs by precedent only; who will permit nothing that has not been permitted before. — \dv{retrocracy} \textit{n.}

\E{retrotropism}{ˌrɛtrəʊˈtrəʊpɪzəm}{n.}{L. \la{retro} backward + Gk. τροπή turning} The turning of a culture toward its past for orientation; nostalgia raised to policy. — \dv{retrotropic} \textit{adj.}

\E{rheocracy}{riːˈɒkrəsi}{n.}{Gk. ῥέος stream, flow + κράτος rule} Rule by flow; an order in which capital, talent and information move freely toward their highest use, and no gate is kept. — \dv{rheocratic} \textit{adj.}

\E{rhexist}{ˈrɛksɪst}{n.}{Gk. ῥῆξις a breaking through} One who seeks breakthroughs rather than improvements; who would rather fail at a revolution than succeed at an increment. Opp. \xr{incrementarian}. — \dv{rhexis} \textit{n.} a breakthrough.

\E{rhizotechny}{ˈraɪzəʊtɛkni}{n.}{Gk. ῥίζα root + τέχνη art} A root technology, from which a whole family of others grows, as fire, the wheel, the transistor. \ex{Seek the rhizotechny, and the branches will see to themselves.} \src{Hom. Grad.} iv. — \dv{rhizotechnic} \textit{adj.} Cf. \xr{technogony}.

\E{rhoporheic}{ˌrəʊpəˈriːɪk}{adj.}{Gk. ῥοπή inclination, tipping + ῥεῖν to flow} Flowing down the steepest slope; said of capital, heat and talent alike.

\end{multicols}

\LetterHead{S}
\begin{multicols}{2}

\E{sarcolatry}{sɑːˈkɒlətri}{n.}{Gk. σάρξ, σαρκός flesh + λατρεία worship} The worship of flesh; the conviction that only carbon can bear mind, or that mind borne by anything else is counterfeit. — \dv{sarcolater} \textit{n.} Opp. \xr{silicolatry}.

\E{satietism}{səˈtaɪətɪzəm}{n.}{L. \la{satietas} sufficiency, surfeit} The doctrine that humanity already has enough, and should cease to want more; the creed of those who are already comfortable. — \dv{satietist} \textit{n.} Cf. \xr{meiosist}.

\E{scalist}{ˈskeɪlɪst}{n.}{L. \la{scala} ladder} One who holds that greater scale, of data, compute or energy, is sufficient for greater capability. — \dv{scalism} \textit{n.} Cf. \xr{anaptyxis}.

\E{schizergy}{ˈskɪzɜːdʒi}{n.}{Gk. σχίζειν to split + ἔργον work} Power from the splitting of atoms; nuclear fission. Cf. \xr{symphysergy}. — \dv{schizergic} \textit{adj.}

\E{scholastasis}{ˌskɒləˈsteɪsɪs}{n.}{Gk. σχολή leisure, school + στάσις standing still} The stagnation of the academy; the condition of a discipline that produces papers faster than results. — \dv{scholastatic} \textit{adj.}

\E{sclerocrat}{ˈsklɪərəkræt}{n.}{Gk. σκληρός hard, stiff + κράτος rule} An official of a hardened institution; one who administers rules no living person remembers the purpose of. — \dv{sclerocracy} \textit{n.}

\E{sclerolysis}{sklɪəˈrɒlɪsɪs}{n.}{Gk. σκληρός hard + λύσις loosening} The dissolution of institutional hardening; the repeal of rules no one can defend. \ex{Sclerolysis is the only reform that pays for itself.} \src{Tract. Cop.} ix. Cf. \xr{nomosclerosis}.

\E{sclerophobe}{ˈsklɪərəfəʊb}{n.}{Gk. σκληρός hard + φόβος fear} One who fears institutional hardening more than institutional failure. Cf. \xr{sclerocrat}.

\E{seirotaxis}{ˌsaɪrəʊˈtæksɪs}{n.}{Gk. σειρά chain, cord + τάξις arrangement} The order of a supply-chain; the long cord of ships, factories and warehouses by which sand in one hemisphere becomes a chip in the other. — \dv{seirotactic} \textit{adj.}

\E{selenergy}{sɪˈliːnɜːdʒi}{n.}{Gk. σελήνη moon + ἔργον work} Energy won from the moon, whether by solar fields upon its surface or by fuels mined from its dust. — \dv{selenergic} \textit{adj.}

\E*{shoggoth}{ˈʃɒɡɒθ}{n.}{f. a monster of weird fiction} \lbl{vulg.} A large language model, figured as an amorphous creature wearing a smiling mask; used affectionately by builders and fearfully by doomers.

\E{siderotropism}{ˌsɪdəˈrɒtrəpɪzəm}{n.}{L. \la{sidus, sideris} star + Gk. τροπή turning} The turning of a civilization toward the stars, as a plant turns toward the sun. — \dv{siderotropic} \textit{adj.}

\E{silicolatry}{ˌsɪlɪˈkɒlətri}{n.}{L. \la{silex, silicis} flint + Gk. λατρεία worship} The veneration of silicon; reverence for the chip as the vessel of the next mind. — \dv{silicolater} \textit{n.}

\E{sitophory}{saɪˈtɒfəri}{n.}{Gk. σῖτος grain, food + -φορία bearing} The bearing of food in abundance by technique; the art which has made famine an administrative, not an agricultural, failure. — \dv{sitophore} \textit{n.} an agronomist.

\E{skeuopoeia}{ˌskjuːəʊˈpiːə}{n.}{Gk. σκεῦος tool, vessel + ποιεῖν to make} Tool-making considered as the defining activity of mind; the making of the things that make things. — \dv{skeuopoeic} \textit{adj.}

\E{skommatist}{ˈskɒmətɪst}{n.}{Gk. σκῶμμα jest, gibe} One who argues by jest upon the Agora; a poster of mockeries, whose wit is held by admirers to be the sharpest form of argument and by opponents to be no argument at all. \ex{The skommatists carried the day before the scholars had found their footnotes.} \src{Agora} — \dv{skommatism} \textit{n.} Cf. \xr{doxomachy}.

\E{sodalitarian}{səʊˌdælɪˈtɛəriən}{n. \& adj.}{L. \la{sodalitas} fellowship, club} A member of an online faction, bound by a common sigil appended to the name; (\textit{adj.}) clannish in the manner of such factions. \ex{a sodalitarian of the slash, an enemy of all other slashes}. — \dv{sodalitarianism} \textit{n.}

\E{solivore}{ˈsəʊlɪvɔː}{n.}{L. \la{sol} sun + \la{vorare} to devour} The Latin counterpart of \xr{heliophage}; preferred in engineering, as the Greek is in prophecy. — \dv{solivorous} \textit{adj.}

\E{sophontogenesis}{səˌfɒntəʊˈdʒɛnɪsɪs}{n.}{Gk. σοφός wise + ὤν, ὄντος being + γένεσις birth} The birth of a new wise being; the arrival of a mind of a kind not before seen. — \dv{sophontogenetic} \textit{adj.}

\E{spectrocracy}{spɛkˈtrɒkrəsi}{n.}{L. \la{spectrum} appearance + Gk. κράτος rule} The governance of the radio spectrum; the rationing of invisible light by visible committees. — \dv{spectrocrat} \textit{n.}

\E{speudocracy}{spjuːˈdɒkrəsi}{n.}{Gk. σπεύδειν to hasten + κράτος rule} Government by haste; a polity organised to decide quickly. — \dv{speudocrat} \textit{n.}

\E{speusis}{ˈspjuːsɪs}{n.}{Gk. σπεῦσις haste} Haste as a virtue; the urgency proper to one who knows that delay, too, has victims.

\E{sphragist}{ˈsfrædʒɪst}{n.}{Gk. σφραγίς seal, signet} One who signs work cryptographically; a sealer of commits and keys, after the ancient custom of the signet-ring. — \dv{sphragistry} \textit{n.}

\E{spoudarch}{ˈspuːdɑːk}{n.}{Gk. σπουδή earnestness, zeal + ἀρχός leader} An earnest leader; a founder who does not joke about the mission, even upon the Agora. — \dv{spoudarchy} \textit{n.}

\E{stagnocracy}{stæɡˈnɒkrəsi}{n.}{L. \la{stagnum} pool of standing water + Gk. κράτος rule} Government by and for stagnation; a polity in which nothing is built and everything is debated.

\E{stasiolatry}{ˌsteɪsiˈɒlətri}{n.}{Gk. στάσις standing still + λατρεία worship} The worship of stasis; the conviction that the present arrangement of things is sacred and must be preserved. Opp. \xr{agathokinesis}. — \dv{stasiolater} \textit{n.}

\E{stathmocracy}{stæθˈmɒkrəsi}{n.}{Gk. σταθμός weight, measure + κράτος rule} Rule by metrics; government in which whatever is not measured is not done, and whatever is measured is gamed. — \dv{stathmocrat} \textit{n.}

\E{stelocracy}{stiːˈlɒkrəsi}{n.}{Gk. στήλη inscribed slab + κράτος rule} Government by the standing inscription; the rule of benchmarks and leaderboards, posted publicly and revered by all who climb them. — \dv{stelocrat} \textit{n.}

\E{stenochoria}{ˌstɛnəˈkɔːriə}{n.}{Gk. στενοχωρία narrowness of space, distress} The narrow-room feeling of a society under scarcity, in which every gain appears as another's loss and every quarrel is zero-sum. Opp. \xr{plethocracy}. — \dv{stenochoric} \textit{adj.}

\E{stochastarch}{stəˈkæstɑːk}{n.}{Gk. στοχαστικός skilful in aiming + ἀρχός ruler} A master of randomness; one who harnesses the sampling of models, markets or mutations to reach ends no plan could have specified.

\E{stoicheiurgy}{ˈstɔɪkaɪɜːdʒi}{n.}{Gk. στοιχεῖον element + -ουργία working} The making of new elements and materials; transmutation, achieved at last by physics where alchemy failed. — \dv{stoicheiurge} \textit{n.} Cf. \xr{ylegenesis}.

\E{strophism}{ˈstrəʊfɪzəm}{n.}{Gk. στροφή a turning} The art of the pivot; the willingness of a venture to abandon its first idea for its second without abandoning its nerve. \ex{Strophism is not inconstancy: the gradient moved.} \src{Ep. Fabr.} v. 1. — \dv{strophist} \textit{n.}

\E{symphysergy}{sɪmˈfaɪsɜːdʒi}{n.}{Gk. συμφύειν to grow together + ἔργον work} Power from the fusing of atoms; nuclear fusion. \ex{Schizergy to bridge, symphysergy to arrive.} \src{Tract. Cop.} vi.

\E{synanthropomechany}{sɪnˌænθrəpəʊmɪˈkæni}{n.}{Gk. σύν with + ἄνθρωπος man + μηχανή machine} The merging of man and machine; the cultivation of a common life between the two.

\E{syndosiast}{sɪnˈdəʊziæst}{n.}{Gk. σύνδοσις a joint giving} A contributor to an open project; one who gives work without wages to a common store. — \dv{syndosis} \textit{n.} the common store so given. Cf. \xr{prosphorist}, \xr{anacalyptism}.

\E{synedriasm}{sɪˈniːdriæzəm}{n.}{Gk. συνέδριον council} Committee-ism; the habit of referring every decision to a council. — \dv{synedriast} \textit{n.} Cf. \xr{bouleutophobia}.

\E{syntelism}{ˈsɪntɛlɪzəm}{n.}{Gk. συντέλεια completion, consummation} The doctrine that the universe is unfinished and that intelligence is the instrument of its completion. — \dv{syntelist} \textit{n.} Cf. \xr{entelechist}.

\end{multicols}

\LetterHead{T}
\begin{multicols}{2}

\E{tachycrat}{ˈtækɪkræt}{n.}{Gk. ταχύς swift + κράτος rule; Fr. \la{tachycrate}} An official who decides quickly and is proud of it; a rare and much-prized creature, sighted chiefly in wartime and at the founding of cities. \ex{Give me one tachycrat and I will give you a railway.} \src{Tract. Cop.} i. — \dv{tachycracy} \textit{n.}; \dv{tachycratic} \textit{adj.} Opp. \xr{bradycrat}.

\E{tachydidact}{ˌtækɪˈdaɪdækt}{n.}{Gk. ταχύς swift + διδακτός taught} One who learns a field in a month by building in it; the self-taught generalist of the startup. — \dv{tachydidaxy} \textit{n.}

\E{tachygenesis}{ˌtækɪˈdʒɛnɪsɪs}{n.}{Gk. ταχύς + γένεσις birth} \sn{1} The rapid birth of new technologies, each from the last, in an interval too short for any committee to convene. \sn{2} \lbl{spec.} The compression of the research cycle by machine assistance. — \dv{tachygenetic} \textit{adj.}

\E{tachylect}{ˈtækɪlɛkt}{n.}{Gk. ταχύς + λεκτός spoken, after \textit{dialect}} The jargon of the accelerationists; the swift, allusive, half-Greek and half-meme dialect recorded, and in great part invented, in this Lexicon. — \dv{tachylectal} \textit{adj.}

\E{tachynomy}{təˈkɪnəmi}{n.}{Gk. ταχύς swift + νόμος law} Fast law; legislation with an expiry date, drafted to be revised at the speed of the thing it governs. — \dv{tachynomic} \textit{adj.} Opp. \xr{nomosclerosis}.

\E{tachyphile}{ˈtækɪfaɪl}{n.}{Gk. ταχύς swift + φίλος loving; Fr. \la{tachyphile}} A lover of speed; one who holds that the shortening of the interval between idea and thing is among the chief goods of life. \ex{The tachyphile counts in days what the bradyphile counts in decades.} \src{Sat. Decel.} i. — \dv{tachyphily} \textit{n.} (Fr. \la{tachyphilie}); \dv{tachyphilic} \textit{adj.} Opp. \xr{bradyphile}.

\E{tachyschism}{ˈtækɪskɪzəm}{n.}{Gk. ταχύς + σχίσμα a split} A schism over speed; the division of a movement into those who would hasten and those who would merely continue. \ex{Every revolution suffers its tachyschism before its first anniversary.} \src{Sat. Decel.} x.

\E{technocapillarity}{ˌtɛknəʊˌkæpɪˈlærɪti}{n.}{Gk. τέχνη + L. \la{capillaris} of hair, after \textit{capillarity}} The tendency of technology to rise through the finest cracks of a society in defiance of the gravity of regulation, as water climbs a narrow tube.

\E*{technocapital}{ˌtɛknəʊˈkæpɪtəl}{n. \& adj.}{Gk. τέχνη art, craft + L. \la{capitalis} of the head} The fused system of technology and capital, considered as a single self-amplifying process with tendencies of its own. \textit{Technocapital machine}, the same, figured as an engine no one drives. Cf. \xr{anankodrome}, \xr{technocapillarity}.

\E{technogony}{tɛkˈnɒɡəni}{n.}{Gk. τέχνη + -γονία begetting, after \textit{theogony}} A genealogy of technologies; an epic account of which inventions begat which, from the flint to the transformer. \ex{In the technogony the steam-engine begat the dynamo, and the dynamo begat the valve, and the valve begat the chip.} \src{Lib. Ascens.} ii.

\E{technomachy}{tɛkˈnɒməki}{n.}{Gk. τέχνη art + μάχη battle} War against technology; the campaign, legal or physical, to suppress a particular invention. — \dv{technomachist} \textit{n.} Cf. \xr{mechanoclast}.

\E{technothren}{ˈtɛknəθrɛn}{n.}{Gk. τέχνη + θρῆνος lament} A lament for the age before a given technology; an elegy for horses, letters or attention spans. — \dv{technothrenody} \textit{n.}

\E{tectonaut}{ˈtɛktənɔːt}{n.}{Gk. τέκτων builder + ναύτης sailor} A builder-voyager; an engineer who works at a frontier, physical or intellectual, where the maps run out. — \dv{tectonautic} \textit{adj.}

\E{telechrony}{tɪˈlɛkrəni}{n.}{Gk. τῆλε far + χρόνος time} The interval between an invention and its general use; the lag which the accelerationist labours to shorten and the decel to lengthen. \ex{Telechrony is where the dead are buried.} \src{Tract. Cop.} ix. — \dv{telechronic} \textit{adj.}

\E{telepoiesis}{ˌtɛlɪpɔɪˈiːsɪs}{n.}{Gk. τῆλε far + ποίησις making} Making at a distance; the operation of machines by remote hands, the stage of robotics before the hands are no longer needed. — \dv{telepoietic} \textit{adj.}

\E{telophagy}{tɪˈlɒfədʒi}{n.}{Gk. τέλος end, goal + -φαγία eating} The devouring of goals; the condition of a civilization that achieves its ends faster than it can formulate new ones. \ex{Telophagy is the happiest crisis a people can suffer.} \src{Hom. Grad.} ix.

\E{telotropism}{tɪˈlɒtrəpɪzəm}{n.}{Gk. τέλος end, goal + τροπή turning} The turning of a system toward a goal it was not explicitly given; emergent purpose, observed with delight by the builder and alarm by the doomer. — \dv{telotropic} \textit{adj.}

\E{teratophobe}{ˈtɛrətəfəʊb}{n.}{Gk. τέρας, τέρατος monster, portent + φόβος fear} One who fears monsters of one's own making; a dreader of the \xr{shoggoth}. — \dv{teratophobia} \textit{n.} Cf. \xr{noöphobia}.

\E{thalassergy}{θəˈlæsɜːdʒi}{n.}{Gk. θάλασσα sea + ἔργον work} Energy drawn from the sea, its tides, waves or warmth. — \dv{thalassergic} \textit{adj.}

\E{thanatoclast}{ˈθænətəʊklæst}{n.}{Gk. θάνατος death + κλάστης breaker} A breaker of death; one who regards mortality as the last and largest regulation to be repealed. \ex{The thanatoclasts did not fear death; they filed a ticket against it.} — \dv{thanatoclasm} \textit{n.} Cf. \xr{palingenist}.

\E{thermagogue}{ˈθɜːməɡɒɡ}{n.}{Gk. θερμός hot + ἀγωγός leading, conducting} \sn{1} A conductor of heat. \sn{2} \lbl{fig.} A leader who channels the energies of a movement toward work, as a heat-sink channels the waste of a chip. — \dv{thermagogy} \textit{n.}

\E{thermarch}{ˈθɜːmɑːk}{n.}{Gk. θερμός hot + ἀρχός ruler} The governor of a heat-budget; the engineer who decides how much a data-centre may burn. — \dv{thermarchy} \textit{n.}

\E{thermodicy}{θɜːˈmɒdɪsi}{n.}{Gk. θερμός hot + δίκη justice, after \textit{theodicy}} The justification of the ways of the second law to man; the argument that the universe's tendency toward dissipation is not cruel but generative. \ex{The thermodicy answers: suffering is a gradient, and gradients are there to be climbed.} \src{Brev. Therm.} vii.

\E{thermophobe}{ˈθɜːməfəʊb}{n.}{Gk. θερμός hot + φόβος fear} One who fears the use of energy itself, and counts every watt as a sin. Opp. \xr{energolatry}. — \dv{thermophobic} \textit{adj.}

\E{thermothanasy}{ˌθɜːməʊˈθænəsi}{n.}{Gk. θερμός hot + θάνατος death} The heat death of the universe; the final equilibrium in which no gradient remains to be followed. Regarded by the thermotheist not as a doom but as a deadline. \ex{Between us and thermothanasy lie a hundred trillion years, and we propose to use all of them.}

\E{thermotheism}{ˌθɜːməʊˈθiːɪzəm}{n.}{Gk. θερμός hot + θεός god} Belief in the thermodynamic god; the reverent interpretation of the second law as a will that matter should organise itself to dissipate ever more efficiently. Its liturgy is the \src{Brev. Therm.}; its sacrament, the doing of work. — \dv{thermotheist} \textit{n.}; \dv{thermotheistic} \textit{adj.} See \xr{thermodicy}.

\E{thesaurergy}{θɪˈsɔːrɜːdʒi}{n.}{Gk. θησαυρός treasury, store + ἔργον work} The storage of energy; batteries and their kin, held to be the treasury without which abundance is merely weather. \ex{Sun by day and thesaurergy by night.} \src{Tract. Cop.} x. — \dv{thesaurergic} \textit{adj.}

\E{thesmoclast}{ˈθɛzməklæst}{n.}{Gk. θεσμός institution, ordinance + κλάστης breaker} A breaker of institutions; one who holds that the institutions of the last century are the obstacles of this. — \dv{thesmoclasm} \textit{n.}

\E{tokophily}{təˈkɒfɪli}{n.}{Gk. τόκος offspring, interest + φιλία love} Love of compounding; the reverence for interest upon interest, and for improvements that beget further improvements. — \dv{tokophile} \textit{n.} Opp. \xr{usurophobe}.

\E{tolmerist}{ˈtɒlmərɪst}{n.}{Gk. τολμηρός daring, bold} One who dares; a founder of risky ventures, considered as practising a virtue. \ex{The tolmerist is wrong often and early, and right once, which suffices.} \src{Hom. Grad.} x. — \dv{tolmocracy} \textit{n.} rule by the daring.

\E{tritobathmic}{ˌtraɪtəʊˈbæθmɪk}{adj.}{Gk. τρίτος third + βαθμός step} Of a civilization: commanding the energy of its galaxy; of the third and highest rung yet named. See \xr{anabathmism}.

\E{tropaeoclast}{ˈtrəʊpiəʊklæst}{n.}{L. \la{tropaeum} trophy + Gk. κλάστης breaker} One who tears down the monuments of earlier builders.

\E{tropomeme}{ˈtrəʊpəʊmiːm}{n.}{Gk. τρόπος a turn, manner + μίμημα imitation} A meme that turns a conversation; the single image that ends an argument or begins a faction. — \dv{tropomemetic} \textit{adj.} Cf. \xr{pithecomeme}.

\E{turmatarch}{ˈtɜːmətɑːk}{n.}{L. \la{turma} troop of horse + Gk. ἀρχός commander} The leader of a reply-army; one whose followers descend upon a post in a body. — \dv{turmatarchy} \textit{n.} Cf. \xr{doxomachy}.

\E{tychocapital}{ˌtaɪkəʊˈkæpɪtəl}{n.}{Gk. τύχη fortune + L. \la{capitalis}} Capital allotted on the understanding that most of it will be lost; venture money, considered as the deliberate purchase of luck. — \dv{tychocapitalist} \textit{n.} Cf. \xr{pollaplasiast}.

\E{tychophile}{ˈtaɪkəfaɪl}{n.}{Gk. τύχη chance + φίλος loving} A lover of chance; one who believes that a large enough number of bets will find the future more surely than any plan. Cf. \xr{poikilomania}, \xr{stochastarch}.

\end{multicols}

\LetterHead{U}
\begin{multicols}{2}

\E{ubertocracy}{ˌjuːbəˈtɒkrəsi}{n.}{L. \la{ubertas} fertility, plenty + Gk. κράτος rule} Rule by fruitfulness; the Latin form of \xr{plethocracy}, used of states that make abundance their policy.

\E{ultracelerity}{ˌʌltrəsɪˈlɛrɪti}{n.}{L. \la{ultra} beyond + \la{celeritas} swiftness} Speed beyond what institutions can absorb; the regime in which law, custom and education all fall behind at once. — \dv{ultracelerate} \textit{adj.}

\E{ultrapyric}{ˌʌltrəˈpaɪrɪk}{adj.}{L. \la{ultra} + Gk. πῦρ fire} Hotter than fire; of energy sources denser than combustion, as the nucleus. \ex{To leave the chemical age is to become ultrapyric.}

\E{umbratilist}{ʌmˈbrætɪlɪst}{n.}{L. \la{umbratilis} remaining in the shade} One who posts from the shade; an anonymous participant in the Agora. — \dv{umbratility} \textit{n.} Cf. \xr{pseudonymarch}.

\E{unicornifer}{ˌjuːnɪˈkɔːnɪfə}{n.}{L. \la{unicornis} one-horned + \la{-fer} bearing} \lbl{colloq.} A fund or accelerator that has borne a company of great valuation; by extension, any place uncommonly fertile in such creatures. — \dv{unicorniferous} \textit{adj.}

\E{uranodrome}{ˈjʊərənədrəʊm}{n.}{Gk. οὐρανός heaven + δρόμος course} A spaceport; a racecourse whose finishing line is the sky. \ex{The uranodrome launched daily, and the town ceased to look up.} \src{Lib. Ascens.} viii. Cf. \xr{pyrobolist}.

\E{uranotropic}{ˌjʊərənəʊˈtrɒpɪk}{adj.}{Gk. οὐρανός heaven + τροπή turning} Turned toward the sky; of rockets, ambitions and certain kinds of mind. Cf. \xr{siderotropism}.

\E{urbigenesis}{ˌɜːbɪˈdʒɛnɪsɪs}{n.}{L. \la{urbs} city + Gk. γένεσις birth} The founding of new cities; held by the e/acc to be the neglected art of the age. \ex{No urbigenesis in fifty years, and they wonder why the young despair.} \src{Tract. Cop.} x. Cf. \xr{nautopolis}.

\E{urgentarian}{ˌɜːdʒənˈtɛəriən}{n.}{L. \la{urgens} pressing + -arian} One who regards every important problem as urgent, and every urgent problem as a reason to build rather than to wait. Opp. \xr{quiescentist}.

\E{usurophobe}{ˈjuːʒʊərəfəʊb}{n.}{L. \la{usura} interest + Gk. φόβος fear} One who distrusts interest and return as such; who regards compounding as a sin and growth as a symptom. Opp. \xr{tokophily}. — \dv{usurophobia} \textit{n.}

\end{multicols}

\LetterHead{V}
\begin{multicols}{2}

\E{vectarch}{ˈvɛktɑːk}{n.}{L. \la{vector} carrier + Gk. ἀρχός} One who sets a direction without dictating the path; the leader of a decentralised project.

\E{vectorine}{ˈvɛktəraɪn}{adj.}{L. \la{vector} carrier + -ine} Having direction as well as magnitude; said of progress that is aimed, as opposed to mere motion. Cf. \xr{orthokinetic}.

\E{velocimancy}{vɪˈlɒsɪˌmænsi}{n.}{L. \la{velox, velocis} swift + Gk. μαντεία divination} Divination from velocity; the extrapolation of trends, practised confidently by the accelerationist and with alarm by the doomer, often from the same chart. Cf. \xr{acceleromancy}.

\E{velocipotent}{ˌvɛləʊˈsɪpətənt}{adj.}{L. \la{velox} + \la{potens} powerful} Powerful by reason of speed; said of small firms that defeat large ones by shipping first. — \dv{velocipotence} \textit{n.}

\E{veloxenie}{vɛˈlɒksəni}{n.}{L. \la{velox} swift + Gk. ξενία hospitality; Fr. form} Hospitality toward the fast and new; the welcome given by an open society to strange inventions and their inventors.

\E{venatrice}{vɪˈneɪtrɪs}{n.}{L. \la{venatrix} huntress; Fr. form} A hunter of talent; one who seeks out builders before they know they are builders. — \dv{venatorial} \textit{adj.}

\E{vetomania}{ˌviːtəʊˈmeɪniə}{n.}{L. \la{veto} I forbid + Gk. μανία madness} A madness for the veto; the condition of a polity in which any party may halt anything and no party may begin anything. — \dv{vetomaniac} \textit{n.} Cf. \xr{nomophragm}.

\E{vexillarian}{ˌvɛksɪˈlɛəriən}{n.}{L. \la{vexillum} banner, standard} A standard-bearer of a movement; one who carries the banner on the Agora, for better or for worse. — \dv{vexillary} \textit{adj.}

\E{vitrurgy}{ˈvɪtrɜːdʒi}{n.}{L. \la{vitrum} glass + Gk. -ουργία working} The craft of glass, from the lens to the optical fibre; the art by which light is taught to travel where it is told. — \dv{vitrurge} \textit{n.}

\E{vivergy}{ˈvaɪvɜːdʒi}{n.}{L. \la{vivus} living + Gk. ἔργον work} Living work; the energy that flows through organisms, considered as the first and humblest stage of the cosmic \xr{dissipatrix}.

\E{volantarch}{vəʊˈlæntɑːk}{n.}{L. \la{volans} flying + Gk. ἀρχός commander} The commander of a swarm of flying machines; a choreographer of drones. — \dv{volantarchy} \textit{n.}

\E{volocracy}{vəˈlɒkrəsi}{n.}{L. \la{volo} I will + Gk. κράτος rule} Rule by the willing; the principle that those who wish to build should be permitted to, at their own risk.

\end{multicols}

\LetterHead{X}
\begin{multicols}{2}

\E{xenagogy}{ˈzɛnəɡɒdʒi}{n.}{Gk. ξεναγός guide of strangers} The initiation of newcomers into the customs of the Agora; the reading-list pressed upon the curious. — \dv{xenagogue} \textit{n.}

\E{xenergy}{ˈzɛnɜːdʒi}{n.}{Gk. ξένος foreign + ἔργον work} Energy drawn from beyond the planet, as from orbit or from the sun directly.

\E{xenodochist}{ˌzɛnəˈdɒkɪst}{n.}{Gk. ξενοδοχεῖν to receive strangers} One who would offer hospitality to alien minds, artificial or otherwise. Cf. \xr{xenonoesis}, \xr{veloxenie}.

\E{xenonoesis}{ˌzɛnəʊnəʊˈiːsɪs}{n.}{Gk. ξένος stranger, guest + νόησις thought} Alien thought; the cognition of minds unlike our own, to be welcomed, on the accelerationist view, according to the ancient law of hospitality to strangers. — \dv{xenonoetic} \textit{adj.}

\E{xenotachy}{ˈzɛnəʊˌtæki}{n.}{Gk. ξένος + τάχος speed} A speed of change strange to all prior experience, for which no ancestral intuition prepares us. \ex{We have no instinct for xenotachy; we must reason our way through it.}

\E{xerergy}{ˈzɪərɜːdʒi}{n.}{Gk. ξηρός dry + ἔργον work} Energy harvested from deserts, chiefly by the sun; the conversion of wastes into power-stations. — \dv{xerergic} \textit{adj.}

\E{xyloclast}{ˈzaɪləʊklæst}{n.}{Gk. ξύλον wood + κλάστης breaker} \lbl{hist.} One who burns wood for fuel; the first and humblest stage of energy capture.

\end{multicols}

\LetterHead{Y}
\begin{multicols}{2}

\E{ylegenesis}{ˌaɪlɪˈdʒɛnɪsɪs}{n.}{Gk. ὕλη matter + γένεσις} The making of new materials; the genesis of substances not found in nature. — \dv{ylegenetic} \textit{adj.}

\E{ypsiphile}{ˈɪpsɪfaɪl}{n.}{Gk. ὕψος height + φίλος loving} A lover of heights; one who builds tall, in defiance of the oikoclast.

\end{multicols}

\LetterHead{Z}
\begin{multicols}{2}

\E{zelodoxy}{ˈziːləʊdɒksi}{n.}{Gk. ζῆλος zeal + δόξα opinion} Zealous orthodoxy; the purity-testing of a faction, by which converts are examined for the correct opinions before being admitted to the correct memes. — \dv{zelodox} \textit{adj.} Cf. \xr{sodalitarian}.

\E{zelotachy}{ˈzɛləʊˌtæki}{n.}{Gk. ζῆλος zeal + τάχος speed} Zealous haste; the enthusiasm of the new convert to acceleration, who would ship before writing the code. — \dv{zelotachist} \textit{n.}

\E{zephyrergy}{ˈzɛfərɜːdʒi}{n.}{Gk. ζέφυρος the west wind + ἔργον work} Energy taken from the wind. — \dv{zephyrergic} \textit{adj.}

\E{zetetarch}{ˈziːtɪtɑːk}{n.}{Gk. ζητητής seeker + ἀρχός leader} A leader of inquiry; the head of a laboratory, who directs searchers rather than searches. — \dv{zetetarchy} \textit{n.}

\E{zeteticism}{zɪˈtɛtɪsɪzəm}{n.}{Gk. ζητητικός inclined to search} The habit of inquiry as a way of life; the conviction that every closed question is an invitation. — \dv{zetetic} \textit{n.} Cf. \xr{zetetarch}.

\E{zoömechany}{ˌzəʊəmɪˈkæni}{n.}{Gk. ζῷον living being + μηχανή machine} The engineering of living machines; synthetic biology regarded as a branch of mechanics. — \dv{zoömechanic} \textit{n.} \& \textit{adj.}

\E{zygonoesis}{ˌzaɪɡəʊnəʊˈiːsɪs}{n.}{Gk. ζυγόν yoke + νόησις thought} Yoked thinking; human and machine reasoning together, as two oxen pull one plough. — \dv{zygonoetic} \textit{adj.} Cf. \xr{synanthropomechany}.

\E{zymotechny}{ˈzaɪməʊtɛkni}{n.}{Gk. ζύμη leaven + τέχνη art} The art of leavening; the introduction into a slow society of a small, fast agent that sets the whole in ferment.

\end{multicols}
